% ============================================================================
% BaoCaoDATN.tex - Graduation thesis - canonical complete source
% Build: latexmk -xelatex BaoCaoDATN.tex
% All LaTeX artifacts are written to build/.
% ============================================================================
\documentclass[12pt, a4paper, oneside]{report}

% ============================================================================
% preamble.tex - formatting per Appendix 2 (university regulations)
% Times New Roman 13pt, 1.5 line spacing, margins top 3 / bottom 3.5 / left 3.5 / right 2 cm.
% Build with XeLaTeX (system fonts + Vietnamese Unicode).
% ============================================================================

\usepackage{fontspec}          % system fonts (XeLaTeX/LuaLaTeX)
\usepackage{polyglossia}       % multilingual, Vietnamese support
\setmainlanguage{vietnamese}
\setotherlanguage{english}

% --- Font: Times New Roman (falls back to Liberation Serif when TNR is absent) ---
% Detection is automatic below; hard-set \setmainfont{Times New Roman} to force it.
\IfFontExistsTF{Times New Roman}{
  \setmainfont{Times New Roman}
}{
  \setmainfont{Liberation Serif}  % metric-compatible substitute for Times New Roman
}

% report.cls ignores the 13pt option, so set the body size explicitly.
\usepackage{anyfontsize}
\usepackage{setspace}
\AtBeginDocument{\fontsize{13}{19.5}\selectfont}
\setlength{\emergencystretch}{3em}

% --- Margins: top 3cm, bottom 3.5cm, left 3.5cm, right 2cm ---
\usepackage[
  top=3cm, bottom=3.5cm, left=3.5cm, right=2cm,
  a4paper
]{geometry}

% --- 1.5 line spacing ---
\onehalfspacing

% --- Number headings to 4 levels (Chapter / x.y / x.y.z / x.y.z.t) ---
\usepackage{titlesec}
\setcounter{secnumdepth}{4}
\setcounter{tocdepth}{3}

% Chapter: bold ~14pt; subsections: bold 13pt (per regulations).
\titleformat{\chapter}[block]
  {\normalfont\fontsize{14}{18}\bfseries}
  {Chương \thechapter.}{0.6em}{}
\titlespacing*{\chapter}{0pt}{0pt}{12pt}

\titleformat{\section}
  {\normalfont\fontsize{13}{18}\bfseries}{\thesection.}{0.5em}{}
\titleformat{\subsection}
  {\normalfont\fontsize{13}{18}\bfseries}{\thesubsection.}{0.5em}{}
\titleformat{\subsubsection}
  {\normalfont\fontsize{13}{18}\bfseries}{\thesubsubsection.}{0.5em}{}

% --- Tables, figures, captions numbered per chapter (Figure 1.1, Table 2.1) ---
\usepackage{graphicx}
\usepackage{float}
\usepackage{booktabs}
\usepackage{longtable}
\usepackage{array}
\usepackage[labelsep=colon]{caption}
\captionsetup{font=small, labelfont=bf}
\usepackage{chngcntr}
\counterwithin{figure}{chapter}
\counterwithin{table}{chapter}

% --- Page numbers: bottom-center (regulation) ---
\usepackage{fancyhdr}
\pagestyle{fancy}
\fancyhf{}
\fancyfoot[C]{\thepage}
\renewcommand{\headrulewidth}{0pt}

% --- IEEE citations (numeric [1], [2]) ---
\usepackage[backend=biber, style=ieee, sorting=nyt]{biblatex}

% --- Utilities ---
\usepackage{hyperref}
\hypersetup{colorlinks=false, hidelinks, pdfborder={0 0 0}}
\usepackage{xurl}
\usepackage{enumitem}
\usepackage{amsmath, amssymb}
\usepackage{tikz}
\usetikzlibrary{shapes.geometric, arrows.meta, positioning, fit, backgrounds, calc}

% Framed callout used for the methodological caveat blocks in the results chapter.
\usepackage[most]{tcolorbox}

% Page frames copied from the two cover sections in docs/BieuMau.docx.
% Main cover: thin outer rule and thick inner rule, 24 pt from the page edge.
\newcommand{\maincoverframe}{%
  \begin{tikzpicture}[remember picture, overlay]
    \draw[line width=0.5pt]
      ($(current page.north west)+(0.85cm,-0.85cm)$)
      rectangle
      ($(current page.south east)+(-0.85cm,0.85cm)$);
    \draw[line width=1.5pt]
      ($(current page.north west)+(0.98cm,-0.98cm)$)
      rectangle
      ($(current page.south east)+(-0.98cm,0.98cm)$);
  \end{tikzpicture}%
}

% Secondary cover: the template's narrow double rule, also 24 pt from the edge.
\newcommand{\secondarycoverframe}{%
  \begin{tikzpicture}[remember picture, overlay]
    \draw[line width=0.5pt]
      ($(current page.north west)+(0.85cm,-0.85cm)$)
      rectangle
      ($(current page.south east)+(-0.85cm,0.85cm)$);
    \draw[line width=0.5pt]
      ($(current page.north west)+(0.90cm,-0.90cm)$)
      rectangle
      ($(current page.south east)+(-0.90cm,0.90cm)$);
  \end{tikzpicture}%
}

% --- Custom commands for the title pages ---
\newcommand{\thesistitle}{Xây dựng hệ thống dự báo xu hướng biến động giá cổ phiếu dựa trên phân tích cảm xúc tin tức tài chính bằng mô hình Transformer và chuỗi thời gian}
\newcommand{\thesistitleEN}{Building a stock price movement forecasting system based on financial news sentiment analysis using Transformer and time-series models}
\newcommand{\authorname}{Nguyễn Phước Thọ}
\newcommand{\studentid}{25410139}
\newcommand{\advisor}{TS. Đặng Văn Thìn}
\newcommand{\major}{Trí tuệ nhân tạo}
\newcommand{\thesisyear}{2026}
% Keep the PDF metadata limited to the academic report identity.
\AtBeginDocument{%
  \hypersetup{
    pdftitle={\thesistitle},
    pdfauthor={\authorname},
    pdfsubject={Đồ án tốt nghiệp ngành Trí tuệ nhân tạo}
  }%
}

\addbibresource{references.bib}

% Architecture diagrams are included in the canonical report.
\newif\ifShowDiagram
\ShowDiagramtrue

\begin{document}

% ------------------------------------------------------------------ FRONT
\pagenumbering{gobble}

% The official cover template uses wider cover margins than the report body.
\newgeometry{top=2cm,bottom=1.5cm,left=2.54cm,right=2cm}

% Main title page
\begin{titlepage}
\maincoverframe
\centering
{\fontsize{15}{18}\bfseries ĐẠI HỌC QUỐC GIA TP. HỒ CHÍ MINH\par}
\vspace{0.15cm}
{\fontsize{16}{20}\bfseries TRƯỜNG ĐẠI HỌC CÔNG NGHỆ THÔNG TIN\par}
\vspace{0.15cm}
{\fontsize{16}{20}\bfseries KHOA KHOA HỌC MÁY TÍNH\par}

\vspace{2.5cm}
{\fontsize{14}{18}\bfseries \MakeUppercase{\authorname}\par}
\vspace{2.4cm}

{\fontsize{16}{20}\bfseries ĐỒ ÁN TỐT NGHIỆP\par}
\vspace{0.35cm}
{\fontsize{18}{22}\bfseries \thesistitle\par}
\vspace{0.3cm}
{\fontsize{16}{20}\bfseries \thesistitleEN\par}

\vspace{1.5cm}
{\fontsize{14}{18}\bfseries CỬ NHÂN NGÀNH \MakeUppercase{\major}\par}

\vfill
{\fontsize{13}{16}\bfseries TP. HỒ CHÍ MINH, \thesisyear\par}
\vspace*{1cm}
\end{titlepage}

% Secondary title page (adds the supervisor)
\begin{titlepage}
\secondarycoverframe
\centering
{\fontsize{14}{18}\bfseries ĐẠI HỌC QUỐC GIA TP. HỒ CHÍ MINH\par}
\vspace{0.15cm}
{\fontsize{16}{20}\bfseries TRƯỜNG ĐẠI HỌC CÔNG NGHỆ THÔNG TIN\par}
\vspace{0.15cm}
{\fontsize{16}{20}\bfseries KHOA KHOA HỌC MÁY TÍNH\par}

\vspace{2.5cm}
{\fontsize{14}{18}\bfseries \MakeUppercase{\authorname}{} -- \studentid\par}
\vspace{2.4cm}

{\fontsize{16}{20}\bfseries ĐỒ ÁN TỐT NGHIỆP\par}
\vspace{0.35cm}
{\fontsize{18}{22}\bfseries \thesistitle\par}
\vspace{0.3cm}
{\fontsize{16}{20}\bfseries \thesistitleEN\par}

\vspace{1.5cm}
{\fontsize{14}{18}\bfseries CỬ NHÂN NGÀNH \MakeUppercase{\major}\par}

\vspace{1.2cm}
{\fontsize{14}{18}\bfseries GIẢNG VIÊN HƯỚNG DẪN\par}
{\fontsize{14}{18}\bfseries \advisor\par}

\vfill
{\fontsize{13}{16}\bfseries TP. HỒ CHÍ MINH, \thesisyear\par}
\vspace*{1cm}
\end{titlepage}

\restoregeometry

% Examination committee info (unnumbered page)
\chapter*{THÔNG TIN HỘI ĐỒNG CHẤM ĐỒ ÁN TỐT NGHIỆP}
\addcontentsline{toc}{chapter}{Thông tin hội đồng chấm đồ án tốt nghiệp}
\thispagestyle{empty}

% TODO: điền số quyết định thành lập hội đồng (do Khoa/Nhà trường cấp) trước khi in nộp.
Hội đồng chấm Đồ án tốt nghiệp, thành lập theo Quyết định số \underline{\hspace{3cm}}
% TODO: điền ngày ký quyết định thành lập hội đồng trước khi in nộp.
ngày \underline{\hspace{2.5cm}} của Hiệu trưởng Trường Đại học Công nghệ Thông tin.

\vspace{1cm}
\begin{enumerate}
% TODO: điền học hàm/học vị và họ tên Chủ tịch hội đồng trước khi in nộp.
  \item \underline{\hspace{7cm}}, Chủ tịch.
% TODO: điền học hàm/học vị và họ tên Thư ký hội đồng trước khi in nộp.
  \item \underline{\hspace{7cm}}, Thư ký.
% TODO: điền học hàm/học vị và họ tên Ủy viên hội đồng trước khi in nộp.
  \item \underline{\hspace{7cm}}, Ủy viên.
\end{enumerate}

% Acknowledgements (unnumbered page)
\chapter*{LỜI CẢM ƠN}
\addcontentsline{toc}{chapter}{Lời cảm ơn}
\thispagestyle{empty}

Trong quá trình học tập tại Trường Đại học Công nghệ Thông tin và thực hiện đồ án
tốt nghiệp, em đã nhận được sự hướng dẫn và hỗ trợ từ Nhà trường cùng quý thầy cô.

Em xin chân thành cảm ơn Trường Đại học Công nghệ Thông tin, Khoa Khoa học Máy
tính và quý thầy cô đã tạo điều kiện học tập, truyền đạt kiến thức và giúp em xây dựng
nền tảng chuyên môn để thực hiện đề tài này.

Em xin bày tỏ lòng biết ơn đến Tiến sĩ Đặng Văn Thìn, giảng viên trực tiếp hướng dẫn
đồ án. Trong suốt quá trình thực hiện, thầy đã dành thời gian trao đổi, định hướng và
góp ý cụ thể cho từng phần của đề tài. Những nhận xét của thầy giúp em nhận ra các
thiếu sót, điều chỉnh phương pháp thực hiện và hoàn thiện báo cáo.

Do kiến thức và kinh nghiệm nghiên cứu còn hạn chế, báo cáo khó tránh khỏi những
thiếu sót. Em mong nhận được góp ý từ Nhà trường và giảng viên hướng dẫn để tiếp tục
hoàn thiện kiến thức và rút kinh nghiệm cho các công việc sau này.

\vspace{1cm}
\begin{flushright}
TP. Hồ Chí Minh, \thesisyear\\
Sinh viên thực hiện\\[1cm]
\authorname
\end{flushright}

\tableofcontents
\clearpage
\listoffigures
\clearpage
\listoftables
\clearpage
% List of abbreviations (alphabetical order)
\chapter*{DANH MỤC TỪ VIẾT TẮT}
\addcontentsline{toc}{chapter}{Danh mục từ viết tắt}
\thispagestyle{empty}

\begin{longtable}{p{3cm} p{10.5cm}}
\toprule
Từ viết tắt & Diễn giải \\
\midrule
\endhead
AUC     & Diện tích dưới đường cong ROC \\
BiLSTM  & Mạng ghi nhớ dài-ngắn hai chiều \\
CI      & Khoảng tin cậy \\
HOSE    & Sở Giao dịch Chứng khoán Thành phố Hồ Chí Minh \\
LSTM    & Mạng ghi nhớ dài-ngắn \\
NLP     & Xử lý ngôn ngữ tự nhiên \\
OHLCV   & Giá mở cửa - cao nhất - thấp nhất - đóng cửa - khối lượng \\
PhoBERT & Mô hình ngôn ngữ tiền huấn luyện cho tiếng Việt \\
TFT     & Bộ biến đổi hợp nhất theo thời gian \\
VN30    & Chỉ số 30 cổ phiếu vốn hóa lớn trên HOSE \\
\bottomrule
\end{longtable}

% ------------------------------------------------------------------ BODY
% Page numbering starts at the Vietnamese summary: Arabic, bottom-center.
\clearpage
\pagenumbering{arabic}
\setcounter{page}{1}

% Vietnamese thesis summary (page numbering starts here)
\chapter*{TÓM TẮT ĐỒ ÁN}
\addcontentsline{toc}{chapter}{Tóm tắt đồ án}

Đồ án nghiên cứu bài toán dự báo xu hướng biến động giá cổ phiếu trên thị trường chứng
khoán Việt Nam thông qua việc kết hợp tin tức tài chính tiếng Việt với dữ liệu giá theo
chuỗi thời gian. Bài toán được đặt ra từ thực tế rằng giá cổ phiếu vừa phản ánh lịch sử
giao dịch, vừa phản ứng với thông tin về kết quả kinh doanh, kế hoạch đầu tư, thay đổi
quản trị, chính sách, rủi ro pháp lý và các sự kiện bất thường. Phạm vi nghiên cứu tập
trung vào sàn HOSE và một nhóm mã thuộc VN30. Mục tiêu của nghiên cứu là xác định hướng
biến động trong phiên giao dịch kế tiếp theo ba lớp tăng, đi ngang và giảm, thay vì dự
báo trực tiếp giá trị tuyệt đối.

Hệ thống gồm hai nhánh. Nhánh thứ nhất sử dụng PhoBERT, một mô hình Transformer được
tiền huấn luyện cho tiếng Việt, để phân loại tin tức tài chính thành ba lớp tích cực,
trung tính và tiêu cực \cite{Nguyen2020_200300744v3}. Nhãn biểu thị tác động kỳ vọng của
nội dung đối với doanh nghiệp hoặc cổ phiếu liên quan, thay vì chỉ phản ánh sắc thái
chung của câu văn. Sau bước suy luận, hệ thống lưu xác suất của cả ba lớp thay vì chỉ
lưu nhãn có xác suất cao nhất. Các xác suất này được tổng hợp theo mã cổ phiếu và ngày
giao dịch, cùng với số lượng tin, tỷ lệ lớp, độ phân tán và cờ biểu thị sự hiện diện của
tin.

Nhánh thứ hai sử dụng dữ liệu OHLCV đã điều chỉnh và các đặc trưng nhân quả, bao gồm lợi
suất, lợi suất log, tỷ lệ giữa giá đóng cửa và đường trung bình động, thay đổi khối lượng
và độ biến động. Hai nhánh được căn chỉnh theo mã và ngày sau khi áp dụng quy tắc về thời
điểm. Thời điểm công bố tin được chuyển sang múi giờ Việt Nam; tin xuất hiện sau mốc cắt
thông tin, vào cuối tuần hoặc ngày nghỉ được ánh xạ sang phiên giao dịch kế tiếp. Quy tắc
này bảo đảm thời điểm của từng bài tin phù hợp với phiên sử dụng tin đó. Giới hạn thời
điểm của điểm kiểm cảm xúc hồi cứu được trình bày rõ trong phần kết quả.

Về kiến trúc, các mô hình được phân thành từng tầng nhằm tách biệt vấn đề dữ liệu với
mức độ phức tạp của mô hình. Tầng đầu gồm mô hình luôn dự báo lớp phổ biến nhất và mô
hình dự báo ngẫu nhiên theo phân bố lớp. Tầng cơ sở gồm LSTM chỉ sử dụng dữ liệu giá và
LSTM hai nhánh kết hợp dữ liệu giá với cảm xúc. Bộ biến đổi hợp nhất theo thời gian chỉ
được xem xét như một phương án mở rộng khi dữ liệu, tài nguyên và kết quả của cấu hình
cơ sở cho phép.

Về phương pháp, đồ án xây dựng một quy trình xử lý gồm khảo sát tài liệu, thu thập và làm
sạch dữ liệu, loại bỏ bản ghi trùng, lưu thông tin truy nguyên, gắn tin với mã cổ phiếu,
đồng bộ theo phiên, tinh chỉnh mô hình cảm xúc, tạo đặc trưng, huấn luyện và đánh giá.
Bộ dữ liệu tin tức gồm tiêu đề, nội dung khi có thể thu thập, nguồn, URL, thời điểm công
bố và mã liên quan. Các bài viết đề cập đến nhiều mã được lưu theo quan hệ nhiều mã; các
bản ghi thiếu thời điểm hoặc không xác định được mã được đánh dấu và không được đưa trực
tiếp vào bảng đặc trưng chính.

Dữ liệu giá được chia theo thứ tự thời gian bằng kiểm định cuốn chiếu với năm đoạn kiểm
thử tách biệt. Nghiên cứu không sử dụng tập kiểm thử cuối riêng biệt vì cấu hình được cố
định trước lần chạy và mỗi đoạn kiểm thử chỉ được đánh giá một lần. Bộ chuẩn hóa, bộ điền
khuyết, ngưỡng tạo nhãn, lựa chọn đặc trưng và các tham số của nhánh giá chỉ được học từ
phần huấn luyện trong từng cửa sổ. Mô hình sử dụng cảm xúc được so sánh với mô hình chỉ
sử dụng giá trên cùng các ngày huấn luyện, xác thực và kiểm thử, đồng thời áp dụng cùng
phương pháp tiền xử lý, hạt giống ngẫu nhiên hoặc tập hạt giống và tiêu chí dừng. Đối
chứng giữ nguyên kiến trúc và độ phủ tin nhằm tách chênh lệch hồi cứu do xác suất cảm
xúc khỏi ảnh hưởng của kiến trúc hoặc phương pháp chia mẫu.

Về dữ liệu, kho tin lưu trữ 10.695 bản ghi có URL, thời điểm đăng và nội dung trong cửa
sổ 2020--2025, cùng 12.624 liên kết tin--mã trên 10 mã VN30 và 14.990 dòng giá
OHLCV trong cùng giai đoạn. Tập nhãn cảm xúc trong miền dữ liệu đã được rà soát
gồm 1.306 mẫu, bên cạnh tập hạt giống CafeF gồm 999 mẫu. Kết quả xác thực chéo phân tầng
5 lượt trên tập 1.306 mẫu đạt F1 vĩ mô 0,7298, độ chính xác 0,8125 và độ chính xác cân
bằng 0,7689; tổng số mẫu NEGATIVE trong năm tập kiểm thử ngoài là 59, với độ bao phủ
0,6803.

Sau khi khóa cấu hình, một điểm kiểm triển khai được tinh chỉnh lại trên toàn bộ 1.306
nhãn với năm epoch cố định và được dùng để chấm lại 14.256 liên kết tin--mã, trong đó
12.624 liên kết thuộc cửa sổ 2020--2025 và 12.613 liên kết được neo vào một phiên bên
trong cửa sổ. Bảng dự báo gồm 14.990 dòng mã--phiên, trong đó 6.948 dòng có tin. Đối với
bài toán dự báo, mô hình LSTM chỉ sử dụng giá đạt F1 vĩ mô 0,3756, độ chính xác cân bằng
0,3915 và AUC-ROC một-chống-còn-lại 0,5714. Giá trị F1 vĩ mô này cao hơn mốc ngẫu nhiên
0,3289 và mốc lớp phổ biến 0,1453. Tầng TFT tối giản đạt 0,3619, thấp hơn kết quả của
LSTM. Do nhãn ba lớp được xác định theo phân vị một phần ba, các lớp gần cân bằng theo
thiết kế và mức ngẫu nhiên của bài toán là 0,333.

Đóng góp riêng của tín hiệu cảm xúc được đo bằng một đối chứng giữ nguyên kiến
trúc hai nhánh, bài tin, thời điểm và số lượng tin, nhưng thay xác suất cảm
xúc thực bằng phân phối tiên nghiệm trung tính. Trong lượt hồi cứu, hiệu số trực tiếp
là $+0{,}0113$ F1 vĩ mô, trong đó hiệu ứng của kiến trúc, sự hiện diện và khối lượng
tin là $+0{,}0034$, còn đóng góp thông tin theo cửa sổ là $+0{,}0080$; khoảng tin
cậy 95\% theo ngày là $[-0{,}0025; +0{,}0186]$, chứa 0.

Do điểm kiểm của lượt hồi cứu được học từ cả những nhãn công bố sau các ngày kiểm
thử, các số liệu trên chỉ phản ánh chênh lệch hồi cứu. Lượt ngoài mẫu sử dụng điểm kiểm
được đóng băng chỉ từ 1.049 nhãn công bố trước ngày quan sát kiểm thử đầu tiên, đồng thời
sử dụng một điểm kiểm đối chứng có kích thước và phân bố lớp tương ứng để tách ảnh hưởng
của kích thước tập huấn luyện. Lượt này đã được đo và xác minh. Đóng góp thông tin ngoài
mẫu của LSTM hai nhánh là $+0{,}0094$ F1 vĩ mô theo cửa sổ, với khoảng tin cậy 95\%
theo ngày $[-0{,}0009; +0{,}0196]$ chứa 0; lượt đối chứng cùng kích thước (hạt
giống 7) đạt $+0{,}0028$ theo cửa sổ, với khoảng tin cậy theo ngày $[-0{,}0074;
+0{,}0133]$ chứa 0. Cả ba ước lượng điểm đều dương và nằm trong khoảng
$+0{,}003$ đến $+0{,}009$, nhưng tất cả các khoảng tin cậy đều chứa 0. Ngoài ra, các lượt
đo này là nhiều phép so sánh trên cùng một bộ dữ liệu, nên không có cơ sở tuyên bố ý nghĩa thống kê. Bốn phân tích bổ sung được thực hiện để khảo sát
độ vững của kết luận, gồm tương quan hạng giữa điểm cảm xúc và lợi suất không qua
mô hình, nhãn lợi suất vượt trội so với các mã cùng phiên, ablation chân trời
dự báo và chẩn đoán kinh tế trong khung thăm dò không giao dịch được. Cả bốn phân tích
đều không xác định được cải thiện có thể tái lập. Vì vậy, trong phạm vi đã đo, chưa có
bằng chứng đủ mạnh cho thấy cảm xúc cải thiện kết quả dự báo; đồng thời, dữ liệu cũng
không loại trừ khả năng tồn tại một cải thiện nhỏ. Kết quả không được diễn giải
thành quan hệ nhân quả hay khuyến nghị đầu tư.

\vspace{0.5cm}
\noindent Từ khóa: dự báo xu hướng giá cổ phiếu, phân tích cảm xúc, PhoBERT,
Transformer, chuỗi thời gian, HOSE, VN30.


\chapter{MỞ ĐẦU}
\label{ch:mo-dau}
\section{Tính cấp thiết của đề tài}
\label{sec:ly-do}

\subsection{Tính cấp thiết về mặt thực tiễn}

Thị trường chứng khoán là một hệ thống có mức độ biến động cao. Giá cổ phiếu được hình
thành từ sự tương tác giữa tình hình hoạt động của doanh nghiệp, điều kiện kinh tế vĩ
mô, diễn biến ngành, dòng tiền trên thị trường và kỳ vọng của nhà đầu tư. Bên cạnh dữ
liệu giao dịch, tin tức tài chính là nguồn thông tin có thể tác động trực tiếp đến cách
nhà đầu tư đánh giá doanh nghiệp và đưa ra quyết định mua, bán hoặc nắm giữ. Thông tin
về kết quả kinh doanh, kế hoạch đầu tư, thay đổi nhân sự cấp cao, chính sách quản lý,
hoạt động mua bán và sáp nhập, rủi ro pháp lý hoặc biến động bất thường của thị trường
đều có thể làm thay đổi tâm lý nhà đầu tư trong thời gian ngắn. Do đó, nghiên cứu và
xây dựng phương pháp phân tích tác động của tin tức tài chính đến biến động giá cổ
phiếu có ý nghĩa đối với cả xử lý ngôn ngữ tự nhiên và phân tích chuỗi thời gian
\cite{Chen2019_191005078v1,Chen2021_210708721v1}.


Trong bối cảnh này, đề tài được lựa chọn dựa trên các lý do sau.

Thứ nhất, lượng thông tin tài chính được công bố hằng ngày ngày càng lớn và phân tán
trên nhiều nguồn. Trong phạm vi đồ án, các nguồn tin có thể khác nhau về cách trình
bày, thời điểm đăng tải và mức độ chi tiết. Vì vậy, cần xây dựng một quy trình tự động
để thu thập, làm sạch, gắn tin với mã cổ phiếu, kiểm tra thời điểm và tổng hợp nội dung,
qua đó bảo đảm tính nhất quán của dữ liệu và hỗ trợ phân tích ở quy mô lớn.


Thứ hai, dữ liệu giá lịch sử chỉ phản ánh kết quả giao dịch đã xảy ra và không thể hiện
đầy đủ nội dung thông tin dẫn đến thay đổi trong kỳ vọng của nhà đầu tư. Hai phiên giao
dịch có thể có đặc điểm kỹ thuật tương tự nhưng cho kết quả khác nhau nếu bối cảnh tin
tức khác nhau. Tin tức về lợi nhuận, dự án mới, thay đổi quản trị hoặc rủi ro pháp lý
có thể tạo ra các hướng phản ứng khác nhau, ngay cả khi các chỉ báo giá tại thời điểm
trước đó gần giống nhau. Việc bổ sung tín hiệu cảm xúc từ tin tức tài chính có thể cung
cấp thêm một lớp thông tin về mức độ tích cực, trung tính hoặc tiêu cực của các sự kiện
liên quan đến doanh nghiệp và thị trường. Tuy nhiên, giá trị của tín hiệu này được xem
là một giả thuyết cần kiểm chứng bằng thực nghiệm, không phải kết luận được xác lập
trước \cite{Jiang2023_230602136v3,Siala2026_260200086v3}.

Đối với Việt Nam, tin tức tài chính cần được xử lý bằng tiếng Việt. Quá trình xử lý
cần xem xét các đặc điểm của ngôn ngữ này, chẳng hạn như từ ghép, quy tắc tách từ, dấu
thanh, tên doanh nghiệp, mã cổ phiếu và cách biểu đạt số liệu. Cùng một doanh nghiệp
có thể được gọi bằng tên đầy đủ, tên viết tắt hoặc mã giao dịch; một bài viết cũng có
thể đồng thời đề cập đến nhiều doanh nghiệp. Nếu các đặc điểm này không được xử lý phù
hợp, bước gắn tin với mã cổ phiếu có thể tạo ra nhiễu trước khi mô hình được huấn luyện
\cite{Nguyen2020_200300744v3}.

PhoBERT là mô hình ngôn ngữ tiền huấn luyện dành cho tiếng Việt, do đó phù hợp để làm
mô hình nền cho bước phân loại văn bản trong đề tài \cite{Nguyen2020_200300744v3}.
Trong miền tiếng Anh, FinBERT cho thấy việc thích nghi BERT với ngữ liệu tài chính có
thể tạo ra biểu diễn hữu ích cho phân tích cảm xúc \cite{Yang2020_200608097v2}. Hai
hướng tiếp cận này gợi ý cách thiết kế đồ án: lựa chọn mô hình nền phù hợp với ngôn ngữ
của dữ liệu, tinh chỉnh mô hình và sử dụng đầu ra cho nhiệm vụ dự báo. Việc lựa chọn
PhoBERT không đồng nghĩa với việc mô hình tự động hiểu chính xác mọi thuật ngữ tài
chính. Các trường hợp liên quan đến tên riêng, số liệu, câu phủ định và bài viết chứa
nhiều sự kiện vẫn cần được kiểm tra bằng dữ liệu gán nhãn và phân tích lỗi.

HOSE và nhóm cổ phiếu VN30 được lựa chọn làm phạm vi khảo sát của đồ án vì cho phép
tổ chức dữ liệu theo một danh sách mã xác định và một lịch giao dịch thống nhất. Quyết
định này nhằm kiểm soát phạm vi thu thập, đồng bộ và đánh giá dữ liệu, không khẳng định
rằng mẫu nghiên cứu đại diện cho toàn bộ thị trường. Danh sách mã và ngày bắt đầu sử
dụng phải được ghi trong bản kê.

Đặc điểm ngành, thanh khoản, mức độ phản ứng với tin tức và mật độ bài viết của từng
mã có thể khác nhau. Thành phần rổ VN30 cũng có thể thay đổi theo thời gian, trong khi
một số doanh nghiệp có số lượng tin nhiều hơn đáng kể so với các doanh nghiệp khác.
Vì vậy, kết quả chỉ được diễn giải như bằng chứng thực nghiệm trong phạm vi mẫu, giai
đoạn và nguồn dữ liệu đã chọn.

Dữ liệu giá phản ánh kết quả giao dịch thông qua giá mở cửa, giá cao nhất, giá thấp
nhất, giá đóng cửa, khối lượng và lợi suất. Tin tức mô tả sự kiện hoặc kỳ vọng và
thường xuất hiện trước khi phản ứng tương ứng được phản ánh đầy đủ vào giá. Vì vậy, hai
nguồn dữ liệu có khả năng bổ sung cho nhau, nhưng khác nhau về dạng biểu diễn, tần suất,
độ dài và thời điểm xuất hiện. Một ngày có thể có nhiều bài viết, trong khi một ngày
khác không có bài nào, mặc dù giá vẫn được ghi nhận ở cả hai ngày. Tin tức có thể xuất
hiện trước giờ mở cửa, trong phiên, sau giờ đóng cửa hoặc vào ngày nghỉ. Trong đồ án,
quy tắc ánh xạ theo thời điểm công bố được sử dụng để hạn chế việc đưa thông tin chưa
xuất hiện vào thời điểm dự báo.

Mô hình chỉ sử dụng giá có thể bỏ qua thông tin về sự kiện và kỳ vọng, trong khi mô
hình chỉ sử dụng cảm xúc có thể không nắm bắt được xu hướng, biến động, thanh khoản và
trạng thái kỹ thuật của cổ phiếu. Đề tài xử lý riêng hai nguồn dữ liệu, tổng hợp tin
theo mã và ngày, sau đó so sánh mô hình có cảm xúc với mô hình chỉ sử dụng giá trong
cùng điều kiện. Câu hỏi cần trả lời không phải là mô hình nhiều đầu vào có điểm cao hơn, mà là nhóm đặc
trưng cảm xúc có tạo ra thông tin gia tăng sau khi lịch sử giá đã được sử dụng hay không.

Các kiến trúc kết hợp tin tức và giá trong nghiên cứu quốc tế có thể trích xuất cảm
xúc bằng FinBERT rồi đưa kết quả vào LSTM, hoặc giữ lại biểu diễn ngữ cảnh và sử dụng
mô hình ngôn ngữ lớn để biểu diễn văn bản chi tiết hơn
\cite{Halder2022_221107392v1,Gu2024_240716150v1,Elahi2024_241101368v1}. Ưu điểm của
việc sử dụng vector xác suất cảm xúc là duy trì mức độ không chắc chắn của bộ phân loại
và cho phép tổng hợp nhiều bài viết trong cùng ngày. Tuy nhiên, quá trình tổng hợp có
thể làm mất thứ tự bài viết, khác biệt giữa các nguồn và chi tiết của từng sự kiện.
Nếu giữ biểu diễn toàn văn, kích thước đặc trưng và chi phí xử lý sẽ tăng, trong khi
dữ liệu tài chính tiếng Việt có thể chưa đủ lớn để kiểm soát hiện tượng quá khớp.

Mức cải thiện quan sát được trên dữ liệu nước ngoài không thể tự động chuyển thành hiệu
quả trên HOSE. Ngôn ngữ, cấu trúc thị trường, nguồn tin, thời gian giao dịch, quy tắc gắn nhãn
giữa mô hình chỉ dùng giá và mô hình có thêm cảm xúc ở vị trí trung tâm, với cùng cửa sổ
thời gian, cách tiền xử lý, hạt giống hoặc tập hạt giống, ngân sách huấn luyện và tập kiểm thử.

Khó khăn của bài toán không chỉ liên quan đến việc lựa chọn mô hình có nhiều tham số
hơn. Có ít nhất bốn vấn đề phương pháp cần được xử lý đồng thời.

Thứ nhất, văn bản phải được gắn với đúng mã cổ phiếu và đúng thời điểm. Một bài viết
được đăng hoặc cập nhật sau thời điểm ra quyết định không thể được sử dụng để giải thích
phiên trước đó. Nếu tin sau giờ giao dịch, tin cuối tuần hoặc tin thiếu thông tin thời
gian không được xử lý thống nhất, mô hình có thể tiếp cận thông tin tương lai và tạo ra
kết quả cao nhưng không phản ánh điều kiện dự báo thực tế. Vì vậy, mỗi bản ghi cần lưu
giữ mốc thời gian, nguồn, URL và trạng thái ánh xạ để phục vụ việc kiểm tra lại.

Thứ hai, tín hiệu ở cấp bài viết phải được tổng hợp thành đặc trưng ở cấp mã và ngày.
Nếu chỉ lựa chọn một bài nổi bật hoặc chỉ tính giá trị trung bình mà bỏ qua số lượng,
tỷ lệ lớp, độ phân tán và ngày không có tin, đặc trưng có thể không phản ánh cường độ
cũng như độ phủ của thông tin. Ngày có nhiều tin trung tính phải được phân biệt với
ngày hoàn toàn không có tin vì hai tình huống này có mức độ quan sát khác nhau.

Thứ ba, dữ liệu tài chính thường mất cân bằng và phân phối có thể thay đổi theo thời
gian. Một mô hình có thể đạt độ chính xác tương đối cao chỉ bằng cách dự đoán lớp phổ
biến nhất, nhưng vẫn bỏ sót các phiên giảm hoặc đi ngang. Bài toán ba lớp cũng nhạy
cảm với cách lựa chọn ngưỡng cho vùng đi ngang. Vì vậy, đồ án sử dụng F1 vĩ mô và độ
chính xác cân bằng, kèm theo độ chính xác dương tính, độ bao phủ và F1 theo từng lớp,
ma trận nhầm lẫn và AUC-ROC một-chống-còn-lại khi đầu ra có xác suất.

Thứ tư, phép đánh giá phải tuân thủ thứ tự thời gian. Việc chuẩn hóa, điền khuyết, chọn
đặc trưng, chọn ngưỡng nhãn hoặc tinh chỉnh mô hình bằng dữ liệu tương lai đều làm giảm
giá trị của kết quả ngoài mẫu. Mô hình cũng có thể ghi nhớ bản sao của cùng một bài nếu
việc loại trùng chỉ được thực hiện sau khi chia tập. Các nghiên cứu về dự báo biến động
giá cổ phiếu cho thấy benchmark, dữ liệu và giao thức đánh giá có vai trò không kém
việc thay đổi kiến trúc \cite{Feng2018_181009936v2,Chen2021_210708721v1}.


\subsection{Tính cấp thiết về mặt lý luận}

Về mặt lý luận, ba lý do còn lại xuất phát từ sự phát triển của các mô hình ngôn ngữ
và mô hình chuỗi thời gian.

Thứ ba, sự phát triển của các mô hình ngôn ngữ tiền huấn luyện tạo điều kiện thuận lợi
cho việc xử lý văn bản theo ngữ cảnh. BERT đặt nền tảng cho việc học biểu diễn hai
chiều và tinh chỉnh một mô hình chung cho nhiều nhiệm vụ xử lý ngôn ngữ tự nhiên. Trên
cơ sở đó, PhoBERT được xây dựng như một mô hình ngôn ngữ đơn ngữ tiền huấn luyện dành
cho tiếng Việt \cite{Nguyen2020_200300744v3}. Đối với văn bản tài chính, biểu diễn theo
ngữ cảnh đặc biệt quan trọng vì sắc thái của thông tin không chỉ phụ thuộc vào từng từ
riêng lẻ mà còn phụ thuộc vào chủ thể được đề cập, hiện tượng phủ định, số liệu, thời
hạn và quan hệ giữa các sự kiện. Chẳng hạn, từ ``tăng'' có thể mô tả doanh thu, lợi
nhuận, chi phí hoặc rủi ro, do đó tác động không thể được xác định chỉ bằng cách đếm
từ khóa. Trong miền tiếng Anh, FinBERT cũng cho thấy việc thích nghi BERT với ngữ liệu
tài chính là một hướng tiếp cận có cơ sở \cite{Yang2020_200608097v2}.

Thứ tư, các mô hình học sâu dành cho dữ liệu tuần tự như LSTM, Transformer và Bộ biến
đổi hợp nhất theo thời gian có khả năng khai thác quan hệ theo thời gian giữa nhiều
nhóm đặc trưng. LSTM được thiết kế để học các phụ thuộc dài hạn trong chuỗi; Transformer
sử dụng cơ chế chú ý để mô hình hóa quan hệ giữa các phần tử; còn Bộ biến đổi hợp nhất
theo thời gian kết hợp xử lý tuần tự, lựa chọn biến và cơ chế chú ý để hỗ trợ dự báo đa
biến có khả năng giải thích \cite{Feng2018_181009936v2,Omranpour2024_241210540v1}.
Những đặc điểm này phù hợp với dữ liệu thị trường, trong đó lịch sử nhiều phiên, khối
lượng, độ biến động, chỉ báo kỹ thuật và thông tin tin tức có thể đồng thời liên quan
đến xu hướng của phiên giao dịch kế tiếp. Việc đặt mô hình chuỗi thời gian cạnh mô hình
văn bản cũng cho phép tách biệt câu hỏi về chất lượng dữ liệu với câu hỏi về độ phức
tạp của kiến trúc.

Thứ năm, các công trình đã khảo sát cho thấy những hướng kết hợp tin tức với dữ liệu
giá thường khác nhau về ngôn ngữ, quy mô dữ liệu, số lượng mã cổ phiếu, nhiệm vụ dự báo
và giao thức đánh giá \cite{Halder2022_221107392v1,Gu2024_240716150v1,
Chen2021_210708721v1,Omranpour2024_241210540v1}. Vì vậy, việc so sánh trực tiếp các
kết quả đã công bố không cho phép đưa ra kết luận ngay về hiệu quả của một tín hiệu
riêng lẻ. Đồ án lựa chọn kiểm soát các điều kiện về dữ liệu và đánh giá nhằm xác định
liệu đặc trưng cảm xúc có tạo ra giá trị bổ sung trong phạm vi nghiên cứu hay không.

Từ năm lý do trên, đề tài lựa chọn xây dựng hệ thống dự báo xu hướng biến động giá cổ
phiếu bằng cách kết hợp dữ liệu giá lịch sử với đặc trưng cảm xúc được trích xuất từ
tin tức tài chính tiếng Việt. Nghiên cứu không giả định trước rằng tin tức luôn giúp
cải thiện độ chính xác dự báo. Trọng tâm của nghiên cứu là xây dựng một quy trình thực
nghiệm có kiểm soát để xác định liệu tín hiệu cảm xúc có mang lại giá trị bổ sung hay
không. Nếu có, đề tài tiếp tục xem xét sự cải thiện xuất hiện trong những điều kiện nào
và có ổn định giữa các mã cổ phiếu hay không. Nếu không có cải thiện, kết quả vẫn có
ý nghĩa khi nguyên nhân được phân tích trên cơ sở dữ liệu và giao thức đánh giá hợp lệ.
\section{Mục tiêu nghiên cứu}
\label{sec:muc-tieu}

\subsection{Mục tiêu tổng quát}

Mục tiêu tổng quát của đề tài là xây dựng và đánh giá hệ thống dự báo xu hướng biến
động giá cổ phiếu trên thị trường chứng khoán Việt Nam thông qua việc kết hợp dữ liệu giá
lịch sử, chỉ báo kỹ thuật và đặc trưng cảm xúc trích xuất từ tin tức tài chính
tiếng Việt. Hệ thống được xây dựng để kiểm chứng mức độ hữu ích của tín hiệu cảm xúc
đối với bài toán dự báo phiên giao dịch kế tiếp, đồng thời phân tích các điều kiện mà nhóm
đặc trưng này có hoặc không tạo ra cải thiện so với mô hình chỉ sử dụng
dữ liệu thị trường. Mục tiêu của đề tài không phải tối ưu một con số hiệu năng duy nhất, mà là hình thành
một quy trình có thể kiểm tra từ dữ liệu đầu vào đến kết luận cuối cùng.

\subsection{Mục tiêu cụ thể}

Để trả lời các câu hỏi nghiên cứu, đồ án xác định bốn mục tiêu cụ thể có thể kiểm chứng:

\begin{enumerate}[nosep]
  \item Xây dựng quy trình thu thập, làm sạch, loại bỏ trùng lặp,
    chuẩn hóa và đồng bộ tin tức tài chính tiếng Việt với dữ liệu giá cổ phiếu theo ngày
    giao dịch; xây dựng hoặc hoàn thiện bộ dữ liệu có nhãn cảm xúc gồm ba lớp tích cực,
    trung tính và tiêu cực theo tiêu chí gán nhãn rõ ràng; tinh chỉnh PhoBERT, so sánh tầng phân loại
    tuyến tính với ít nhất một kiến trúc kết hợp như PhoBERT với CNN hoặc PhoBERT với
    BiLSTM, sau đó đánh giá khả năng phân biệt ba lớp trên cùng một tập dữ liệu.
  \item Xây dựng tập đặc trưng cảm xúc cho từng mã cổ phiếu và ngày
    giao dịch dựa trên phân phối xác suất dự đoán thay vì chỉ sử dụng nhãn có xác suất cao nhất;
    xây dựng mô hình dự báo xu hướng phiên kế tiếp bằng các mô hình chỉ sử dụng giá, LSTM
    và Transformer hoặc Bộ biến đổi hợp nhất theo thời gian khi đủ điều kiện; thiết kế cơ chế hợp
    nhất hai nhóm đặc trưng và so sánh mô hình kết hợp với đối chứng lớp phổ biến, đối chứng ngẫu nhiên, cùng các mô
    hình chỉ sử dụng dữ liệu giá và chỉ báo kỹ thuật.
  \item Thiết lập quy trình đánh giá theo thứ tự thời gian bằng phương pháp cuốn chiếu theo thời
    gian; kiểm soát rò rỉ trong quá trình chuẩn hóa, xử lý giá trị thiếu, xác định ngưỡng nhãn,
    lựa chọn đặc trưng và điều chỉnh tham số; phân tích xem tín hiệu cảm xúc có biểu hiện khác nhau
    giữa các giai đoạn thị trường, giữa ngày có nhiều tin và ngày ít tin, hoặc
    trong những ngày gắn với sự kiện.
  \item Xây dựng hệ thống thử nghiệm có khả năng tái lập để trực quan hóa
    kết quả; sử dụng cơ chế lựa chọn biến, cơ chế chú ý hoặc các phương pháp giải thích phù
    hợp nhằm nhận diện vai trò và giới hạn của đặc trưng cảm xúc trong quyết định dự
    báo.
\end{enumerate}
\section{Câu hỏi nghiên cứu}

Đề tài tập trung trả lời bốn câu hỏi nghiên cứu sau:

\begin{enumerate}[nosep]
  \item Mô hình PhoBERT được tinh chỉnh trên dữ liệu tin tức tài chính tiếng Việt có khả năng phân
    biệt ba lớp cảm xúc tích cực, trung tính và tiêu cực ở mức độ nào? Chất lượng phân loại thay đổi như thế nào
    khi so sánh tầng phân loại tuyến tính với một biến thể kết hợp trên cùng một tập dữ liệu?
  \item Việc bổ sung đặc trưng cảm xúc có cải thiện F1 vĩ mô, độ chính xác cân bằng và
    AUC-ROC một-chống-còn-lại so với mô hình chỉ sử dụng giá hay không? Nếu có, mức chênh lệch này có
    ổn định giữa các cửa sổ thời gian hay chỉ xuất hiện trong một số phần chia?
  \item Tín hiệu cảm xúc có khác biệt giữa các giai đoạn thị trường, giữa ngày
    có nhiều tin và ngày có ít tin, hoặc giữa ngày gắn với sự kiện và ngày thông thường hay
    không?
  \item Các cơ chế lựa chọn biến, cơ chế chú ý hoặc phương pháp giải thích phù hợp có thể
    hỗ trợ xác định vai trò của đặc trưng cảm xúc trong quyết định dự báo của mô hình
    hay không? Việc diễn giải các tín hiệu này có những giới hạn nào?
\end{enumerate}
\section{Đối tượng và phạm vi nghiên cứu}
\label{sec:pham-vi}

\subsection{Đối tượng nghiên cứu}

Đối tượng nghiên cứu của đề tài là mối quan hệ giữa cảm xúc được thể hiện trong tin tức
tài chính tiếng Việt và xu hướng biến động giá cổ phiếu trên thị trường chứng khoán Việt
Nam. Đề tài tập trung vào ba nhóm vấn đề: phương pháp biểu diễn và phân loại cảm xúc
trong văn bản tài chính, phương pháp mô hình hóa dữ liệu giá theo chuỗi thời gian, và cơ
chế hợp nhất hai nguồn dữ liệu để phục vụ bài toán dự báo.

Thứ nhất, đối với dữ liệu văn bản, đề tài nghiên cứu tiêu đề, nội dung nếu có thể thu
thập, thời điểm công bố, nguồn bài viết và mã cổ phiếu có liên quan. Kết quả xử lý văn bản
là xác suất thuộc ba lớp cảm xúc: tích cực, trung tính và tiêu cực. Nhãn cảm xúc được xác
định theo tác động kỳ vọng đối với doanh nghiệp hoặc cổ phiếu liên quan, thay vì chỉ dựa
trên sắc thái chung của câu văn.

Thứ hai, đối với dữ liệu thị trường, đề tài nghiên cứu giá mở cửa, giá cao nhất, giá
thấp nhất, giá đóng cửa, giá điều chỉnh khi có dữ liệu phù hợp, khối lượng giao dịch,
lợi suất và các chỉ báo kỹ thuật được tính theo ngày. Các biến này mô tả trạng thái đã
được quan sát tại thời điểm dự báo. Giá của phiên kế tiếp chỉ được sử dụng để tạo nhãn.

Thứ ba, đối với mô hình, đề tài nghiên cứu PhoBERT và các biến thể tầng phân loại cho
nhiệm vụ phân tích cảm xúc; nhóm mô hình chuỗi thời gian gồm mô hình cổ điển, LSTM và
Bộ biến đổi hợp nhất theo thời gian; cùng với kiến trúc hợp nhất đặc trưng phục vụ phân loại
xu hướng giá. PhoBERT được lựa chọn vì mô hình này được tiền huấn luyện cho tiếng Việt.
FinBERT chỉ được sử dụng làm điểm đối chiếu khái niệm nhằm đặt lựa chọn này trong dòng
nghiên cứu NLP tài chính quốc tế, không thay thế mô hình nền tiếng Việt của đề tài
\cite{Nguyen2020_200300744v3,Yang2020_200608097v2}.

\subsection{Phạm vi nghiên cứu}

\subsubsection{Phạm vi thị trường và mã cổ phiếu}

Phạm vi thị trường bao gồm Sở Giao dịch Chứng khoán Thành phố Hồ Chí Minh, chỉ số
VN-Index ở mức tham chiếu và một nhóm cổ phiếu thuộc rổ VN30. Việc lựa chọn cổ phiếu
ưu tiên các mã có thanh khoản tương đối cao, dữ liệu giá đầy đủ và xuất hiện thường
xuyên trên các nguồn tin tài chính. Theo đề cương, số lượng mã cuối cùng dự kiến nằm trong
khoảng 10 đến 30 mã và chỉ được xác định sau bước khảo sát chất lượng dữ liệu.

Trong cấu hình mức sàn hiện tại, rổ nghiên cứu được quản lý theo danh sách mã trong
bản kê. Danh sách mã và ngày bắt đầu sử dụng phải được lưu lại do thành phần VN30 có
thể thay đổi theo thời gian. Khi chia dữ liệu, các mã được căn chỉnh theo cùng một lịch
phiên để tránh sai lệch khi so sánh theo ngày giữa các mã.

Việc giới hạn số mã là một lựa chọn có chủ đích. Phạm vi này đủ để kiểm tra quy trình
trên nhiều mã và giảm mức độ phụ thuộc vào một trường hợp riêng lẻ, nhưng chưa đủ để đại
diện cho mọi nhóm ngành hoặc mọi mức thanh khoản trên HOSE. Vì vậy, kết quả phải được
diễn giải trong phạm vi rổ mã và khung thời gian cụ thể, không được xem là kết luận cho
toàn thị trường.

\subsubsection{Phạm vi thời gian và tần suất}

Dữ liệu dự kiến được sử dụng từ năm 2020 đến hết năm 2025, tùy theo ngày cuối có dữ liệu
hợp lệ của từng nguồn. Đơn vị quan sát là ngày giao dịch. Tại ngày $t$, hệ thống dự báo lớp
của lợi suất từ ngày $t$ đến phiên kế tiếp. Tần suất ngày phù hợp với mật độ tin tức và
giúp đề tài tránh phải xử lý thêm các vấn đề vi cấu trúc của dữ liệu trong ngày.

Trong trường hợp bản kê thực nghiệm mở rộng thời gian quan sát sau mốc đề cương,
mốc ngày cuối, nguồn dữ liệu và lý do thay đổi phải được ghi rõ cùng phiên bản dữ liệu.
Quy tắc này cho phép phân biệt phạm vi dự kiến với phạm vi thực tế đã được đo, qua đó
tránh đưa một khoảng thời gian chưa được kiểm tra vào kết luận.

Nguồn tin dự kiến gồm CafeF, Vietstock, VnEconomy và Stockbiz. Danh sách nguồn có thể
được điều chỉnh dựa trên điều kiện sử dụng, khả năng truy cập hợp lệ, mức độ đầy đủ
của nội dung và tính ổn định của cấu trúc trang. Mốc cắt được xác định theo giờ giao
dịch và múi giờ Việt Nam. Tin được công bố sau giờ đóng cửa, vào cuối tuần hoặc ngày nghỉ
được ánh xạ sang phiên giao dịch kế tiếp để bảo đảm thông tin chỉ được sử dụng sau thời điểm
thực sự có thể quan sát. Mốc thời gian bị thiếu hoặc mâu thuẫn không được tự động suy đoán.
Các bản ghi này phải được đánh dấu, loại khỏi phân tích chính hoặc báo cáo riêng.

\subsubsection{Phạm vi bài toán và giới hạn}

Bài toán dự báo được phát biểu dưới dạng phân loại ba lớp:

\begin{itemize}[nosep]
  \item UP: lợi suất của phiên kế tiếp lớn hơn ngưỡng tăng;
  \item FLAT: lợi suất nằm trong vùng trung tính;
  \item DOWN: lợi suất của phiên kế tiếp nhỏ hơn ngưỡng giảm.
\end{itemize}

Ngưỡng được ước lượng trên tập huấn luyện của từng cửa sổ, không sử dụng phân vị của toàn
bộ giai đoạn. Quy ước này ngăn phân phối lợi suất tương lai ảnh hưởng đến nhãn của dữ liệu
quá khứ. Đề tài không dự báo giá tuyệt đối, không nghiên cứu giao dịch trong ngày, sổ lệnh,
giao dịch tần suất cao, sản phẩm phái sinh hoặc khuyến nghị đầu tư. Việc đánh giá lợi
nhuận tích lũy, tỷ số Sharpe, mức sụt giảm tối đa và chi phí giao dịch không thuộc phạm
vi bắt buộc. Các nội dung này chỉ được thực hiện như phần mở rộng khi dữ liệu và tiến
độ cho phép.
\section{Phương pháp nghiên cứu}

Nghiên cứu sử dụng phương pháp định lượng để xây dựng và đánh giá hệ thống dự báo
xu hướng biến động giá cổ phiếu. Quy trình gồm hai nhánh. Nhánh PhoBERT xử lý và tổng hợp
xác suất cảm xúc từ tin tức tài chính tiếng Việt, trong khi nhánh chuỗi thời gian xử lý dữ liệu
giá và các chỉ báo kỹ thuật. Hai nhóm đặc trưng được căn chỉnh theo mã cổ phiếu, ngày giao
dịch và mốc cắt trước khi đưa vào mô hình dự báo. Các cấu hình được so sánh bao gồm đối chứng cơ sở,
mô hình chỉ sử dụng dữ liệu giá và LSTM hai nhánh; Bộ biến đổi hợp nhất theo thời gian được xem là phần mở
rộng. Hiệu năng được đánh giá bằng phương pháp cuốn chiếu theo thời gian, trong đó mọi phép
biến đổi và ngưỡng nhãn chỉ được học từ phần huấn luyện của từng cửa sổ.
\section{Đóng góp của nghiên cứu}

\subsection{Ý nghĩa khoa học}

Ý nghĩa của đề tài được xem xét trên hai phương diện. Về học thuật, đồ án kết hợp mô hình
ngôn ngữ tiếng Việt với mô hình chuỗi thời gian nhằm kiểm tra định lượng liệu cảm xúc có
cung cấp thông tin gia tăng cho dự báo xu hướng của phiên kế tiếp hay không. Việc tách
tầng phân loại cảm xúc khỏi tầng dự báo cho phép phân biệt lỗi phát sinh từ văn bản, lỗi
trong ánh xạ thời gian và lỗi của mô hình chuỗi thời gian. Do đó, kết quả có thể được
phân tích theo từng tầng thay vì chỉ dựa trên một điểm số cuối cùng.

Các mốc hiệu năng trong tài liệu nguồn, chẳng hạn kết quả phân loại trên một tập dữ
liệu cụ thể hoặc độ chính xác hướng trong nghiên cứu nước ngoài, chỉ được sử dụng làm
tham chiếu kỹ thuật. Các mốc này được đo trên ngôn ngữ, thị trường, khoảng thời gian và
phương pháp chia dữ liệu khác với đề tài, nên không được sử dụng làm ngưỡng đạt hoặc
cam kết về hiệu năng trên HOSE. Tương tự, cơ chế chú ý và mức độ quan trọng của biến
chỉ mô tả các tín hiệu mà mô hình đang sử dụng; các thông tin này không đủ để kết luận
rằng cảm xúc là nguyên nhân gây ra biến động giá.

Do mục tiêu của đề tài là kiểm tra giá trị gia tăng, mọi nhận định về tính hữu ích hoặc
không hữu ích của cảm xúc phải dựa trên cặp mô hình được huấn luyện và kiểm tra theo
cùng một giao thức. Một kết quả có sức thuyết phục cần báo cáo rõ chênh lệch điểm số,
độ biến thiên giữa các phần chia, khoảng tin cậy, phân bố lớp, phạm vi mã và các giới
hạn của dữ liệu.

\subsection{Ý nghĩa thực tiễn}

Về kỹ thuật, đồ án xây dựng một quy trình có khả năng tái lập, bao gồm thu thập dữ liệu,
kiểm tra chất lượng, ghi nhận và lưu trữ nguồn gốc, loại bỏ dữ liệu trùng lặp, xử lý theo
thời điểm, trích xuất xác suất cảm xúc, tổng hợp theo mã và ngày, huấn luyện, đánh giá
cuốn chiếu theo thời gian và lưu cấu hình. Quy trình này cũng tạo nền tảng để thay đổi
mô hình nền hoặc cơ chế hợp nhất trong khi vẫn duy trì khả năng so sánh với mức sàn đã
khóa.

Trong phạm vi được xác định, các đóng góp dự kiến gồm:

\begin{enumerate}[nosep]
  \item Một quy trình thực nghiệm tiếng Việt liên kết tin tức với dữ liệu giá theo mã và
    thời điểm, có cơ chế kiểm tra trùng lặp, nguồn gốc, khả năng truy xuất và mốc cắt.
  \item Một quy trình sử dụng PhoBERT để tạo phân phối xác suất cảm xúc thay vì chỉ giữ
    nhãn có xác suất cao nhất, sau đó tổng hợp tín hiệu theo mã và ngày.
  \item Một kiến trúc sàn gồm LSTM hai nhánh và cơ chế hợp nhất bằng phép nối, được đặt
    cạnh cấu hình chỉ sử dụng giá để đo lường phần đóng góp biên của cảm xúc.
  \item Một giao thức cuốn chiếu theo thời gian, trong đó tiền xử lý chỉ được thực hiện
    trên tập huấn luyện, gồm năm đoạn kiểm thử tách biệt, F1 vĩ mô và khoảng tin cậy
    được ước lượng bằng lấy mẫu lặp, qua đó giúp tách hiệu năng mô hình khỏi lỗi rò rỉ
    dữ liệu.
\end{enumerate}

Các đóng góp trên tập trung vào hệ thống thực nghiệm và khả năng kiểm chứng trong phạm
vi dữ liệu đã chọn. Đồ án không tuyên bố xây dựng mô hình tối ưu cho mọi thị trường,
không khẳng định quan hệ nhân quả giữa tin tức và giá, đồng thời không xem điểm số trên
dữ liệu nước ngoài là bằng chứng trực tiếp cho dữ liệu Việt Nam.

Đề tài không đặt mục tiêu chứng minh rằng cảm xúc luôn cải thiện mô hình. Trong dữ liệu
tài chính, tín hiệu văn bản có thể yếu, xuất hiện trễ, trùng với phản ứng đã được phản
ánh trong giá, đã được thị trường biết trước hoặc chỉ hữu ích trong một số giai đoạn.
Vì vậy, kết quả không cải thiện, hoặc chỉ cải thiện ở mức nhỏ nhưng không ổn định, vẫn
có giá trị nếu được thu nhận bằng một giao thức hợp lệ và được trình bày cùng các giới
hạn. Theo cách tiếp cận này, giá trị gia tăng của cảm xúc được xem là một giả thuyết cần
kiểm tra, thay vì một kết luận được giả định trước
\cite{Jiang2023_230602136v3,Siala2026_260200086v3}.

\chapter{TỔNG QUAN}
\label{ch:tong-quan}

Chương này tổng quan các công trình liên quan đến việc sử dụng tin tức tài chính để dự báo xu hướng biến động giá cổ phiếu. Nội dung được tổ chức theo ba hướng nghiên cứu chính. Hướng thứ nhất tập trung vào biểu diễn và phân loại cảm xúc trong văn bản tài chính. Hướng thứ hai xem xét việc mô hình hóa dữ liệu giá và các chỉ báo kỹ thuật dưới dạng chuỗi thời gian. Hướng thứ ba nghiên cứu phương pháp hợp nhất dữ liệu văn bản và dữ liệu thị trường trong mô hình dự báo đa phương thức. Bên cạnh các công trình quốc tế, chương này trình bày riêng các nghiên cứu sử dụng dữ liệu Việt Nam, do sự khác biệt về ngôn ngữ, nguồn tin, cấu trúc thị trường và cách tổ chức phiên giao dịch có thể ảnh hưởng đến kết quả thực nghiệm. Cuối chương, các kết quả khảo sát được tổng hợp thành khoảng trống nghiên cứu và cơ sở định vị đề tài.

Phần tổng quan không nhằm khẳng định một kiến trúc luôn vượt trội so với các kiến trúc khác. Một phương pháp có thể đạt kết quả cao trên thị trường, giai đoạn hoặc bộ dữ liệu được sử dụng trong công trình gốc, nhưng chưa chắc duy trì được ưu thế khi áp dụng cho dữ liệu tiếng Việt và thị trường HOSE. Vì vậy, mỗi công trình được xem xét theo dữ liệu sử dụng, nhiệm vụ dự báo, phương pháp xây dựng đặc trưng, phương pháp học, kết quả được báo cáo và các hạn chế còn tồn tại. Các kết quả định lượng từ những nghiên cứu trước chỉ được sử dụng làm căn cứ tham khảo về phương pháp, không được xem là bằng chứng trực tiếp về hiệu quả của mô hình trên dữ liệu của đề tài.
\section{Phân tích cảm xúc tin tức tài chính}
\label{sec:tq-sentiment}

Phân tích cảm xúc tin tức tài chính là quá trình xác định hướng tác động mà một bài
viết có thể tạo ra đối với doanh nghiệp, mã cổ phiếu hoặc toàn bộ thị trường. Trong
bối cảnh dự báo, cảm xúc không chỉ phản ánh thái độ tích cực hay tiêu cực được thể hiện
qua câu chữ. Một bài viết có thể sử dụng nhiều từ ngữ tích cực nhưng mô tả rủi ro đối
với doanh nghiệp, hoặc trình bày một sự kiện trung tính nhưng làm thay đổi kỳ vọng của
nhà đầu tư. Vì vậy, đề tài xác định nhãn theo tác động kỳ vọng trong ngắn hạn: tích cực
khi thông tin gợi ra triển vọng thuận lợi, tiêu cực khi thông tin cho thấy rủi ro hoặc
diễn biến bất lợi, và trung tính khi bài viết chủ yếu mô tả sự kiện nhưng chưa thể hiện
rõ hướng tác động.

\subsection{Từ mô hình ngôn ngữ tổng quát đến mô hình miền tài chính}

BERT đặt nền tảng cho phương pháp tiền huấn luyện mô hình ngôn ngữ trên lượng văn bản
lớn trước khi tinh chỉnh cho một nhiệm vụ cụ thể \cite{Devlin2019_bert}. Nhờ kiến trúc
Transformer và cơ chế tự chú ý, mô hình có thể tạo ra các biểu diễn phụ thuộc vào ngữ
cảnh thay vì gán ý nghĩa cố định cho từng từ. Khi được tinh chỉnh trên dữ liệu có nhãn,
các biểu diễn này có thể được sử dụng cho phân loại văn bản, nhận dạng thực thể hoặc
trích xuất thông tin. Tuy nhiên, mô hình ngôn ngữ tổng quát thường được huấn luyện trên
nhiều loại văn bản khác nhau, trong đó thuật ngữ tài chính chỉ chiếm một phần. Hạn chế
này thúc đẩy các nghiên cứu phát triển mô hình thích nghi với miền tài chính.

FinBERT là một đại diện tiêu biểu của hướng nghiên cứu này. Yang và cộng sự công bố
FinBERT như một mô hình ngôn ngữ được tiền huấn luyện và tinh chỉnh cho truyền thông
tài chính \cite{Yang2020_200608097v2}. Công trình hướng đến giải quyết sự khác biệt về
ngữ nghĩa giữa ngôn ngữ thông thường và ngôn ngữ tài chính. Trong lĩnh vực tài chính,
một từ có thể mang sắc thái khác nhau tùy thuộc vào chủ thể và chỉ tiêu được đề cập.
Chẳng hạn, ``tăng'' có thể mô tả doanh thu tăng, chi phí tăng hoặc rủi ro tăng. Do đó,
phương pháp phân loại chỉ dựa trên từ khóa khó phản ánh đầy đủ tác động của bản tin.
FinBERT cung cấp một mô hình nền phù hợp hơn cho các nhiệm vụ phân tích cảm xúc và dự
báo chuyển động trong những công trình sử dụng tin tức tiếng Anh.

Kết quả quan trọng của hướng FinBERT không chỉ nằm ở một giá trị độ chính xác cụ thể,
mà còn ở kết quả cho thấy việc thích nghi mô hình nền với miền dữ liệu có thể cải thiện
chất lượng biểu diễn khi dữ liệu gán nhãn có quy mô hạn chế. Đối với đề tài, FinBERT có
hạn chế là được xây dựng cho tiếng Anh và không trực tiếp xử lý các hiện tượng tách từ,
viết tắt, tên doanh nghiệp hoặc cách diễn đạt đặc thù trong tin tức tiếng Việt.

Đối với tiếng Việt, PhoBERT là mô hình đơn ngữ được tiền huấn luyện trên ngữ liệu tiếng
Việt và sử dụng bộ tách từ phù hợp với đặc trưng của ngôn ngữ này
\cite{Nguyen2020_200300744v3}. Sự khác biệt về phương pháp tách từ có ý nghĩa thực tế
vì một đơn vị ngữ nghĩa trong tiếng Việt thường gồm nhiều âm tiết được phân cách bằng
khoảng trắng. Bộ tách từ không phù hợp có thể làm phân mảnh tên riêng, thuật ngữ hoặc
mã cổ phiếu trước khi văn bản được đưa vào bộ mã hóa. Vì vậy, PhoBERT là một lựa chọn
hợp lý cho bước phân tích cảm xúc tiếng Việt.

Tuy nhiên, PhoBERT chỉ cung cấp nền tảng ngôn ngữ và không tự động bảo đảm rằng mô
hình hiểu chính xác miền tài chính. Tên mã, tên doanh nghiệp, số liệu tăng trưởng,
thông tin về phát hành cổ phiếu, thay đổi nhân sự và các câu phủ định có thể không xuất
hiện đầy đủ trong dữ liệu tiền huấn luyện. Do đó, mô hình vẫn cần được tinh chỉnh trên
tập dữ liệu tài chính với hướng dẫn gán nhãn rõ ràng. Kết quả của bước này cũng phải
được đánh giá riêng trước khi xác suất cảm xúc được sử dụng làm đặc trưng cho nhánh dự
báo giá.

\subsection{Các công trình cảm xúc và chuyển động cổ phiếu quốc tế}

Một hướng nghiên cứu gần với đề tài là trích xuất tín hiệu từ tin tức, sau đó đưa các
tín hiệu này vào mô hình chuỗi thời gian. Halder (2022) xây dựng quy trình
FinBERT-LSTM cho bài toán dự báo giá cổ phiếu từ cảm xúc của tin tức
\cite{Halder2022_221107392v1}. Tác giả mô tả quá trình phát triển từ các mô hình chuỗi
thời gian đơn giản đến LSTM, sau đó bổ sung FinBERT để phân tích cảm xúc tin tức trước
bước dự báo. Dữ liệu số cung cấp lịch sử giá, trong khi dữ liệu văn bản cung cấp thông
tin được trích xuất từ tin tức tài chính tiếng Anh. Các biểu diễn hoặc điểm cảm xúc
được kết hợp với dữ liệu giá trước khi đưa vào LSTM.

Kết quả chính của Halder cho thấy cấu hình FinBERT-LSTM đạt kết quả tốt hơn các cấu
hình chỉ sử dụng MLP, RNN hoặc LSTM trong phạm vi thí nghiệm của công trình. Nghiên cứu
này cung cấp một mẫu thiết kế gần với kiến trúc hai nhánh của đề tài: mô hình ngôn ngữ
xử lý tin tức, LSTM xử lý chuỗi thị trường, và hai nguồn dữ liệu được tích hợp trong
cùng một quy trình dự báo. Tuy nhiên, bài toán của Halder tập trung vào dự báo giá và
sử dụng dữ liệu tiếng Anh, thay vì phân loại ba lớp tăng, giảm và đi ngang của phiên kế
tiếp. Khi áp dụng ý tưởng này cho HOSE, cần kiểm soát sự khác biệt về thị trường, múi
giờ, lịch tin và cơ chế hình thành giá. Phương pháp ghép đặc trưng của công trình cũng
chưa giải quyết đầy đủ vấn đề tin tức được công bố trước hay sau giờ đóng cửa trong
bối cảnh của đề tài.

Gu và cộng sự (2024) tiếp tục phát triển hướng FinBERT-LSTM cho bài toán dự báo giá
\cite{Gu2024_240716150v1}. Công trình sử dụng dữ liệu giá cổ phiếu cùng với tin tức tài
chính tiếng Anh, trong đó FinBERT tạo biểu diễn hoặc điểm cảm xúc, còn LSTM học các
quan hệ theo thời gian. Nghiên cứu trình bày một phiên bản hoàn thiện hơn về mặt kỹ
thuật của quy trình, qua đó củng cố phương pháp tách bước xử lý văn bản khỏi bước mô
hình hóa chuỗi giá. Kết quả được công bố cho thấy việc bổ sung cảm xúc vào cấu hình
FinBERT-LSTM có thể cải thiện khả năng dự báo so với các cấu hình đối chứng trên bộ dữ
liệu riêng của công trình.

Hạn chế của nghiên cứu do Gu và cộng sự thực hiện là kết quả vẫn phụ thuộc vào dữ liệu
tiếng Anh, bài toán giá được lựa chọn và phương pháp ghép tin tức với phiên giao dịch.
Công trình dự báo giá tuyệt đối, không thực hiện nhiệm vụ phân loại xu hướng ba lớp như
đề tài. Ngoài ra, nếu quy trình chia dữ liệu, chuẩn hóa hoặc ghép mốc thời gian không
được mô tả theo cùng một chuẩn, chênh lệch giữa mô hình có và không có tin tức có thể
không phản ánh hoàn toàn giá trị của cảm xúc. Vì vậy, công trình cung cấp cơ sở tham
khảo về kiến trúc nhưng không thể được sử dụng để xác lập trước kỳ vọng hiệu năng đối
với cổ phiếu Việt Nam.

Chen (2021) nghiên cứu dự báo chuyển động cổ phiếu từ tin tức bằng biểu diễn ngữ cảnh
được tinh chỉnh từ BERT \cite{Chen2021_210708721v1}. Thay vì nén mỗi bài viết thành
một nhãn tích cực hoặc tiêu cực, phương pháp này giữ lại biểu diễn ngữ cảnh của tin tức,
sau đó tổng hợp theo đơn vị thời gian để đưa vào mô hình dự báo hướng biến động. Dữ liệu
của công trình gồm các cặp tin tức tài chính tiếng Anh và dữ liệu chuyển động giá trong
bối cảnh thị trường được nghiên cứu. Cách tiếp cận này cho phép mô hình khai thác những
sắc thái không thể hiện trong một điểm cảm xúc duy nhất.

Kết quả do Chen báo cáo cho thấy biểu diễn ngữ cảnh có thể hỗ trợ dự báo chuyển động;
công trình cũng xem xét giá trị của mô hình thông qua mô phỏng giao dịch. Khía cạnh này
có ý nghĩa về phương pháp vì hiệu năng phân loại không phải lúc nào cũng tương ứng với
giá trị kinh tế. Tuy nhiên, bộ dữ liệu và môi trường tiếng Anh khác đáng kể so với dữ
liệu tiếng Việt. Mô phỏng giao dịch còn phụ thuộc vào các giả định về thời điểm khớp
lệnh, phí giao dịch và thanh khoản, nên không được chuyển trực tiếp thành tiêu chí bắt
buộc trong mức sàn của đề tài. Trong phạm vi hiện tại, công trình chủ yếu gợi ý phương
pháp duy trì thông tin ngữ cảnh và phân tích kết quả, còn đặc trưng đầu vào của đề tài
vẫn ưu tiên vector xác suất cảm xúc dễ tái lập.

Chen và cộng sự (2019) áp dụng cách tiếp cận giàu thông tin hơn bằng việc trích xuất
các sự kiện chi tiết trong tin tức thay vì chỉ phân loại cảm xúc
\cite{Chen2019_191005078v1}. Các đơn vị thông tin được xem xét gồm sự kiện, chủ thể,
loại sự kiện và hướng tác động đối với cổ phiếu. Những tín hiệu ở mức sự kiện sau đó
được đưa vào mô hình dự báo chuyển động. Kết quả của công trình cho thấy biểu diễn sự
kiện có thể bổ sung cho dữ liệu giá vì hai bài viết có cùng sắc thái chung vẫn có thể
mô tả các sự kiện kinh tế khác nhau. Tuy nhiên, việc xác định lược đồ sự kiện và tạo
nhãn ở mức chi tiết đòi hỏi nhiều công sức hơn so với phân loại cảm xúc ba lớp. Sai sót
trong nhận dạng chủ thể hoặc thời điểm sự kiện cũng có thể lan truyền sang bước dự báo.
Do đó, tín hiệu ở mức sự kiện phù hợp với vai trò là một hướng mở rộng sau khi quy trình
phân tích cảm xúc cơ bản đã được kiểm soát.

Jiang và Zeng (2023) xem xét việc sử dụng FinBERT để phân tích cảm xúc tài chính và áp
dụng đầu ra vào dự báo chuyển động cổ phiếu \cite{Jiang2023_230602136v3}. Công trình
đại diện cho một quy trình gọn từ đầu đến cuối: văn bản tài chính được đưa qua FinBERT,
kết quả cảm xúc được liên kết với dữ liệu chuyển động, và mô hình dự báo được đánh giá
trên dữ liệu riêng của nghiên cứu. Kết quả công bố củng cố nhận định rằng cảm xúc miền
tài chính có thể được sử dụng như một nguồn tín hiệu cho bài toán cổ phiếu. Tuy nhiên,
mô hình, thị trường và ngôn ngữ của công trình không tương ứng với dữ liệu HOSE. Bên
cạnh đó, khả năng phân loại tốt của một mô hình FinBERT không tự động cho thấy xác suất cảm xúc có giá trị gia tăng sau khi lịch sử giá đã được sử dụng. Đề tài
cần kiểm tra vấn đề này bằng phương pháp loại bỏ đặc trưng.

Các nghiên cứu gần đây mở rộng quy trình thông qua mô hình ngôn ngữ lớn. Elahi và
Taghvaei (2024) nghiên cứu việc kết hợp dữ liệu tài chính có cấu trúc với bài viết bằng
LLM \cite{Elahi2024_241101368v1}. Dữ liệu số và dữ liệu văn bản được xử lý dưới các
dạng khác nhau, sau đó LLM được sử dụng để tạo hoặc hợp nhất thông tin phục vụ dự báo
chuyển động giá. Kết quả công bố cho thấy phương pháp hợp nhất bằng LLM có tiềm năng
khai thác nhiều thông tin ngữ cảnh hơn so với việc chỉ sử dụng một điểm cảm xúc. Tuy
nhiên, mức cải thiện phụ thuộc vào điểm kiểm, câu lệnh, thời điểm truy cập và bộ dữ liệu
của công trình. Nếu mô hình thay đổi phiên bản, kết quả khó được tái lập hoàn toàn; chi phí suy luận cũng có thể lớn đối với quy trình thực nghiệm trên nhiều mã.

Liang và cộng sự (2024) đề xuất FinGPT theo hướng làm giàu ngữ cảnh và xem xét quá
trình lan truyền của tin tức \cite{Liang2024_241210823v2}. Thay vì xem mỗi tin tức là
một quan sát độc lập, công trình quan tâm đến mức độ phổ biến của tin và mối liên hệ
giữa ngữ cảnh tin tức với dự báo biến động giá cổ phiếu dựa trên tín hiệu cảm xúc. Đây
là một hướng mở rộng hợp lý vì cùng một nội dung có thể tạo ra tác động khác nhau tùy
theo thời điểm và mức độ thị trường đã tiếp nhận thông tin. Kết quả của công trình cho
thấy việc bổ sung ngữ cảnh lan truyền là một hướng có thể cải thiện chất lượng biểu
diễn. Tuy nhiên, cách tiếp cận này đòi hỏi nhiều siêu dữ liệu hơn, mô hình lớn hơn và
quy trình kiểm soát dữ liệu phức tạp hơn mức sàn của đề tài.

Siala, Khanfir và Papadakis (2026) xem xét trực tiếp mức cải thiện mà cảm xúc do LLM
tạo ra có thể mang lại cho dự báo chuyển động giá \cite{Siala2026_260200086v3}. Công
trình không giả định rằng tin tức luôn hữu ích, mà so sánh dự báo có tín hiệu cảm xúc
với dự báo đối chứng để đo lường lượng thông tin gia tăng. Kết quả được công bố trong
nghiên cứu mới cho thấy tác động của cảm xúc LLM cần được đo lường trên từng khung dữ
liệu và không nên được suy rộng thành lợi ích ổn định cho mọi thị trường. Tuy nhiên,
đây là một bản thảo tiền công bố mới; dữ liệu và giao thức không trùng với đề tài, đồng
thời việc sử dụng LLM làm tăng yêu cầu về chi phí và khả năng tái lập. Quan điểm kiểm
tra đóng góp biên của cảm xúc cung cấp cơ sở trực tiếp cho thiết kế loại bỏ đặc trưng
của đề tài.

\subsection{Khả năng áp dụng cho dữ liệu tiếng Việt}

Các tài liệu đã được xác minh cung cấp cơ sở cho việc lựa chọn mô hình ngôn ngữ, phương
pháp biểu diễn tin tức và thiết kế mô hình kết hợp. Tuy nhiên, kết quả của các công
trình không thể được chuyển trực tiếp sang dữ liệu của đồ án do khác biệt về ngôn ngữ,
nguồn dữ liệu, nhiệm vụ dự báo và giao thức đánh giá. Vì vậy, các tài liệu này được sử
dụng để xây dựng cơ sở phương pháp, không thay thế kết quả thực nghiệm của đề tài.

Đối với tiếng Việt, PhoBERT là mô hình ngôn ngữ đơn ngữ được tiền huấn luyện trên ngữ
liệu tiếng Việt và có bộ tách từ phù hợp với đặc trưng của ngôn ngữ này
\cite{Nguyen2020_200300744v3}. Đặc điểm này phù hợp với mục tiêu xử lý tên doanh nghiệp,
mã cổ phiếu và các cách biểu đạt thường gặp trong tin tức tài chính. Tuy nhiên, việc
lựa chọn PhoBERT chỉ cung cấp cơ sở ngôn ngữ ban đầu; chất lượng phân loại vẫn phải
được đánh giá trên dữ liệu có nhãn của đề tài.

Trong miền tiếng Anh, FinBERT được phát triển bằng cách thích nghi BERT với ngữ liệu
tài chính \cite{Yang2020_200608097v2}. Các nghiên cứu sử dụng FinBERT hoặc biểu diễn
BERT để kết hợp tin tức với dữ liệu giá cho thấy có thể tách bước xử lý văn bản khỏi
bước mô hình hóa chuỗi \cite{Halder2022_221107392v1,Gu2024_240716150v1,
Chen2021_210708721v1}. Cách phân tách này gợi ý thiết kế hai nhánh của đồ án, trong đó
nhánh văn bản tạo đặc trưng cảm xúc, còn nhánh giá học các quan hệ theo thời gian.

Một hướng mở rộng là duy trì biểu diễn ngữ cảnh hoặc thông tin sự kiện thay vì nén mỗi
bài viết thành một nhãn cảm xúc duy nhất. Chen và cộng sự đề xuất khai thác các sự kiện
chi tiết trong tin tức \cite{Chen2019_191005078v1}, trong khi các nghiên cứu khác xem
xét biểu diễn ngữ cảnh và mô hình ngôn ngữ lớn cho dự báo chuyển động cổ phiếu
\cite{Chen2021_210708721v1,Elahi2024_241101368v1,Liang2024_241210823v2}. Các hướng
tiếp cận này có thể duy trì nhiều thông tin hơn nhưng đồng thời làm tăng yêu cầu về dữ
liệu, chi phí tính toán và khả năng tái lập. Vì vậy, ở mức sàn, đồ án sử dụng xác suất
của ba lớp cảm xúc thay vì mặc định áp dụng một kiến trúc ngôn ngữ lớn và phức tạp.

Các công trình về mô hình hóa chuỗi cũng cho thấy sự khác biệt giữa các kiến trúc cần được đánh giá cùng với dữ liệu và giao thức. LSTM, Transformer và Bộ biến
đổi hợp nhất theo thời gian lần lượt cung cấp các phương pháp khác nhau để mô hình hóa
phụ thuộc theo thời gian, quan hệ giữa nhiều biến và cơ chế lựa chọn biến
\cite{Hochreiter1997_lstm,Vaswani2017_attention,Lim2021_tft}. Các kiến trúc đa phương
thức hoặc mô hình hóa quan hệ giữa nhiều mã tiếp tục mở rộng khả năng biểu diễn
\cite{Omranpour2024_241210540v1,Boyle2023_230503835v1,Xu2021_211013716v2,
Zhao2023_230609862v3}. Tuy nhiên, độ phức tạp cao hơn không tự động tạo ra kết quả tốt
hơn nếu dữ liệu, số lượng mẫu hoặc quy tắc đánh giá chưa phù hợp.

Khoảng trống mà đồ án tập trung kiểm tra không dựa trên giả định rằng mọi tín hiệu cảm
xúc đều có giá trị dự báo. Vấn đề cần giải quyết là liệu đặc trưng cảm xúc được trích
xuất từ tin tức tiếng Việt có cung cấp thông tin bổ sung sau khi lịch sử giá và chỉ báo
kỹ thuật đã được sử dụng hay không. Để trả lời câu hỏi này, hai cấu hình có và không có
cảm xúc phải được xây dựng trong cùng điều kiện về cửa sổ dữ liệu, tiền xử lý, hạt giống,
ngân sách huấn luyện và tập đánh giá. Kết luận chỉ được rút ra từ kết quả ngoài mẫu
trên dữ liệu của đề tài.

\subsection{Đánh giá chất lượng nhãn và đầu ra cảm xúc}

Đánh giá cảm xúc cần được thực hiện riêng với đánh giá dự báo xu hướng. Ở tầng phân
loại văn bản, mô hình được đánh giá bằng độ chính xác, độ chính xác dương tính, độ bao
phủ và F1 theo từng lớp. F1 vĩ mô nên được ưu tiên khi số lượng mẫu tích cực, trung
tính và tiêu cực không cân bằng. Độ tin cậy của bộ nhãn được kiểm soát thông qua hướng
dẫn thống nhất, lưu vết phiên bản, kiểm tra nhãn thiếu và rà soát thủ công toàn bộ mẫu
bởi một người. Nếu mô hình ngôn ngữ lớn được sử dụng để tạo nhãn sơ bộ, các nhãn này
chỉ đóng vai trò đầu vào hỗ trợ và không thay thế nhãn đã được con người kiểm tra;
không tính độ đồng thuận liên người gán nhãn.

Đầu ra chuyển sang nhánh dự báo cũng không nên chỉ gồm nhãn có xác suất cao nhất. Với
mỗi bài viết, mô hình tạo vector:
\[
  \mathbf{p}_{i,j,t}=(p_{\mathrm{NEG}},p_{\mathrm{NEU}},p_{\mathrm{POS}}),
  \qquad \sum_c p_c=1,
\]
trong đó $i$ là mã cổ phiếu, $j$ là bài viết và $t$ là ngày được ánh xạ. Vector này
duy trì mức độ không chắc chắn của mô hình. Các vector được tổng hợp cùng số tin,
độ phân tán, tỷ lệ lớp và cờ có tin theo từng mã và ngày. Nhờ đó, ngày có nhiều tin
trung tính được phân biệt với ngày không có tin, trong khi ngày có một bài tích cực cao
được phân biệt với ngày có nhiều bài mang tín hiệu trái chiều.
\section{Dự báo xu hướng giá bằng mô hình chuỗi thời gian}
\label{sec:tq-timeseries}

Nhánh chuỗi thời gian biểu diễn trạng thái thị trường thông qua dữ liệu OHLCV, lợi suất và các chỉ báo kỹ thuật \cite{Wilder1978_rsi,Appel2005_macd}. Mục tiêu của nhánh này là xây dựng một đối chứng cơ sở đủ mạnh để đánh giá liệu cảm xúc có cung cấp thông tin mới hay chỉ phản ánh lại thông tin đã tồn tại trong giá. Do dữ liệu thị trường có thứ tự thời gian và chịu ảnh hưởng của sự thay đổi chế độ, quy trình đánh giá phải bảo toàn quan hệ trước sau. Các nghiên cứu về dự báo xu hướng thường bắt đầu với lịch sử giá, sau đó mở rộng sang các mô hình sâu hoặc các mô hình biểu diễn quan hệ giữa nhiều mã \cite{Feng2018_181009936v2,Xu2021_211013716v2}.

\subsection{Đối chứng cơ sở chỉ giá và các mô hình cổ điển}

Đối chứng lớp phổ biến luôn dự đoán lớp xuất hiện nhiều nhất trong tập huấn luyện. Đối chứng ngẫu nhiên đưa ra dự đoán theo phân bố lớp của tập huấn luyện với hạt giống cố định. Hai mốc đối chứng này cho biết mức hiệu năng có thể đạt được mà không cần học các quan hệ phức tạp. Nếu mô hình sâu chỉ vượt đối chứng lớp phổ biến về độ chính xác nhưng có F1 vĩ mô thấp hơn, kết luận rằng mô hình tốt hơn cần được xem xét lại. Các đối chứng cơ sở này cũng hỗ trợ phát hiện trường hợp phân bố nhãn bị lệch quá mức hoặc quy tắc tạo nhãn làm cho một lớp chiếm ưu thế.

Nhóm mô hình cổ điển gồm hồi quy logistic, rừng ngẫu nhiên và XGBoost, với đầu vào chỉ bao gồm các đặc trưng giá và chỉ báo kỹ thuật. Hồi quy logistic cung cấp một mốc tuyến tính dễ kiểm tra; rừng ngẫu nhiên mô hình hóa các ngưỡng và tương tác phi tuyến; XGBoost xây dựng tuần tự nhiều cây nhằm giảm sai số. Các mô hình này không trực tiếp xử lý thứ tự của toàn bộ cửa sổ, nhưng vẫn có thể tạo ra tín hiệu mạnh nếu việc tạo đặc trưng đã tóm tắt đủ thông tin. Việc so sánh các mô hình này với LSTM giúp phân biệt mức cải thiện do đặc trưng với mức cải thiện do kiến trúc.

\subsection{LSTM cho phụ thuộc ngắn hạn và dài hạn}

LSTM xử lý chuỗi theo thứ tự thời gian, đồng thời duy trì trạng thái ẩn và trạng thái ô nhớ \cite{Hochreiter1997_lstm}. Các cổng vào, cổng quên và cổng ra cho phép mô hình quyết định thông tin cần giữ lại hoặc loại bỏ qua nhiều bước thời gian. Với cửa sổ giá, LSTM có thể học sự thay đổi của lợi suất, khối lượng và các chỉ báo qua những ngày liên tiếp. Với cửa sổ cảm xúc, LSTM có thể học ảnh hưởng của một chuỗi tín hiệu được tổng hợp theo ngày. Vì vậy, LSTM thường được sử dụng làm bộ mã hóa chung trong các quy trình kết hợp tin tức và giá \cite{Halder2022_221107392v1,Gu2024_240716150v1}.

Ưu điểm của LSTM là cấu trúc tương đối gọn và phù hợp khi quy mô dữ liệu không đủ lớn để huấn luyện một Transformer có nhiều tham số. Tuy nhiên, mô hình xử lý dữ liệu tuần tự nên khó xem trực tiếp quan hệ giữa các thời điểm xa, đồng thời có thể nhạy với độ dài cửa sổ và phương pháp chuẩn hóa. Do đó, LSTM trong đề tài được sử dụng làm mức sàn thực nghiệm và không được giả định là mô hình tối ưu.

Trong đề tài, LSTM chỉ giá và LSTM hai nhánh phải sử dụng cùng cửa sổ, phương pháp chia dữ liệu, ngân sách huấn luyện và tiêu chí dừng. Mô hình hai nhánh xử lý riêng chuỗi giá và chuỗi cảm xúc trước khi hợp nhất, nhưng số bước quan sát và thông tin được phép nhìn thấy phải tương đương với mô hình chỉ giá. Điều kiện này bảo đảm chênh lệch giữa hai mô hình chủ yếu phản ánh tác động của việc bổ sung cảm xúc, thay vì xuất phát từ việc một mô hình được cung cấp lịch sử dài hơn hoặc có nhiều cơ hội điều chỉnh hơn.

\subsection{Transformer cho chuỗi thời gian}

Transformer sử dụng cơ chế chú ý để kết nối các vị trí trong chuỗi \cite{Vaswani2017_attention}. So với cơ chế truyền trạng thái tuần tự, cơ chế chú ý cho phép mô hình xem xét trực tiếp các thời điểm cách xa nhau và nhiều nhóm biến. Đối với dữ liệu giá, khả năng này hữu ích khi một biến động hiện tại có liên quan đến một mẫu đã xuất hiện từ nhiều phiên trước. Tuy nhiên, cơ chế chú ý không tự loại bỏ rò rỉ dữ liệu và cũng không tự chứng minh quan hệ nhân quả. Mô hình vẫn phải được huấn luyện theo thứ tự thời gian và kiểm tra trong điều kiện ngoài mẫu.

Boyle và Kalita (2023) nghiên cứu Spatiotemporal Transformer để dự báo chuyển động cổ phiếu \cite{Boyle2023_230503835v1}. Công trình mô hình hóa đồng thời chiều thời gian và chiều không gian, trong đó chiều không gian biểu diễn quan hệ giữa nhiều cổ phiếu. Phương pháp này phù hợp với giả thuyết rằng một sự kiện hoặc biến động ở một mã có thể liên quan đến các mã khác trong cùng ngành hoặc cùng thị trường. Kết quả của công trình cho thấy Transformer không gian thời gian có thể khai thác cả hai loại phụ thuộc thay vì xử lý từng mã như một chuỗi độc lập. Hạn chế của phương pháp là yêu cầu xác định hoặc học quan hệ liên mã, qua đó làm tăng yêu cầu về dữ liệu và thiết kế đánh giá. Các quan hệ được học trên thị trường trong nghiên cứu cũng chưa có bằng chứng trực tiếp trên HOSE, vì vậy đề tài chỉ xem đây là một hướng tham khảo trong tương lai.

Omranpour, Rabusseau và Rabbany (2024) đề xuất Transformer bậc cao cho dữ liệu chuỗi thời gian đa phương thức \cite{Omranpour2024_241210540v1}. Thay vì chỉ ghép hai vector tại cuối mạng, hướng này tìm cách mô hình hóa tương tác bậc cao giữa nhiều chuỗi hoặc nhiều phương thức trong cùng một bộ Transformer. Về mặt ý tưởng, kiến trúc này gần với trường hợp kết hợp chuỗi giá, chỉ báo và chuỗi cảm xúc. Kết quả được công bố cho thấy Transformer bậc cao có thể hỗ trợ bài toán dự báo biến động giá cổ phiếu trên dữ liệu đa phương thức. Hạn chế của phương pháp là số lượng tham số, yêu cầu tính toán và mức độ phức tạp trong kiểm soát quá khớp đều cao hơn so với mô hình chỉ sử dụng phép nối. Vì vậy, mô hình này phù hợp với hướng phát triển sau khi mức sàn LSTM hai nhánh đã được đánh giá.

Bộ biến đổi hợp nhất theo thời gian, thường được gọi là TFT, là một tầng mở rộng kết hợp xử lý chuỗi, lựa chọn biến, các khối có cổng và cơ chế chú ý. Cấu trúc này có thể sử dụng biến tĩnh, biến được biết trước và biến quan sát theo thời gian, đồng thời tạo ra các biểu diễn phục vụ dự báo. Nếu được áp dụng trong đề tài, TFT phải được cấu hình cho bài toán phân loại ba lớp, có kích thước đầu ra bằng 3 và sử dụng hàm mất mát phân loại. Tầm quan trọng của biến hoặc chú ý chỉ cho biết mô hình sử dụng biến như thế nào trong một điều kiện cụ thể, không đủ để suy ra rằng cảm xúc gây ra biến động giá \cite{Lim2021_tft}.

Trong phạm vi đồ án, TFT chỉ được xem là mô hình mở rộng sau khi cấu hình LSTM hai nhánh đã được đánh giá. Việc sử dụng TFT phải giữ nguyên nhiệm vụ phân loại ba lớp, phạm vi dữ liệu, tập đánh giá và tiêu chí so sánh. Do đó, kết quả của các nghiên cứu về TFT trên nhiệm vụ dự báo giá hoặc trên các thị trường khác không được chuyển trực tiếp thành kết luận của đồ án.

\subsection{Dịch chuyển phân phối và quan hệ giữa các mã}

Một nhóm nghiên cứu khác không xem mỗi cổ phiếu là một chuỗi hoàn toàn độc lập. HIST của Xu và cộng sự (2021) khai thác thông tin chung giữa các mã thông qua các khái niệm được xác định trước và các khái niệm ẩn \cite{Xu2021_211013716v2}. Mô hình kết hợp thông tin riêng của từng mã với thông tin được chia sẻ, đồng thời cho phép mức độ liên quan giữa mã và khái niệm thay đổi theo thời gian. Dữ liệu được tổ chức theo các chuỗi chuyển động của nhiều cổ phiếu và các quan hệ khái niệm. Kết quả công bố cho thấy HIST đạt kết quả đầu tư cao hơn các đối chứng cơ sở trên bộ dữ liệu của công trình. Hạn chế của phương pháp là yêu cầu thông tin về quan hệ giữa các mã và cách định nghĩa khái niệm, trong khi đề tài hiện tập trung vào dữ liệu của 10 mã với quy trình độc lập để thuận tiện cho việc kiểm soát.

REST của Xu và cộng sự (2021) tiếp cận bài toán dự báo xu hướng theo sự kiện \cite{Xu2021_210207372v2}. Tác động của một sự kiện được điều kiện hóa theo trạng thái của cổ phiếu, sau đó lan truyền qua đồ thị gồm các mã có liên quan. Phương pháp này phù hợp với nhận định rằng một tin tức có thể ảnh hưởng đến nhiều doanh nghiệp và mức độ ảnh hưởng phụ thuộc vào bối cảnh của từng mã. Kết quả của REST cho thấy việc đưa quan hệ sự kiện và quan hệ giữa các mã vào mô hình có thể cải thiện hiệu quả dự báo trên bộ dữ liệu được nghiên cứu. Hạn chế là phương pháp yêu cầu cấu trúc sự kiện, đồ thị và dữ liệu đa mã đủ dày, đồng thời khó tách biệt tác động của tin tức khỏi tác động của quan hệ liên mã.

DoubleAdapt của Zhao, Kong và Shen (2023) tập trung xử lý trôi dạt khái niệm trong dữ liệu liên tục \cite{Zhao2023_230609862v3}. Phương pháp sử dụng một bộ thích nghi dữ liệu để ánh xạ dữ liệu mới về một phân phối cục bộ gần ổn định hơn, cùng với một bộ thích nghi mô hình để khởi tạo các tham số phù hợp với giai đoạn mới. Cơ chế học siêu tham số này hướng đến việc cập nhật mô hình khi quan hệ giữa đặc trưng và xu hướng thay đổi. Kết quả công bố cho thấy DoubleAdapt có thể cải thiện khả năng dự báo trong bối cảnh phân phối thay đổi so với các phương pháp cập nhật thông thường. Hạn chế của phương pháp là thiết kế phức tạp và yêu cầu một giao thức liên tục rõ ràng. Trong phạm vi đề tài, trôi dạt khái niệm được ghi nhận là một rủi ro cần phân tích, nhưng chưa phải thành phần bắt buộc của mức sàn.

MTMD của Wang và cộng sự (2022) sử dụng bộ nhớ thời gian đa quy mô và cơ chế giảm thiên lệch cho bài toán dự báo xu hướng cổ phiếu \cite{Wang2022_221208656v2}. Mô hình học các biểu diễn bộ nhớ ở nhiều thang thời gian, sử dụng cơ chế chú ý để nhận diện các điểm tương đồng và kết hợp tín hiệu cục bộ với tín hiệu toàn cục thông qua cấu trúc đồ thị. Mục tiêu của mô hình là duy trì tính nhất quán theo thời gian và giảm ảnh hưởng của nhiễu. Kết quả của công trình cho thấy việc kết hợp bộ nhớ đa quy mô với quan hệ giữa các mã có thể vượt qua các đối chứng cơ sở trước đó trên dữ liệu gốc. Hạn chế là mô hình có nhiều thành phần, khó tái lập trong một đồ án có dữ liệu hạn chế và khó xác định phần cải thiện đến từ bộ nhớ hay đồ thị.

Liu và cộng sự (2025) đề xuất DishFT-GNN, một mạng đồ thị có nhận biết tương lai dựa trên chưng cất mô hình \cite{Liu2025_250210776v1}. Mô hình giáo viên học tương quan giữa sự dịch chuyển của phân phối lịch sử và phân phối tương lai, sau đó truyền thông tin cho mô hình học sinh để tạo ra biểu diễn không gian thời gian có khả năng dự báo tốt hơn. Phương pháp này hướng đến xử lý sự thay đổi của phân phối thị trường, một vấn đề thường bị bỏ qua khi mô hình chỉ được huấn luyện trên một phân phối cố định. Kết quả được công bố cho thấy chưng cất mô hình có nhận biết tương lai có thể hỗ trợ dự báo xu hướng. Hạn chế là cấu trúc giáo viên, học sinh và cơ chế quan sát tương lai phải được thiết kế cẩn thận để tránh biến thông tin tương lai thành rò rỉ trong quá trình huấn luyện.

Lu, Hu và Zhang (2026) giới thiệu S$^{3}$G, một đồ thị không gian trạng thái cổ phiếu cho bài toán dự báo xu hướng \cite{Lu2026_260324236v1}. Công trình kết hợp khử nhiễu bằng wavelet, đồ thị thay đổi theo thời gian và mô hình không gian trạng thái để theo dõi sự tiến hóa của quan hệ giữa các mã. Trong S$^{3}$G, quan hệ không được giả định là tĩnh mà phụ thuộc vào dữ liệu và được cập nhật theo trạng thái thị trường. Kết quả được báo cáo cho thấy phương pháp này đạt kết quả tốt trên các bộ dữ liệu quốc tế được sử dụng trong nghiên cứu. Hạn chế là phương pháp đòi hỏi lượng lớn dữ liệu liên mã, nhiều bước huấn luyện và không trực tiếp xử lý cảm xúc từ tin tức tiếng Việt. Đây là một bối cảnh nghiên cứu mới và không phải đối thủ trực tiếp của quy trình hai nhánh trong phạm vi đề tài.

Nhìn chung, các nghiên cứu HIST, REST, DoubleAdapt, MTMD, DishFT-GNN và S$^{3}$G cho thấy các phương pháp dự báo xu hướng hiện đại đang mở rộng từ chuỗi đơn sang mô hình hóa quan hệ liên mã và thích nghi với sự thay đổi phân phối. Tuy nhiên, việc bổ sung đồ thị, bộ nhớ và học siêu tham số làm tăng khó khăn trong việc bảo đảm so sánh công bằng và truy nguyên nguyên nhân của mức cải thiện. Vì vậy, đề tài trước hết lựa chọn mô hình hai nhánh cho từng mã, đặt các mã trên cùng lịch cuốn chiếu theo thời gian, đồng thời dành quan hệ liên mã và thích nghi trực tuyến cho hướng phát triển tiếp theo.
\section{Kết hợp tin tức và giá (đa phương thức)}
\label{sec:tq-fusion}

Hợp nhất đa phương thức là quá trình kết hợp thông tin từ văn bản và chuỗi số trong
một mô hình dự báo. Hai nguồn dữ liệu này khác nhau về kích thước, tần suất, thời điểm
xuất hiện và mức độ tin cậy. Giá được quan sát theo phiên, trong khi tin tức có thể xuất
hiện trước giờ mở cửa, trong phiên, sau giờ đóng cửa hoặc vào ngày nghỉ. Do đó, cách ghép đặc trưng không chỉ là phép nối vector, mà còn là quyết định về mốc cắt, ngày ánh xạ và cách xử lý ngày không có tin. Các hướng nghiên cứu quốc tế có thể được phân thành các phương pháp sử
dụng điểm cảm xúc, biểu diễn ngữ cảnh, sự kiện chi tiết hoặc LLM
\cite{Halder2022_221107392v1,Elahi2024_241101368v1}.

\subsection{Hợp nhất bằng điểm cảm xúc}

Trong cách tiếp cận đơn giản nhất, mỗi bài viết được chuyển thành nhãn hoặc vector xác
suất. Các vector sau đó được tổng hợp theo mã và ngày rồi nối với đặc trưng kỹ thuật.
Đối với bài viết $j$ của mã $i$ trong ngày $t$, vector xác suất có thể được tổng hợp
bằng trung bình, trung vị, tổng có trọng số hoặc các thống kê mô tả số lượng tin. Phương
pháp này dễ kiểm tra, thuận lợi cho việc thực hiện loại bỏ đặc trưng và phù hợp với
dữ liệu có quy mô không quá lớn. Việc lưu vector thay vì chỉ lưu nhãn cũng bảo toàn
thông tin về mức độ không chắc chắn của bộ phân loại.

Hạn chế của phương pháp này là nhiều bài viết có thể bị nén thành một giá trị trung
bình, dẫn đến mất thông tin về thứ tự bài viết, sự khác biệt giữa các nguồn, độ tin cậy,
độ dài và mối quan hệ giữa sự kiện với trạng thái giá. Nếu ngày không có tin được thay
bằng vector zero mà không bổ sung cờ thiếu dữ liệu, mô hình có thể không phân biệt
``không có tin'' với ``có nhiều tin trung tính''. Ngược lại, nếu sử dụng giá trị từ
tương lai để điền dữ liệu hoặc gộp bài viết sau giờ đóng cửa vào cùng phiên, quá trình
hợp nhất có thể gây rò rỉ dữ liệu.

FinBERT-LSTM của Halder và dòng công trình của Gu là các ví dụ điển hình cho quy trình cảm xúc rồi chuỗi thời gian \cite{Halder2022_221107392v1,Gu2024_240716150v1}.
Đề tài kế thừa cách tách các nguồn dữ liệu và sử dụng PhoBERT thay cho FinBERT để xử lý
tin tiếng Việt. Tuy nhiên, một mô hình chỉ có ý nghĩa khi được so sánh với mô hình chỉ giá trên cùng dữ liệu. Chất lượng phân loại cảm xúc cao
không bảo đảm rằng vector cảm xúc được tổng hợp theo ngày sẽ làm tăng F1 vĩ mô của
nhánh dự báo.

\subsection{Hợp nhất bằng biểu diễn ngữ cảnh và sự kiện}

Một hướng tiếp cận khác duy trì biểu diễn của tiêu đề hoặc bài viết thay vì nén văn
bản thành một điểm phân cực. Chen sử dụng biểu diễn ngữ cảnh từ BERT cho bài toán dự
báo biến động giá cổ phiếu và cho thấy biểu diễn này có thể chứa thông tin ngoài nhãn
cảm xúc \cite{Chen2021_210708721v1}. Khi được đưa vào mô hình RNN, biểu diễn này có thể
giúp mô hình phân biệt các câu có cùng sắc thái nhưng đề cập đến những loại sự kiện
khác nhau. Tuy nhiên, các vector có kích thước lớn, yêu cầu nhiều bộ nhớ và có thể chứa
thông tin không liên quan đến mục tiêu dự báo.

Một hướng tiếp cận giàu thông tin hơn là trích xuất sự kiện chi tiết, bao gồm chủ thể,
loại sự kiện, hướng tác động và thời gian. Chen và cộng sự nghiên cứu việc đưa tín hiệu
sự kiện chi tiết vào dự báo chuyển động \cite{Chen2019_191005078v1}. Phương pháp này có
thể hạn chế mất mát thông tin khi một bài viết chứa nhiều sự kiện, nhưng đòi hỏi lược
đồ, bộ gán nhãn và các quy tắc ánh xạ sự kiện bổ sung. Một sai sót nhỏ khi xác định chủ
thể, chẳng hạn nhầm công ty được mua với công ty đi mua, có thể làm sai cả đặc trưng
và nhãn đánh giá.

Do đề tài hướng đến một quy trình có thể tái lập trong điều kiện hạn chế về dữ liệu và
thời gian, cấu hình cơ sở sử dụng xác suất cảm xúc được tổng hợp theo mã và ngày. Biểu
diễn toàn văn và tín hiệu ở mức sự kiện được xem là các hướng phân tích mở rộng. Lựa
chọn này nhằm tách biệt câu hỏi nền tảng về giá trị của cảm xúc với câu hỏi phức tạp
hơn về biểu diễn văn bản, đồng thời không phủ nhận khả năng của các biểu diễn hoặc sự
kiện chi tiết.

\subsection{Hợp nhất động và LLM}

Các nghiên cứu sử dụng LLM nhằm duy trì nhiều thông tin ngữ cảnh hơn trong quá trình
xử lý tin tức. Elahi và Taghvaei kết hợp dữ liệu tài chính với bài viết thông qua LLM,
qua đó đưa việc hợp nhất dữ liệu có cấu trúc và văn bản vào cùng một quy trình
\cite{Elahi2024_241101368v1}. Lợi thế của hướng tiếp cận này là mô hình có thể xử lý
điều kiện, chủ thể và quan hệ giữa các sự kiện thay vì chỉ đếm các từ tích cực hoặc
tiêu cực. Tuy nhiên, kết quả phụ thuộc vào điểm kiểm mô hình, câu lệnh, giới hạn ngữ
cảnh và thiết kế đầu ra. Nếu phiên bản và câu lệnh không được cố định, cùng một bộ dữ
liệu có thể tạo ra các đặc trưng khác nhau giữa những lần chạy.

Liang và cộng sự đưa yếu tố lan truyền và ngữ cảnh vào FinGPT
\cite{Liang2024_241210823v2}. Hướng nghiên cứu này nhấn mạnh rằng tác động của tin phụ
thuộc vào cả nội dung, mức độ phân phối và bối cảnh tiếp nhận của nhà đầu tư. Quan điểm
này phù hợp với điều kiện thị trường thực tế, nhưng yêu cầu siêu dữ liệu về thời điểm,
nguồn và mức độ phổ biến của tin. Các trường dữ liệu này thường không đầy đủ trong kho
ngữ liệu được tổng hợp từ nhiều website.

Siala và cộng sự tập trung đo lường tác động thực tế của cảm xúc do LLM tạo ra
\cite{Siala2026_260200086v3}. Cách tiếp cận này có ý nghĩa đối với đề tài vì văn bản
hoặc nhãn do LLM tạo ra có thể có vẻ thuyết phục nhưng chưa chắc cải thiện kết quả dự
báo ngoài mẫu. Vì vậy, đề tài không giả định trước rằng PhoBERT hoặc bất kỳ mô hình nền
nào luôn cải thiện kết quả. Nếu thí nghiệm cho thấy cảm xúc không mang lại cải thiện,
kết quả vẫn có giá trị trong việc xác định giới hạn của tín hiệu đối với dữ liệu và
giao thức đã chọn.

\subsection{Các cơ chế hợp nhất}

Các cơ chế hợp nhất có thể được chia thành ba nhóm:

\begin{enumerate}[nosep]
  \item Hợp nhất sớm (early fusion): chuẩn hóa rồi nối hai nhóm đặc trưng trước bộ mã hóa.
    Cơ chế này đơn giản, có ít tầng kết hợp và dễ triển khai, nhưng nhóm đặc trưng có
    nhiều chiều hoặc thang đo lớn có thể lấn át nhóm còn lại. Nếu dữ liệu được nối trước
    khi căn chỉnh thời gian, lỗi về mốc thời gian cũng khó được phát hiện.
  \item Hợp nhất muộn (late fusion): mỗi nhánh tạo biểu diễn hoặc dự đoán riêng, sau đó một
    lớp kết hợp các đầu ra. Cơ chế này tách biệt quy trình xử lý văn bản và giá, tạo
    thuận lợi cho việc loại bỏ đặc trưng, nhưng có thể bỏ lỡ tương tác chi tiết giữa
    một tin cụ thể và một trạng thái giá cụ thể.
  \item Hợp nhất có cổng hoặc chú ý chéo: mô hình học cách lựa chọn thông tin
    giữa hai nhánh theo thời điểm và ngữ cảnh. Cơ chế này có khả năng biểu đạt cao hơn
    và có thể giảm ảnh hưởng của tín hiệu nhiễu, nhưng yêu cầu nhiều dữ liệu hơn và dễ
    làm tăng phương sai.
\end{enumerate}

Đề tài sử dụng phép nối sau hai nhánh LSTM làm cấu hình cơ sở vì cấu trúc này rõ ràng,
dễ kiểm tra và phù hợp với phép loại bỏ đặc trưng. Nhánh giá xử lý cửa sổ OHLCV cùng
các chỉ báo kỹ thuật, trong khi nhánh cảm xúc xử lý các thống kê xác suất được tổng hợp
theo cùng lịch ngày. Hai vector sau đó được nối trước lớp phân loại. Hợp nhất có cổng,
chú ý chéo hoặc Transformer hai luồng được dành cho các hướng phát triển tiếp theo.
Transformer bậc cao là một hướng mở rộng khi cần xử lý tương tác giữa nhiều chuỗi trong
cùng một Transformer \cite{Omranpour2024_241210540v1}.

\subsection{Đánh giá đa phương thức và nguy cơ diễn giải quá mức}

Một mô hình đa phương thức có thể đạt điểm cao do nhiều nguyên nhân. Cảm xúc có thể phản ánh lại giá đã xuất hiện trong cùng ngày; tin có thể được xếp sai thời
điểm; bản sao của bài viết có thể xuất hiện trong cả hai tập; hoặc một chỉ báo kỹ thuật
có thể đã mã hóa gần như toàn bộ tín hiệu được cho là do cảm xúc cung cấp. Vì vậy, việc
chỉ so sánh một mô hình hợp nhất với một mô hình chỉ sử dụng giá nhưng có cấu hình khác
là chưa đủ để kết luận về giá trị của văn bản.

Thiết kế của đề tài yêu cầu bốn điều kiện. Thứ nhất, việc ánh xạ tin theo mốc thời gian
phải được kiểm tra; trong đó, tin sau giờ đóng cửa, tin cuối tuần và tin ngày nghỉ được
ánh xạ sang phiên đầu tiên mà nhà đầu tư có thể quan sát. Thứ hai, hai cấu hình loại bỏ
đặc trưng phải sử dụng cùng cửa sổ, hạt giống, quy trình tiền xử lý, ngân sách huấn
luyện và tập giữ lại. Thứ ba, các chỉ số phải được tính trên dự đoán ngoài mẫu theo thứ
tự thời gian. Thứ tư, chênh lệch phải được báo cáo cùng độ bất định, chẳng hạn khoảng
tin cậy được xác định bằng lấy mẫu lặp theo khối ngày. Mô phỏng giao dịch có thể được
bổ sung sau, nhưng không thay thế việc kiểm tra phân loại ngoài mẫu
\cite{Chen2021_210708721v1,Feng2018_181009936v2}.
\section{Khoảng trống nghiên cứu và định vị đề tài}
\label{sec:tq-gap}

\subsection{Cơ sở xác định khoảng trống}

Kết quả khảo sát cho thấy các thành phần riêng lẻ đã được nghiên cứu ở nhiều mức độ khác nhau.
FinBERT cung cấp mô hình nền cho phân tích cảm xúc tài chính bằng tiếng Anh
\cite{Yang2020_200608097v2}, trong khi PhoBERT cung cấp mô hình nền đơn ngữ tiếng Việt
\cite{Nguyen2020_200300744v3}. Các quy trình FinBERT-LSTM kết hợp tin tức với chuỗi giá
đã được Halder và Gu trình bày \cite{Halder2022_221107392v1,Gu2024_240716150v1}.
Các biểu diễn ngữ cảnh và sự kiện mở rộng tín hiệu văn bản
\cite{Chen2021_210708721v1,Chen2019_191005078v1}, còn Transformer, đồ thị và
học siêu tham số mở rộng phương pháp mô hình hóa chuỗi hoặc quan hệ liên mã
\cite{Omranpour2024_241210540v1,Boyle2023_230503835v1,
Xu2021_211013716v2,Zhao2023_230609862v3}.

Các tài liệu đã được xác minh khác nhau về nhiệm vụ, quy mô dữ liệu, đặc trưng đầu vào và giao
thức đánh giá. Do đó, các kết quả công bố không được xem như những phép đo trên cùng một benchmark. Đồ án chỉ kế thừa những cơ sở phương pháp được mô tả rõ trong
các tài liệu này và tiến hành kiểm chứng riêng trên dữ liệu, nhiệm vụ của đề tài.

\subsection{Các khoảng trống nghiên cứu}

Từ các công trình đã khảo sát, đề tài xác định năm khoảng trống cần tiếp tục nghiên
cứu như sau.

Thứ nhất, các tài liệu đã xác minh cho thấy tín hiệu văn bản có thể được biểu diễn bằng
nhãn cảm xúc, biểu diễn ngữ cảnh hoặc thông tin sự kiện
\cite{Chen2019_191005078v1,Chen2021_210708721v1}. Các biểu diễn này có mức độ chi tiết
khác nhau và đặt ra những yêu cầu khác nhau đối với dữ liệu gán nhãn. Đồ án lựa chọn xác suất của
ba lớp cảm xúc ở mức sàn, đồng thời đánh giá riêng chất lượng phân loại trước khi sử dụng
đầu ra làm đầu vào cho mô hình giá. Quy trình này cho phép tách biệt lỗi của bộ phân loại
với lỗi của mô hình dự báo.

Thứ hai, các phương pháp sử dụng mô hình ngôn ngữ lớn có thể tạo ra biểu diễn giàu ngữ cảnh,
nhưng đồng thời làm tăng yêu cầu về điểm kiểm, phiên bản, chi phí suy luận và khả năng tái
lập \cite{Elahi2024_241101368v1,Liang2024_241210823v2,Siala2026_260200086v3}. Vì vậy,
đồ án ưu tiên sử dụng một mô hình mở cho phép ghi nhận cấu hình và đầu ra. Mô hình
phức tạp hơn không được xem là lựa chọn mặc định.

Thứ ba, các công trình đã xác minh sử dụng những nhiệm vụ khác nhau, bao gồm dự báo giá,
dự báo chuyển động và mô hình hóa chuỗi đa phương thức
\cite{Halder2022_221107392v1,Gu2024_240716150v1,Chen2021_210708721v1,
Omranpour2024_241210540v1}. Sự khác biệt này làm hạn chế khả năng so sánh trực tiếp.
Đề tài giới hạn nhiệm vụ ở dự báo hướng của phiên kế tiếp trên nhiều mã, đồng thời
duy trì cùng một định nghĩa nhãn và giao thức để có thể so sánh giữa các mã.

Thứ tư, các phương pháp đánh giá theo thời gian chưa được áp dụng nhất quán. Việc
sử dụng kiểm định chéo mười nếp gấp cho dữ liệu thị trường có thể trộn lẫn các giai đoạn
trước và sau nếu các nếp gấp không bảo toàn thứ tự thời gian. Trong trường hợp đó, các thống kê
chuẩn hóa, ngưỡng tạo nhãn, lựa chọn đặc trưng hoặc thông tin về phân phối tương lai
có thể vô tình được đưa vào quá trình huấn luyện. Bài toán dự báo thực tế đòi hỏi dữ liệu
được phân chia theo thời điểm, mốc cắt của tin tức được xác định rõ và mọi tham số chỉ được
ước lượng trên phần dữ liệu lịch sử mà nhà đầu tư có thể quan sát tại thời điểm đó. Quy trình
cuốn chiếu theo thời gian được lựa chọn nhằm hạn chế rò rỉ và đánh giá tính ổn định qua
các cửa sổ thời gian.

Thứ năm, đóng góp riêng của tín hiệu cảm xúc chưa được đánh giá đầy đủ trong
cùng một kiến trúc và cùng điều kiện huấn luyện. Việc so sánh một mô hình có tin tức,
sử dụng XGBoost hoặc LLM, với một mô hình chỉ sử dụng giá, áp dụng LSTM hoặc một giao thức khác,
không đủ để kết luận mức cải thiện bắt nguồn từ dữ liệu tin tức hay năng lực của mô hình.
Do đó, thí nghiệm loại bỏ đặc trưng cần được thiết kế sao cho mô hình có cảm xúc và mô
hình không có cảm xúc tương đương về cửa sổ dữ liệu, cấu trúc nhánh giá,
quy trình tiền xử lý, hạt giống, ngân sách huấn luyện và tiêu chí lựa chọn. Chênh lệch
giữa hai cấu hình là căn cứ phù hợp hơn để đo lường giá trị gia tăng của cảm xúc.

Tóm lại, trong phạm vi tập tài liệu đã khảo sát và xác minh, các công trình được sử dụng
trong báo cáo chưa cung cấp một thiết lập đánh giá đồng thời thỏa mãn năm điều kiện:
trích xuất cảm xúc từ tin tức tài chính tiếng Việt; khảo sát nhiều mã trên HOSE; dự báo
ba lớp cho phiên kế tiếp; khóa thông tin theo thời điểm quan sát và đánh giá cuốn chiếu
ngoài mẫu; đo đóng góp biên của cảm xúc bằng đối chứng ghép cặp cùng độ bất định. Vì vậy,
khoảng trống trung tâm của đề tài là thiếu bằng chứng thực nghiệm có kiểm soát cho tổ hợp
dữ liệu, nhiệm vụ và giao thức này, không phải thiếu một kiến trúc có khả năng ghép văn
bản với chuỗi giá. Cách phát biểu này giới hạn kết luận trong tập tài liệu đã khảo sát và
không hàm ý rằng chưa từng có nghiên cứu kết hợp tin tức với giá cổ phiếu.

Trên cơ sở năm khoảng trống nêu trên, đề tài tập trung xây dựng mô hình phân tích cảm xúc
cho tin tức tài chính tiếng Việt bằng PhoBERT, tổng hợp xác suất cảm xúc theo từng mã cổ
phiếu và ngày giao dịch, sau đó hợp nhất kết quả với dữ liệu thị trường trong mô hình học sâu
dành cho chuỗi thời gian. Quy trình được đánh giá bằng phương pháp xác thực cuốn chiếu theo thời gian,
kết hợp mô hình đối chứng và thí nghiệm loại bỏ đặc trưng nhằm xác định liệu tín hiệu cảm xúc có
tạo ra giá trị bổ sung ổn định hay không.

\subsection{So sánh các hướng nghiên cứu}

\begin{table}[H]
  \centering
  \caption{So sánh các nhóm nghiên cứu liên quan với phạm vi đề tài.}
  \label{tab:tq-so-sanh}
  \begin{tabular}{p{3.0cm} p{4.1cm} p{4.1cm} p{2.6cm}}
    \toprule
    Nhóm & Đặc điểm chính & Hạn chế khi áp dụng cho đề tài & Vai trò \\
    \midrule
    Cảm xúc tài chính & FinBERT và PhoBERT tạo nhãn hoặc xác suất cảm xúc cho văn bản
    \cite{Yang2020_200608097v2,Nguyen2020_200300744v3} & Khác biệt về ngôn ngữ, miền,
    phương pháp gán nhãn và thời điểm của tin tức & Mô hình nền, tầng tiền xử lý \\
    Chỉ giá LSTM & Mô hình hóa cửa sổ OHLCV và các chỉ báo kỹ thuật theo thời gian
\cite{Halder2022_221107392v1} & Không đánh giá giá trị gia tăng của văn bản & Đối chứng cơ sở học sâu \\
    Transformer dự báo biến động & Sử dụng cơ chế chú ý để mô hình hóa quan hệ phụ thuộc theo thời gian, đa phương thức hoặc
    liên mã \cite{Boyle2023_230503835v1,Omranpour2024_241210540v1} & Có độ phức tạp cao,
    chưa có bằng chứng trực tiếp trên HOSE & TFT mở rộng \\
    Hợp nhất tin tức - giá & Kết nối cảm xúc, biểu diễn, sự kiện hoặc LLM
    \cite{Chen2021_210708721v1,Elahi2024_241101368v1} & Dễ trộn lẫn khác biệt về dữ liệu,
    rò rỉ và giả định thị trường & Mẫu kiến trúc \\
    Đồ thị và học siêu tham số & Khai thác quan hệ liên mã và hiện tượng trôi dạt khái niệm
    \cite{Xu2021_211013716v2,Zhao2023_230609862v3} & Nằm ngoài phạm vi 10 mã độc lập và
    ngân sách của đồ án & Bối cảnh tương lai \\
    \bottomrule
  \end{tabular}
\noindent\textit{Nguồn: Tác giả xây dựng từ các tài liệu được trích dẫn.}
\end{table}

\subsection{Định vị và giả thuyết kiểm chứng}

Đề tài được định vị là một nghiên cứu thực nghiệm đa phương thức trên dữ liệu tiếng
Việt. Đề tài không phải cuộc thi tìm kiến trúc mới. Mức sàn gồm đối chứng lớp phổ biến, đối chứng ngẫu nhiên, các mô hình
cổ điển chỉ sử dụng giá, LSTM chỉ sử dụng giá và LSTM hai nhánh với phép nối. TFT chỉ được triển khai sau khi
mức sàn đã có số liệu, đầu ra được cấu hình cho bài toán phân loại và tập giữ lại đã được khóa. Cách
tổ chức theo tầng này giúp phân tách câu hỏi về chất lượng dữ liệu và đặc trưng khỏi câu hỏi về độ
phức tạp của mô hình. Huấn luyện đối kháng trong các benchmark chuyển động quốc tế cũng cho thấy việc xây dựng đối chứng cơ sở và kiểm tra độ bền quan trọng không kém việc thêm tầng mô hình \cite{Feng2018_181009936v2}.

Giả thuyết chính cho rằng vector xác suất cảm xúc tổng hợp có thể cung cấp thông tin gia
tăng sau khi đã kiểm soát lịch sử giá và các chỉ báo kỹ thuật. Giả thuyết phụ cho rằng phần gia
tăng, nếu tồn tại, có thể biểu hiện rõ hơn vào những ngày có nhiều tin hoặc trong các giai đoạn biến động
cao. Hai giả thuyết này được kiểm tra bằng loại bỏ đặc trưng và phân tầng dữ liệu, không được
trình bày như khẳng định rằng tin tức luôn gây ra biến động. Quá trình kiểm tra cũng cần
phân biệt trường hợp cảm xúc không có tác động với trường hợp quy trình gặp lỗi ánh
xạ thời gian hoặc nhãn không đáng tin cậy.

\subsection{Tiêu chí để tuyên bố đóng góp}

Một tuyên bố đóng góp chỉ được đưa ra khi thỏa mãn các điều kiện sau:

\begin{itemize}[nosep]
  \item Dữ liệu và phạm vi mã được mô tả rõ ràng, bao gồm các bản ghi thiếu hoặc bị loại;
  \item Mốc cắt, công thức nhãn, phép chia theo thời gian và quy trình ước lượng tham số được
    ghi nhận trong bản kê;
  \item Kết quả cảm xúc và kết quả dự báo được trình bày thành hai tầng;
  \item Mô hình có cảm xúc và mô hình chỉ sử dụng giá áp dụng cùng giao thức, hạt giống hoặc tập hạt giống;
  \item F1 vĩ mô, độ chính xác cân bằng, AUC-ROC một-chống-còn-lại và các chỉ số theo lớp được báo cáo
    cùng khoảng tin cậy khi có thể;
  \item Kết luận nêu rõ phạm vi các mã thuộc VN30 được chọn và không suy rộng sang thị
    trường khác.
\end{itemize}

Theo đó, tuyên bố đóng góp thận trọng nhất của đề tài là xây dựng một quy trình có thể tái lập để đo lường phần
đóng góp biên của cảm xúc tiếng Việt thông qua đánh giá cuốn chiếu theo thời gian có kiểm soát. Nếu
thí nghiệm loại bỏ đặc trưng không cho thấy sự cải thiện, kết quả này vẫn trả lời câu hỏi nghiên cứu và giúp
tránh đưa ra kết luận quá mức từ một mô hình phức tạp. Nếu kết quả cho thấy sự cải thiện, báo cáo cần
xác định rõ cải thiện xuất hiện ở mã nào, giai đoạn nào, mật độ tin tức nào và có ổn định
qua các hạt giống hay không.

\chapter{CƠ SỞ LÝ THUYẾT VÀ PHƯƠNG PHÁP NGHIÊN CỨU}
\label{ch:phuong-phap}

Chương này trình bày cơ sở lý thuyết, dữ liệu, kiến trúc hệ thống và giao thức đánh giá của đề tài. Quy trình được thiết kế để kết nối hai nguồn thông tin có bản chất khác nhau: tin tức tài chính tiếng Việt phản ánh các sự kiện và kỳ vọng, trong khi dữ liệu thị trường ghi nhận kết quả giao dịch theo thời gian. Do đó, các bước thu thập, làm sạch, gắn mã cổ phiếu, ánh xạ thời điểm, tạo nhãn và phân chia tập dữ liệu được mô tả trước phần huấn luyện mô hình.

Phương pháp nghiên cứu được tổ chức thành hai tầng. Tầng bắt buộc khóa mức sàn, bao gồm đối chứng cơ sở, mô hình chỉ sử dụng dữ liệu giá, mô hình hai nhánh, phép loại bỏ đặc trưng và đánh giá cuốn chiếu theo thời gian. Tầng mở rộng chỉ được xem xét sau đó, bao gồm các mô hình phức tạp hơn như Bộ biến đổi hợp nhất theo thời gian, hợp nhất có cổng hoặc biểu diễn sự kiện. Cấu trúc này cho phép lần lượt đánh giá tính đầy đủ và hợp lệ của dữ liệu, khả năng học tín hiệu của mô hình giá, và cảm xúc có bổ sung thông tin hay không. Khi một tầng trước chưa được kiểm tra, điểm số của tầng sau không được dùng để che lấp vấn đề ở tầng dữ liệu.
\section{Tổng quan kiến trúc hệ thống}
\label{sec:pp-kien-truc}

Hệ thống gồm hai nhánh xử lý riêng trước khi hợp nhất. Nhánh cảm xúc tiếp nhận tiêu
đề và nội dung tin tức tiếng Việt trong phạm vi được phép sử dụng. PhoBERT được tinh
chỉnh để sinh xác suất cho ba lớp tích cực, trung tính và tiêu cực; các tín hiệu này
sau đó được tổng hợp theo mã cổ phiếu và ngày giao dịch. Nhánh giá tiếp nhận dữ liệu
OHLCV điều chỉnh để tính lợi suất, độ biến động, khối lượng và các chỉ báo kỹ thuật
theo lịch phiên. Sau khi áp dụng quy tắc mốc cắt, hai nhánh được căn chỉnh theo khóa
gồm mã cổ phiếu và ngày giao dịch.

Mỗi dòng trong bảng đặc trưng phải kèm theo các trạng thái dữ liệu cần thiết, chẳng hạn
trạng thái có tin trong ngày, tính hợp lệ của mốc thời gian của tin và giá có đủ lịch sử hay không. Bảng đặc trưng thu được được sử dụng để dự báo lớp UP, FLAT hoặc DOWN
của phiên kế tiếp. Giá của phiên kế tiếp chỉ được đưa vào cột nhãn sau khi dữ liệu đầu
vào đã được khóa và không được sử dụng ngược để tạo đặc trưng cho ngày quan sát.

\subsection{Nguyên tắc thiết kế}

Kiến trúc được xây dựng dựa trên năm nguyên tắc. Thứ nhất, mỗi bản ghi phải có mốc thời
gian, nguồn gốc và trạng thái ánh xạ để kiểm tra khả năng sử dụng tại thời điểm dự báo.
Thứ hai, tín hiệu cảm xúc được lưu dưới dạng xác suất thay vì chỉ sử dụng một nhãn
cứng. Cách biểu diễn này duy trì thông tin về mức độ không chắc chắn và cho phép tính
các đại lượng tổng hợp khác nhau. Thứ ba, mô hình sử dụng cảm xúc phải được so sánh với
mô hình chỉ sử dụng giá trong cùng điều kiện; nếu không, phần chênh lệch không thể được
quy cho thông tin văn bản. Thứ tư, mọi phép biến đổi có tham số cần học phải được ước
lượng trong cửa sổ huấn luyện tương ứng. Thứ năm, mô hình mở rộng chỉ được triển khai
sau khi mức sàn đã vượt qua các kiểm tra về dữ liệu, rò rỉ và độ ổn định. Năm nguyên
tắc này đồng thời là tiêu chí xác định một kết quả được đưa vào phân tích chính hay chỉ
được trình bày trong phần bổ sung.

\subsection{Cơ sở lý thuyết của Transformer và cơ chế chú ý}

Transformer biểu diễn chuỗi token bằng các vector và sử dụng cơ chế chú ý để xác định
mức độ liên quan giữa các vị trí \cite{Vaswani2017_attention}. Với ma trận truy vấn
$Q$, khóa $K$ và giá trị $V$, phép chú ý tích vô hướng có điều chỉnh được xác định như
sau:

\begin{equation}
  \operatorname{Attention}(Q,K,V)=
  \operatorname{softmax}\left(\frac{QK^{\mathsf{T}}}{\sqrt{d_k}}\right)V,
  \label{eq:scaled-attention}
\end{equation}

trong đó $d_k$ là số chiều của khóa. Phép chia cho $\sqrt{d_k}$ giúp duy trì độ lớn của
tích vô hướng trong khoảng ổn định trước khi áp dụng hàm softmax. Cơ chế chú ý nhiều
đầu thực hiện đồng thời nhiều phép chiếu và ghép các kết quả, qua đó cho phép mô hình
học những dạng quan hệ khác nhau giữa các token. Trong mô hình ngôn ngữ, mỗi token
được biểu diễn theo ngữ cảnh của toàn bộ câu. Điều này đặc biệt hữu ích với tin tài chính vì
ý nghĩa của một con số hoặc động từ thường phụ thuộc vào doanh nghiệp, chủ thể và câu đứng trước hoặc sau nó. FinBERT và PhoBERT là hai ví dụ về việc sử dụng mô hình
nền thuộc họ BERT cho miền tài chính hoặc tiếng Việt
\cite{Devlin2019_bert,Yang2020_200608097v2,
Nguyen2020_200300744v3}.

Một bộ mã hóa Transformer thường bao gồm cơ chế chú ý nhiều đầu, mạng truyền thẳng theo
từng vị trí, kết nối tắt và phép chuẩn hóa. Do cơ chế chú ý không tự chứa thông tin về
thứ tự, mô hình cần sử dụng biểu diễn vị trí hoặc một phương thức mã hóa tương đương.
Các khối này cho phép mô hình học quan hệ giữa các vị trí cách xa nhau, đồng thời duy
trì đường truyền gradient ổn định. PhoBERT cung cấp sẵn bộ tách từ và cách biểu diễn
theo thiết kế của mô hình nền. Khi sử dụng TFT, thứ tự ngày được duy trì trong
\texttt{TimeSeriesDataSet}; quy tắc khai báo biến theo thời điểm có thể quan sát được
được trình bày trong tiểu mục về mô hình này.

\subsection{Cơ sở lý thuyết của PhoBERT}

PhoBERT là mô hình ngôn ngữ tiền huấn luyện dành cho tiếng Việt. Văn bản được tách từ
theo quy ước phù hợp, sau đó được xử lý bằng bộ tách từ BPE để tạo chuỗi đơn vị từ con.
Khi tinh chỉnh cho bài toán cảm xúc, biểu diễn từ bộ mã hóa được đưa vào tầng phân
loại. Với ba lớp, điểm logit $\mathbf{z}$ được chuyển thành xác suất như sau:

\begin{equation}
  p_c=\frac{\exp(z_c)}{\sum_{k=1}^{3}\exp(z_k)},
  \qquad c\in\{\mathrm{NEG},\mathrm{NEU},\mathrm{POS}\}.
  \label{eq:sentiment-softmax}
\end{equation}

Việc lựa chọn PhoBERT không thay thế yêu cầu thích nghi mô hình với miền tài chính.
Tên công ty, mã cổ phiếu, thuật ngữ chuyên ngành và cách biểu diễn phần trăm có thể
khác với ngữ liệu tiền huấn luyện. Vì vậy, tập tinh chỉnh phải quy định rõ cách tách
từ, giới hạn độ dài, phương pháp xử lý thực thể và quy tắc loại trùng. Phần dữ liệu
được sử dụng để sinh đặc trưng dự báo cũng phải được xử lý bằng điểm kiểm đã được khóa
theo quy tắc thời gian. Mô hình được kiểm tra trên tập dữ liệu trong miền đã khóa thay
vì chỉ dựa trên độ chính xác của tập hạt giống. Kết quả phân loại cảm xúc cần được báo
cáo riêng với kết quả dự báo để tránh đồng nhất chất lượng của hai nhiệm vụ.

\subsection{Cơ sở lý thuyết của LSTM}

LSTM xử lý chuỗi theo thứ tự thời gian và duy trì trạng thái ẩn $\mathbf{h}_t$ cùng
trạng thái ô nhớ $\mathbf{c}_t$ \cite{Hochreiter1997_lstm}. Với đầu vào
$\mathbf{x}_t$, một bước cập nhật được mô tả như sau:

\begin{align}
  \mathbf{i}_t &= \sigma(W_i\mathbf{x}_t+U_i\mathbf{h}_{t-1}+\mathbf{b}_i),\\
  \mathbf{f}_t &= \sigma(W_f\mathbf{x}_t+U_f\mathbf{h}_{t-1}+\mathbf{b}_f),\\
  \mathbf{o}_t &= \sigma(W_o\mathbf{x}_t+U_o\mathbf{h}_{t-1}+\mathbf{b}_o),\\
  \widetilde{\mathbf{c}}_t &= \tanh(W_c\mathbf{x}_t+U_c\mathbf{h}_{t-1}+\mathbf{b}_c),\\
  \mathbf{c}_t &= \mathbf{f}_t\odot\mathbf{c}_{t-1}+
                 \mathbf{i}_t\odot\widetilde{\mathbf{c}}_t,\\
  \mathbf{h}_t &= \mathbf{o}_t\odot\tanh(\mathbf{c}_t).
  \label{eq:lstm-update}
\end{align}

Trong đó, $\mathbf{i}_t$, $\mathbf{f}_t$ và $\mathbf{o}_t$ lần lượt là cổng vào, cổng
quên và cổng ra. Cổng quên xác định phần thông tin cũ được giữ lại, cổng vào điều chỉnh
lượng thông tin mới được đưa vào ô nhớ, còn cổng ra tạo trạng thái ẩn cho bước tiếp
theo. Cơ chế này cho phép LSTM duy trì các quan hệ phụ thuộc kéo dài qua nhiều phiên mà
không cần đưa toàn bộ chuỗi vào một phép tổng hợp đơn giản.

LSTM phù hợp với mức sàn vì có thể mô hình hóa cửa sổ giá hoặc cửa sổ cảm xúc mà không
yêu cầu đồ thị liên mã. Trong hệ thống, hai nhánh LSTM sử dụng cùng số ngày quan sát,
cùng quy tắc đệm chuỗi hoặc loại mẫu và cùng điều kiện dừng để phép so sánh có ý nghĩa. Việc bổ sung nhánh cảm xúc làm tăng số tham số; do đó, số chiều trạng
thái và thời gian huấn luyện cũng cần được ghi nhận khi phân tích chênh lệch.

\subsection{Cơ sở lý thuyết của Bộ biến đổi hợp nhất theo thời gian}

TFT được đưa vào như mô hình mở rộng cho chuỗi nhiều biến
\cite{Lim2021_tft}. Mô hình kết hợp bộ mã hóa chuỗi với cơ chế chọn biến, các khối có
cổng và cơ chế chú ý. Cấu trúc này cho phép mức độ quan trọng của một biến thay đổi
theo thời điểm và trạng thái dự báo. Trong bài toán này, động lượng, độ biến động, tỷ lệ
cảm xúc và cờ có tin là các nhóm biến cần được phân biệt.

Nếu được triển khai, TFT phải được cấu hình cho bài toán phân loại. Đầu ra cuối cùng là
điểm logit kích thước 3 và hàm mất mát là entropy chéo. Cấu hình hồi quy phân vị có sẵn
trong thư viện không được sử dụng thay cho mục tiêu này. Các biến cũng phải được khai
báo rõ theo trạng thái: biết trước, biến tĩnh hoặc chỉ quan sát được tại thời điểm dự
báo. Tin tức sau mốc cắt không được đưa vào nhóm biến đã biết tại ngày trước đó.

Tầm quan trọng của biến và thành phần chú ý chỉ được lưu như bằng chứng mô tả cách mô
hình sử dụng các biến, không được diễn giải là ảnh hưởng nhân quả. Các nghiên cứu sử
dụng Transformer cho biến động và chuỗi đa phương thức cho thấy tiềm năng của hướng
tiếp cận này. Tuy nhiên, độ phức tạp, số lượng tham số, khả năng quá khớp và sự khác
biệt giữa các thị trường vẫn cần được kiểm tra riêng
\cite{Boyle2023_230503835v1,Omranpour2024_241210540v1}.

\subsection{Luồng xử lý tổng thể}

Luồng xử lý tổng thể của hệ thống gồm các bước sau:

\begin{enumerate}[nosep]
  \item Khóa phạm vi mã, khoảng thời gian, mốc cắt, công thức nhãn và hạt giống;
  \item Thu thập tin tức, OHLCV và lịch giao dịch, đồng thời lưu nguồn gốc và khả năng truy xuất;
  \item Làm sạch, loại trùng, kiểm tra mốc thời gian và ánh xạ bài viết theo mã;
  \item Tạo nhánh cảm xúc và nhánh giá bằng các phép biến đổi độc lập;
  \item Căn chỉnh hai nhánh theo mã và ngày giao dịch, sau đó tạo nhãn ở thời điểm hợp lệ;
  \item Chạy thang mô hình từ đối chứng cơ sở đến LSTM hai nhánh, chỉ chạy TFT khi đủ điều kiện;
  \item Đánh giá bằng cuốn chiếu theo thời gian, loại bỏ đặc trưng và lấy mẫu lặp, rồi lưu bản kê kết quả.
\end{enumerate}

\ifShowDiagram
Hình~\ref{fig:kien-truc} minh họa các bước chính của quy trình xử lý. Các khối TFT và
chú ý chéo có cổng được xác định là phần mở rộng, trong khi mô hình LSTM hai nhánh là
cấu hình cơ sở. Sự phân biệt này làm rõ các thành phần bắt buộc và các thành phần chỉ
được triển khai khi điều kiện dữ liệu và thời gian cho phép.
\else
Các bước chính của quy trình xử lý được trình bày trong phần này; sơ đồ minh họa được
lược bỏ trong biến thể báo cáo không kèm hình.
\fi
\ifShowDiagram
\begin{figure}[H]
  \centering
  \begin{tikzpicture}[
      node distance=6mm and 10mm,
      font=\small,
      stage/.style={
        rectangle, rounded corners=2pt, draw=black!55, thick,
        text width=10.5cm, minimum height=9mm, align=center,
        fill=gray!6, inner sep=4pt},
      model/.style={
        rectangle, rounded corners=2pt, draw=blue!65, very thick,
        text width=10.5cm, minimum height=9mm, align=center,
        fill=blue!8, inner sep=4pt},
      floor/.style={
        rectangle, rounded corners=2pt, draw=green!55!black, very thick,
        text width=10.5cm, minimum height=9mm, align=center,
        fill=green!10, inner sep=4pt},
      ext/.style={
        rectangle, rounded corners=2pt, draw=violet!60, very thick, dashed,
        text width=10.5cm, minimum height=9mm, align=center,
        fill=violet!6, inner sep=4pt},
      arr/.style={-{Stealth[length=2.5mm]}, thick, black!65},
    ]
    \node[stage] (g1) {1. Thu thập dữ liệu\\ Giá OHLCV của 10 mã VN30 và tin tức Vietstock trong giai đoạn 2020--2025};
    \node[model, below=of g1] (g2) {2. Phân tích cảm xúc\\ PhoBERT được tinh chỉnh để sinh xác suất của 3 lớp: tích cực, trung tính và tiêu cực};
    \node[stage, below=of g2] (g3) {3. Đặc trưng giá\\ Lợi suất, lợi suất log, MA, khối lượng và biến động};
    \node[stage, below=of g3] (g4) {4. Căn chỉnh và gán nhãn\\ Gộp cảm xúc và giá theo mã và ngày, mốc cắt 15:00, nhãn UP, FLAT và DOWN};
    \node[floor, below=of g4] (g5) {5. Mô hình dự báo\\ Đối chứng cơ sở, sau đó là LSTM một nhánh và LSTM hai nhánh ở mức sàn tốt nghiệp};
    \node[ext, below=of g5] (g6) {Mở rộng có điều kiện: đặc trưng kỹ thuật và TFT};
    \node[stage, below=of g6] (g7) {6. Đánh giá\\ Cuốn chiếu theo thời gian + loại bỏ đặc trưng có/không cảm xúc, F1 vĩ mô, khoảng tin cậy lấy mẫu lặp};

    \draw[arr] (g1) -- (g2);
    \draw[arr] (g2) -- (g3);
    \draw[arr] (g3) -- (g4);
    \draw[arr] (g4) -- (g5);
    \draw[arr] (g5) -- (g6);
    \draw[arr] (g6) -- (g7);
  \end{tikzpicture}
  \caption{Kiến trúc tổng thể của quy trình dự báo xu hướng giá. Khối xanh dương thể
  hiện mô hình cảm xúc PhoBERT, khối xanh lá là mức sàn tốt nghiệp, khối nét đứt là
  phần mở rộng.}
  \label{fig:kien-truc}
\noindent\textit{Nguồn: Tác giả xây dựng.}
\end{figure}
\fi
\section{Thu thập và tiền xử lý dữ liệu}
\label{sec:pp-du-lieu}

\subsection{Phạm vi dữ liệu và nguồn thu thập}

Dữ liệu gồm hai bảng chính. Bảng tin tức lưu mã bản ghi, URL, nguồn, tiêu đề, nội
dung nếu có, thời điểm đăng hoặc cập nhật, thời điểm thu thập và mã cổ phiếu liên
quan. Bảng giá lưu mã cổ phiếu, ngày giao dịch, dữ liệu OHLCV điều chỉnh và cờ chất
lượng. Theo thiết kế hiện hành, giai đoạn thực nghiệm được cố định từ năm 2020 đến hết
năm 2025, với danh sách 10 mã HOSE thuộc VN30 được ghi trong bản kê. Nếu phạm vi thực
tế thay đổi do mức độ bao phủ của dữ liệu, bản kê phải ghi rõ mốc bắt đầu, mốc kết thúc
và lý do thay đổi, thay vì chỉ trình bày phạm vi trong phần mô tả.

Nguồn tin được ưu tiên hiện nay là Vietstock theo từng mã cổ phiếu, do cách tổ chức
này giúp giảm sự mơ hồ khi ánh xạ bài viết với cổ phiếu. CafeF, VnEconomy hoặc
Stockbiz chỉ được bổ sung sau khi đã kiểm tra độ phủ, nội dung, mốc thời gian và điều
kiện sử dụng. Theo đề cương, đây là các nguồn dự kiến và có thể được điều chỉnh khi
khả năng truy cập hoặc cấu trúc trang thay đổi. Dữ liệu giá được thu thập thông qua
thư viện \texttt{vnstock}, sử dụng nguồn VCI theo cấu hình của đề tài. Nguồn khác như
\texttt{vnquant} chỉ được sử dụng để đối chiếu sau khi đã so sánh ngữ nghĩa của giá
điều chỉnh, độ bao phủ và lược đồ; dữ liệu từ nguồn này không được tự động trộn vào
tập dữ liệu chính. Việc phân tách nguồn chính và nguồn đối chiếu hỗ trợ phát hiện sai
lệch nhưng không làm tăng số lượng quan sát bằng cách ghép các bản ghi chưa thống
nhất về quy ước.

\subsection{Kiểm tra chất lượng, nguồn gốc và khả năng truy xuất}

Mỗi bản ghi tin tức được gắn mã băm của văn bản đã chuẩn hóa, cùng với nguồn và URL.
Các bản ghi trùng tiêu đề, trùng nội dung hoặc có mốc thời gian gần nhau được đánh dấu
để kiểm tra trước khi loại bỏ. Quá trình loại trùng được thực hiện trước khi chia tập
dữ liệu cảm xúc và trước khi tổng hợp đặc trưng, vì việc một bản sao xuất hiện trong
hai phần chia có thể làm tăng kết quả một cách giả tạo. Nếu một bài viết được gắn với
nhiều mã hợp lệ, hệ thống giữ nguyên quan hệ nhiều mã và ghi rõ phương thức mở rộng
bản ghi. Một bài viết có thể được tính cho nhiều mã trong bảng liên kết, nhưng không
được xem là nhiều bài độc lập khi đánh giá độ không chắc chắn.

Nếu không xác định được mã cổ phiếu hoặc thời điểm, bản ghi không được đưa vào bảng
đặc trưng chính. Các bản ghi bị loại vẫn được thống kê theo nguồn và lý do loại để
người đọc có thể đánh giá độ bao phủ. Khi thời điểm đăng và thời điểm cập nhật khác
nhau, cả hai trường phải được giữ lại nếu có; quy tắc lựa chọn thời điểm sử dụng phải
được khóa trước khi tạo mẫu.

Các thống kê phải bao gồm số lượng bài viết theo nguồn, mã cổ phiếu và năm; tỷ lệ
thiếu nội dung; tỷ lệ thiếu mốc thời gian; số bản ghi trùng; độ phủ của các ngày có
tin; và số phiên giá hợp lệ. Báo cáo cũng cần ghi số bài gắn với một mã và nhiều mã,
số ngày không có tin, số bản ghi được chuyển sang phiên kế tiếp và số bản ghi không
thể xác minh. Bản kê đồng thời lưu thời điểm thu thập, phiên bản mã nguồn, mã băm của
tệp dữ liệu và quy tắc làm sạch. Bước này nhằm bảo đảm khả năng truy nguyên của kho
ngữ liệu, thay vì chỉ tạo một tập dữ liệu có quy mô lớn nhưng khó kiểm tra.

\subsection{Dữ liệu giá}

Dữ liệu OHLCV được sử dụng theo giá điều chỉnh để tính lợi suất và hạn chế các điểm
đứt gãy do chia tách, cổ tức hoặc sự kiện doanh nghiệp. Với giá đóng cửa điều chỉnh
$C^{\mathrm{adj}}_{i,t}$ của mã $i$ tại ngày $t$, lợi suất mục tiêu được định nghĩa
như sau:

\begin{equation}
  r_{i,t+1}=\frac{C^{\mathrm{adj}}_{i,t+1}}
  {C^{\mathrm{adj}}_{i,t}}-1.
  \label{eq:next-return}
\end{equation}

Các kiểm tra cơ bản bao gồm thứ tự thời gian, ngày trùng, giá không dương, khối lượng
không hợp lệ và phiên bị thiếu. Dữ liệu từ nguồn chính được đối chiếu với nguồn thay
thế trong phạm vi cần thiết nhằm phát hiện ngày thiếu, bản ghi bất thường hoặc sai lệch
do thay đổi đơn vị. Các sự kiện chia tách, trả cổ tức bằng cổ phiếu hoặc điều chỉnh giá
được xem xét khi xây dựng chuỗi, vì việc sử dụng dữ liệu không nhất quán có thể tạo ra
lợi suất bất thường không phản ánh giao dịch thực tế.

Giá trị tại ngày $t+1$ chỉ được sử dụng để tạo nhãn, không được dùng để tạo đặc trưng
tại ngày $t$. Các ngày đầu chuỗi chưa có đủ dữ liệu để tính trung bình động hoặc độ
biến động được loại khỏi mẫu hoặc được xử lý theo quy tắc ghi trong bản kê. Mọi khác
biệt giữa giá đóng cửa gốc và giá điều chỉnh phải được giải thích trong báo cáo dữ
liệu; một chuỗi không được tự động xem là sai chỉ vì hai nguồn áp dụng phương pháp
điều chỉnh khác nhau.

\subsection{Dữ liệu tin tức và ánh xạ theo phiên}

Mỗi bản tin giữ lại cả tiêu đề và nội dung để phục vụ kiểm tra, nhưng chính sách lựa
chọn đầu vào phải được khóa trước khi huấn luyện. Cấu hình cơ sở ưu tiên sử dụng tiêu
đề vì phù hợp với tập hạt giống công khai và làm giảm khác biệt về độ dài. Nội dung
được lưu để kiểm tra và chỉ được sử dụng trong phân tích độ nhạy khi đã được gán nhãn,
tách theo thời gian và báo cáo riêng. Khi một bài viết chỉ có tiêu đề, hệ thống không
tự suy diễn phần nội dung còn thiếu; trường dữ liệu thiếu phải được lưu lại để phân
tích độ bao phủ.

Tất cả mốc thời gian được chuyển về múi giờ Việt Nam. Đối với mỗi ngày giao dịch $t$,
chỉ những tin có mốc thời gian không muộn hơn 15 giờ trong ngày mới được xem là có thể
quan sát tại thời điểm ra quyết định. Tin sau 15 giờ được chuyển sang phiên giao dịch
kế tiếp theo quy tắc bảo thủ và không được dùng để tạo đặc trưng cho phiên đã đóng.
Tin xuất hiện vào cuối tuần, ngày nghỉ hoặc ngày không giao dịch cũng được gán cho
phiên hợp lệ kế tiếp. Quy tắc này có thể làm trễ một phần tín hiệu, nhưng giúp giảm
nguy cơ đánh giá lạc quan do sử dụng tin xuất hiện sau mốc cắt.

Quá trình ánh xạ được thực hiện sau khi chuẩn hóa múi giờ và sắp xếp lịch giao dịch,
không chỉ dựa trên ngày được ghi trong URL. Đối với tin có cả thời điểm đăng và thời
điểm cập nhật, trường thời gian được sử dụng cho dự báo phải được ghi trong bản kê.
Các bản ghi có mốc thời gian bị thiếu, mâu thuẫn hoặc chỉ có ngày mà không có giờ
không được gán một giờ mặc định rồi đưa trực tiếp vào tập dữ liệu chính; các bản ghi
này phải được đánh dấu để loại bỏ hoặc phân tích riêng.

\subsection{Tiền xử lý văn bản tiếng Việt}

Quy trình tiền xử lý văn bản gồm các bước sau:

\begin{enumerate}[nosep]
  \item Chuẩn hóa Unicode, khoảng trắng, ký hiệu phần trăm và cách biểu diễn số, nhưng
    giữ lại các ký hiệu có thể mang ý nghĩa tài chính;
  \item Tách từ bằng công cụ phù hợp với PhoBERT, giữ nguyên dấu tiếng Việt và các đơn
    vị từ tài chính có ý nghĩa như mã cổ phiếu, tên công ty và đơn vị tiền tệ;
  \item Nhận diện và che một số thực thể như tổ chức, địa danh hoặc mã cổ phiếu khi
    mục tiêu là học ngữ cảnh, đồng thời giữ bản gốc trong các trường nguồn gốc và
    truy xuất;
  \item Loại bỏ HTML, quảng cáo, phần điều hướng, bài rỗng và nội dung không thuộc văn
    bản tin theo quy tắc đã xác định, không tự tạo câu thay thế;
  \item Giới hạn độ dài theo điểm kiểm, sử dụng bộ tách từ BPE của PhoBERT và ghi lại
    chiều dài cùng số đơn vị từ bị cắt trong từng bản ghi.
\end{enumerate}

Số liệu, dấu phần trăm và ký hiệu tăng giảm không bị loại bỏ một cách máy móc vì có thể
chứa tín hiệu tài chính. Nếu áp dụng phép loại dấu câu hoặc che thực thể, quy tắc này
phải được duy trì nhất quán giữa các phần chia. Mọi phép biến đổi sử dụng thông tin từ
toàn bộ kho ngữ liệu, chẳng hạn vốn từ mới hoặc thống kê lựa chọn token, phải được loại
khỏi thực nghiệm chính hoặc chỉ được ước lượng trên phần huấn luyện.

\subsection{Tạo tập cảm xúc và gán nhãn}

Tập hạt giống gồm các tiêu đề tài chính đã được gán nhãn theo ba lớp và được sử dụng
để kiểm tra nhanh quy trình. Tập dữ liệu trong miền được mở rộng theo một hướng dẫn
gán nhãn cố định. Người gán nhãn đọc phần nội dung được phép sử dụng, xác định chủ thể
chính, sự kiện và hướng tác động kỳ vọng đối với giá trong ngắn hạn; nội dung đầy đủ
chỉ được đọc khi thuộc thời điểm và phạm vi được phép sử dụng. Các quy tắc cơ bản gồm:

\begin{itemize}[nosep]
  \item Tích cực nếu thông tin hàm ý triển vọng hoặc tác động có lợi rõ ràng;
  \item Tiêu cực nếu thông tin hàm ý rủi ro, suy giảm hoặc tác động bất lợi rõ ràng;
  \item Trung tính nếu thông tin chỉ mang tính thông báo, mô tả hoặc chưa có đủ bằng
    chứng để xác định chiều tác động.
\end{itemize}

Tin chứa hai chiều tác động trái ngược được gán nhãn theo tác động nổi trội nếu hướng
dẫn cho phép, hoặc được đánh dấu để rà soát. Tin không xác định được chủ thể hoặc hướng
tác động rõ ràng được giữ lại với nhãn trung tính khi có đủ căn cứ; nếu không, mẫu được
loại khỏi tập huấn luyện và lý do loại được lưu lại.

Mỗi mẫu được gán một trong ba lớp NEGATIVE, NEUTRAL hoặc POSITIVE theo hướng dẫn đã
khóa. Báo cáo gán nhãn phải ghi số lượng mẫu, tiêu chí, phiên bản hướng dẫn, số mẫu đã
được rà soát và phân bố lớp. Mô hình ngôn ngữ lớn chỉ được sử dụng để tạo nhãn sơ bộ;
một người duy nhất rà soát toàn bộ nhãn trước khi huấn luyện chính thức. Do không có
người gán nhãn thứ hai, không thể tính Cohen's hoặc Fleiss' kappa. Các mẫu dùng để
đánh giá dự báo phải được tách theo thời gian khỏi dữ liệu dùng để tinh chỉnh cảm xúc
trong phần chia tương ứng.
\section{Nhánh cảm xúc: tinh chỉnh PhoBERT}
\label{sec:pp-phobert}

\subsection{Thiết lập bài toán phân loại}

Mỗi bài viết hợp lệ được xem là một mẫu có nhãn $y_j$ thuộc một trong ba lớp NEGATIVE, NEUTRAL
hoặc POSITIVE. Nhãn biểu thị tác động tiềm năng của nội dung đối với doanh nghiệp hoặc cổ phiếu
liên quan, thay vì chỉ phản ánh cảm xúc chung của người viết. Mô hình sử dụng điểm kiểm
\texttt{vinai/phobert-base} làm bộ mã hóa, sau đó là tầng bỏ ngẫu nhiên và tầng phân loại tuyến
tính. Điểm logit được chuyển thành xác suất theo Phương trình~\ref{eq:sentiment-softmax}.
Tầng phân loại tuyến tính được chọn làm cấu hình chính do có ít tham số bổ sung và thuận tiện
cho việc tái lập. Các cấu hình kết hợp PhoBERT với CNN hoặc BiLSTM chỉ được bổ sung khi dữ
liệu và tiến độ cho phép, đồng thời chỉ được so sánh sau khi cấu hình cơ sở đã được khóa.

PhoBERT được lựa chọn phù hợp với ngôn ngữ của dữ liệu, trong khi FinBERT được sử dụng để
đối chiếu về mặt khái niệm trong miền tiếng Anh
\cite{Nguyen2020_200300744v3,Yang2020_200608097v2}. Cấu hình mức sàn không dịch toàn bộ
tin sang tiếng Anh vì dịch máy có thể làm thay đổi sắc thái, tên riêng và số liệu. Nếu thực
hiện thí nghiệm dịch, thí nghiệm này phải được xem là một đối chứng cơ sở riêng, với phiên bản công cụ và cấu hình được ghi lại.

\subsection{Chia dữ liệu và huấn luyện}

Quy trình xác thực chéo của PhoBERT sử dụng \texttt{StratifiedKFold} ngẫu nhiên với năm phần
chia và duy trì tỷ lệ của ba lớp trong từng tập kiểm thử ngoài. Phép đo này chỉ phản ánh hiệu
năng phân loại trong miền dữ liệu, không phải là đánh giá theo đúng thời điểm của nhánh dự
báo. Một phần của tập huấn luyện ngoài được lấy ngẫu nhiên theo tỷ lệ lớp để lựa chọn điểm
kiểm; tập kiểm thử ngoài không tham gia vào quá trình lựa chọn này. Các bài có đầu vào trùng
khớp hoàn toàn được loại bỏ trước khi chia tập nhằm hạn chế khả năng mô hình ghi nhớ câu chữ
và tạo ra kết quả đánh giá lạc quan. Các bản gần trùng hoặc bài viết cùng nguồn chưa được
phát hiện tự động, vì vậy đây vẫn là một giới hạn của giao thức.

Khi mô hình cảm xúc được sử dụng trong quy trình cuốn chiếu theo thời gian, có hai cách triển
khai hợp lệ: tinh chỉnh lại trong từng phần chia bằng nhãn của dữ liệu quá khứ, hoặc sử dụng
một điểm kiểm đóng băng được huấn luyện hoàn toàn trước giai đoạn kiểm thử. Phương án được
chọn phải được ghi trong bản kê và áp dụng nhất quán cho tất cả mô hình so sánh. Không được
tinh chỉnh mô hình bằng nhãn hoặc văn bản của tương lai rồi xem đầu ra là ngoài mẫu. Nếu dữ
liệu gán nhãn không đủ để tinh chỉnh lại trong từng phần chia, giới hạn này phải được nêu rõ
cùng với phạm vi diễn giải.

Hàm mất mát cơ bản là entropy chéo:

\begin{equation}
  \mathcal{L}_{\mathrm{CE}}=-\frac{1}{N}\sum_{j=1}^{N}
  \log p_{j,y_j}.
  \label{eq:sentiment-loss}
\end{equation}

Nếu các lớp mất cân bằng, trọng số $w_c$ được tính từ phần tập huấn luyện và được sử dụng
như sau:

\begin{equation}
  \mathcal{L}_{\mathrm{WCE}}=-\frac{1}{N}\sum_{j=1}^{N}
  w_{y_j}\log p_{j,y_j}.
  \label{eq:weighted-sentiment-loss}
\end{equation}

Trọng số, phương pháp lấy mẫu lại, hạt giống, kích thước lô, số vòng huấn luyện, bộ tách từ
và điểm kiểm đều được ghi trong bản kê. Điểm kiểm không được lựa chọn dựa trên kết quả của
tập kiểm thử. Bên cạnh các điểm số tổng hợp, cần lưu ma trận nhầm lẫn và các mẫu dự đoán sai
đại diện để phân tích các trường hợp chứa phủ định, số liệu, tên riêng, tin có nhiều sự kiện
và bài trung tính.

\subsection{Đầu ra xác suất và tổng hợp theo mã, ngày}

Gọi $\mathcal{N}_{i,t}$ là tập hợp các tin hợp lệ của mã $i$ được ánh xạ vào ngày $t$, với
$n_{i,t}=|\mathcal{N}_{i,t}|$. Xác suất trung bình của lớp $c$ được tính như sau:

\begin{equation}
  \bar p_{i,t,c}=\frac{1}{n_{i,t}}
  \sum_{j\in\mathcal{N}_{i,t}}p_{i,j,t,c},
  \qquad n_{i,t}>0.
  \label{eq:daily-sentiment}
\end{equation}

Vector đặc trưng cảm xúc bao gồm $\bar p_{\mathrm{NEG}}$, $\bar p_{\mathrm{NEU}}$,
$\bar p_{\mathrm{POS}}$, chênh lệch $\bar p_{\mathrm{POS}}-\bar p_{\mathrm{NEG}}$,
số tin, độ lệch chuẩn của xác suất, tỷ lệ tin tích cực hoặc tiêu cực và biến nhị phân
\texttt{has\_news}. Có thể bổ sung các cửa sổ tổng hợp ngắn hạn, chẳng hạn một ngày, ba ngày
hoặc năm ngày giao dịch gần nhất, để kiểm tra khả năng tín hiệu tiếp tục kéo dài sau ngày
công bố.

Trong cấu hình mức sàn, ngày không có tin được gán vector xác suất theo tiên nghiệm trung
tính; số tin và các tỷ lệ được đặt bằng không, đồng thời \texttt{has\_news}=0. Ngày chỉ có
tin trung tính vẫn có số tin dương và \texttt{has\_news}=1, do đó hai trường hợp này được
phân biệt. Ngoài cấu hình chính, có thể so sánh phương án đặt vector bằng không hoặc giữ giá
trị gần nhất trong một giới hạn thời gian trên tập xác thực. Cách biểu diễn ngày không có tin
không được lựa chọn dựa trên điểm của đoạn kiểm thử.

Các đặc trưng được tổng hợp sau mốc cắt. Số tin hoặc xác suất của bài đăng sau mốc cắt không
được sử dụng để mô tả ngày đã kết thúc. Nếu một bài được gắn với nhiều mã, phương pháp nhân
bản đặc trưng phải được ghi rõ để tránh xem bài đó là nhiều quan sát độc lập khi phân tích
độ không chắc chắn. Mức độ phủ tin, số lượng bài và tỷ lệ bản ghi bị loại cũng được lưu lại,
qua đó giúp người đọc phân biệt giữa ``không có tin'' và ``có tin nhưng mô hình không chắc
chắn'' khi diễn giải một ngày có ít tín hiệu.
\section{Nhánh giá và ghép đặc trưng}
\label{sec:pp-ghep}

\subsection{Đặc trưng kỹ thuật}

Từ dữ liệu OHLCV đã điều chỉnh, hệ thống xây dựng các nhóm đặc trưng sau:

\begin{itemize}[nosep]
  \item Lợi suất ngày và lợi suất log;
  \item Tỷ lệ giữa giá đóng cửa và trung bình động năm phiên;
  \item Mức thay đổi khối lượng trong một phiên;
  \item Độ lệch chuẩn của lợi suất trong cửa sổ năm phiên để biểu diễn độ biến động;
  \item Số lượng tin, xác suất cảm xúc tổng hợp, tỷ lệ lớp và biến chỉ báo có tin.
\end{itemize}

Các đặc trưng được tính từ dữ liệu có sẵn đến ngày quan sát $t$. Nhằm giảm mức độ phụ thuộc vào giá tuyệt đối, các đại lượng có đơn vị giá được chuyển thành tỷ lệ hoặc sai khác tương đối. Mỗi đặc trưng chỉ sử dụng thông tin trong quá khứ và không sử dụng giá của phiên kế tiếp. Các hàng ở đầu chuỗi chưa có đủ dữ liệu lịch sử được giữ ở dạng khuyết và loại bỏ khi tạo cửa sổ. Bộ chuẩn hóa chỉ được ước lượng một lần trên phần huấn luyện, sau đó được áp dụng nguyên trạng cho các phần xác thực và kiểm thử.

\subsection{Căn chỉnh hai nguồn dữ liệu}

Mỗi dòng trong bảng đặc trưng có khóa $(i,t)$, bao gồm vector giá $\mathbf{x}^{p}_{i,t}$, vector cảm xúc $\mathbf{x}^{s}_{i,t}$, cờ chất lượng, thông tin mã và nhãn mục tiêu. Với cửa sổ quan sát có độ dài $L$, đầu vào của mô hình tại ngày $t$ được xác định như sau:

\begin{equation}
  X^{p}_{i,t}=[\mathbf{x}^{p}_{i,t-L+1},\ldots,\mathbf{x}^{p}_{i,t}],
  \qquad
  X^{s}_{i,t}=[\mathbf{x}^{s}_{i,t-L+1},\ldots,\mathbf{x}^{s}_{i,t}].
  \label{eq:input-windows}
\end{equation}

Các mã được phân chia theo ngày nhằm bảo đảm rằng mọi mẫu trong cùng một phần chia sử dụng cùng mốc lịch. Đối với ngày không có tin, vector cảm xúc vẫn được tạo theo quy tắc nêu trên. Không loại bỏ ngày một cách có chọn lọc vì thao tác này có thể làm sai lệch phép so sánh giữa giai đoạn có nhiều tin và giai đoạn có ít tin. Đối với những ngày thiếu dữ liệu giá hoặc chưa có đủ lịch sử, lý do loại bỏ phải được ghi rõ; không loại riêng các ngày khó dự báo.

\subsection{Gán nhãn xu hướng phiên kế tiếp}

Lợi suất mục tiêu được tính theo Phương trình~\ref{eq:next-return}. Với cửa sổ huấn luyện $w$, hai phân vị $q^-_w$ và $q^+_w$ được xác định từ lợi suất chỉ thuộc phần huấn luyện. Quy tắc phân lớp được xác định như sau:

\begin{equation}
 y_{i,t+1}=\begin{cases}
 \mathrm{DOWN}, & r_{i,t+1}<q^-_w,\\
 \mathrm{FLAT}, & q^-_w\leq r_{i,t+1}\leq q^+_w,\\
 \mathrm{UP}, & r_{i,t+1}>q^+_w.
 \end{cases}
 \label{eq:trend-label}
\end{equation}

Trong giao thức chính, $q^-_w$ và $q^+_w$ được lấy từ hai phân vị thấp và cao đã được khóa trong cấu hình. Nếu sử dụng ngưỡng cố định, giá trị ngưỡng phải được ghi rõ và không được điều chỉnh dựa trên kết quả kiểm thử. Việc chọn ngưỡng cũng phải cân bằng hai yêu cầu: lớp đi ngang cần phản ánh các biến động nhỏ chưa có hướng rõ rệt, đồng thời phân bố ba lớp không nên mất cân bằng đến mức mô hình có thể đạt điểm cao chỉ bằng cách dự đoán lớp phổ biến. Việc ước lượng ngưỡng trên tập huấn luyện giúp tránh để phân phối lợi suất tương lai quyết định trước tỷ lệ lớp của giai đoạn quá khứ. Phân bố nhãn được báo cáo theo từng mã và từng giai đoạn để người đọc đánh giá độ khó của bài toán.

\subsection{Hợp nhất bằng LSTM hai nhánh}

Nhánh giá và nhánh cảm xúc sử dụng hai bộ mã hóa LSTM độc lập:

\begin{equation}
  \mathbf{h}^{p}_{i,t}=\operatorname{LSTM}_{p}(X^{p}_{i,t}),
  \qquad
  \mathbf{h}^{s}_{i,t}=\operatorname{LSTM}_{s}(X^{s}_{i,t}).
  \label{eq:two-branch-encoders}
\end{equation}

Hai trạng thái cuối được nối với nhau và truyền qua tầng phân loại hợp nhất:

\begin{equation}
  \mathbf{z}_{i,t}=\left[\mathbf{h}^{p}_{i,t};\mathbf{h}^{s}_{i,t}\right],
  \qquad
  \hat{\mathbf{y}}_{i,t+1}=\operatorname{softmax}
  \left(W_2\,\phi(W_1\mathbf{z}_{i,t}+\mathbf{b}_1)+\mathbf{b}_2\right),
  \label{eq:fusion-classifier}
\end{equation}

trong đó $\phi$ là hàm kích hoạt và $\hat{\mathbf{y}}$ biểu diễn xác suất của ba lớp. Mô hình chỉ sử dụng dữ liệu giá gồm nhánh giá và tầng phân loại tương ứng, nhưng không có $\mathbf{h}^{s}$. Khi so sánh hai mô hình, số chiều ẩn, cửa sổ, hạt giống, ngân sách huấn luyện, tiêu chí dừng sớm và cách ước lượng tiền xử lý phải gần như nhau. Chênh lệch về số lượng tham số do nhánh cảm xúc phải được báo cáo để người đọc đánh giá chính xác mức độ khác biệt giữa hai cấu hình. Nếu hai cấu hình còn có những khác biệt khác ngoài sự hiện diện của thông tin cảm xúc, các khác biệt này phải được xem là yếu tố gây nhiễu trong phép loại bỏ đặc trưng.
\section{Thang mô hình dự báo}
\label{sec:pp-mo-hinh}

Đồ án được tổ chức theo nguyên tắc thang tầng. Mỗi tầng giải quyết
một câu hỏi kiểm chứng riêng và chỉ được mở rộng sau khi tầng trước đã có dữ liệu
hợp lệ.

\subsection{Tầng đối chứng cơ sở}

Tầng đầu tiên gồm đối chứng lớp phổ biến nhất và đối chứng ngẫu nhiên. Đối chứng lớp
phổ biến nhất cho biết mức điểm thu được khi mọi quan sát đều được dự đoán thuộc lớp
phổ biến nhất trong tập huấn luyện. Đối chứng ngẫu nhiên được tạo theo quy tắc đã
khóa, sử dụng hạt giống và được chạy trên cùng tập kiểm thử. Khi cần, đối chứng ngẫu
nhiên đều có thể được báo cáo bổ sung để phân biệt ảnh hưởng của phân bố lớp với ảnh
hưởng của mô hình. Hai mốc này không học quy luật từ dữ liệu, nhưng cần thiết để xác
định liệu một kiến trúc có thực sự vượt qua phương pháp dự đoán đơn giản hay chỉ tận dụng sự mất cân bằng lớp.

\subsection{Tầng mô hình cổ điển chỉ giá}

Tầng này sử dụng một bộ phân loại phi tuần tự, nhận vector đặc trưng giá của phiên
quan sát cuối cùng trong cửa sổ. Mục đích là thiết lập mốc so sánh nhằm đánh giá liệu
dữ liệu có chứa tín hiệu đủ rõ trước khi áp dụng mô hình sâu hay không. Cấu hình đã
được thực hiện là hồi quy logistic đa lớp. Rừng ngẫu nhiên và XGBoost được nêu trong
đề cương nhưng không được chạy. Các mô hình này giải quyết cùng một câu hỏi như hồi
quy logistic, trong khi câu hỏi trung tâm của đề tài là đóng góp của cảm xúc; do đó,
ngân sách thực nghiệm được dành cho phép đối chứng này.

Tầng cổ điển được chạy ở cả hai dạng, chỉ giá và giá kèm cảm xúc, trên đúng các dòng
dữ liệu mà mô hình tuần tự sử dụng, qua đó bảo đảm hai họ mô hình được đánh giá trên
cùng tập mẫu. Dạng chỉ giá không nhận bất kỳ cột cảm xúc nào, bao gồm số tin, cờ có
tin và siêu dữ liệu phát sinh từ văn bản.

\subsection{Tầng LSTM chỉ giá}

LSTM chỉ giá nhận $X^p$ trong một cửa sổ cố định và dự báo ba lớp. Đây là đối chứng
cơ sở theo hướng học sâu, được sử dụng để tách biệt câu hỏi ``LSTM có khai thác được
lịch sử giá hay không'' với câu hỏi ``cảm xúc có bổ sung thông tin hay không''. Siêu
tham số được lựa chọn trên tập xác thực và được giữ nguyên khi thực hiện loại bỏ đặc
trưng. Đầu vào chỉ bao gồm các đặc trưng có thể quan sát đến ngày $t$, trong khi nhãn
được tạo từ lợi suất của phiên kế tiếp.

\subsection{Tầng LSTM hai nhánh có cảm xúc}

Tầng này là mức sàn của đề tài. Nhánh cảm xúc nhận chuỗi đặc trưng được tổng hợp theo
ngày, còn nhánh giá nhận chuỗi chỉ báo kỹ thuật. Hai trạng thái được kết hợp bằng
phép nối trước khi được đưa qua tầng phân loại. Nhánh cảm xúc không nhận giá của
phiên tương lai, còn nhánh giá không nhận siêu dữ liệu của bài viết ngoài các đặc
trưng đã được quy định. Thí nghiệm loại bỏ đặc trưng được thực hiện với hai cấu hình:

\begin{enumerate}[nosep]
  \item Chỉ giá: chỉ sử dụng $X^p$;
  \item Giá kết hợp cảm xúc: sử dụng cả $X^p$ và $X^s$.
\end{enumerate}

Hai cấu hình phải sử dụng cùng các ngày huấn luyện, xác thực và kiểm thử; cùng phương
pháp ước lượng các phép tiền xử lý; cùng hạt giống hoặc tập hạt giống; cùng quy tắc
tạo nhãn; và cùng chỉ số đánh giá. Nếu cảm xúc không cải thiện chỉ số, kết quả vẫn
được giữ lại và phân tích theo độ ổn định. Giao thức không được thay đổi nhằm tìm
kiếm một bảng điểm có lợi hơn.

\subsection{Tầng TFT mở rộng}

TFT chỉ được chạy sau khi các đối chứng cơ sở và LSTM hai nhánh đã có số liệu, đồng
thời mốc cắt không còn lỗi. Mô hình nhận các biến thời gian, biến tĩnh hoặc biến quan
sát theo cấu hình được xác định rõ ràng, sử dụng kích thước đầu ra bằng 3 và hàm mất
mát phân loại. Các biến không được khai báo chỉ vì chúng có sẵn trong một bảng dữ
liệu, mà phải được kiểm tra theo thời điểm có thể quan sát. Kết quả TFT chỉ được so
sánh với các tầng trước khi sử dụng cùng các đoạn kiểm thử và quy tắc tạo nhãn. Trong
đồ án này, TFT được cài đặt tối giản bằng PyTorch thuần
(\texttt{TemporalFusionClassifier}), gồm phép nhúng từng biến, mạng chọn biến theo
thời điểm, bộ mã hóa LSTM một tầng, cơ chế chú ý đa đầu với cổng phần dư và đầu phân
loại GRN. Cấu hình hồi quy phân vị của thư viện không được sử dụng cho mục tiêu ba
lớp.

\subsection{Huấn luyện và kiểm soát độ phức tạp}

Mô hình được huấn luyện bằng entropy chéo cho ba lớp và áp dụng dừng sớm theo F1 vĩ
mô trên tập xác thực. Tầng bỏ ngẫu nhiên, suy giảm trọng số và cắt biên độ gradient
có thể được sử dụng để hạn chế quá khớp, nhưng các lựa chọn và số lần thử phải được
ghi lại. Nếu các lớp mất cân bằng, trọng số của hàm mất mát hoặc phương pháp lấy mẫu
lại chỉ được tính từ tập huấn luyện. Đối với mỗi mô hình, điểm kiểm tốt nhất theo chỉ
số đã chọn được lưu lại cùng với hạt giống, cửa sổ, số tham số và thời gian huấn
luyện. Tập kiểm thử không được sử dụng để điều chỉnh kiến trúc sau mỗi lần chạy.
\section{Giao thức đánh giá chống rò rỉ}
\label{sec:pp-danh-gia}

Quy trình đánh giá là thành phần cốt lõi của phương pháp. Việc phân chia tập dữ liệu
theo thời gian chưa đủ để bảo đảm kết quả không bị rò rỉ nếu bộ chuẩn hóa, ngưỡng,
phép chọn đặc trưng hoặc điểm kiểm đã sử dụng dữ liệu tương lai. Do đó, giao thức
được thiết kế theo phương pháp cuốn chiếu theo thời gian với các đoạn kiểm thử tách
biệt, đồng thời kiểm tra toàn bộ các phép biến đổi.

\subsection{Cuốn chiếu theo thời gian}

Tất cả các mã được phân chia theo cùng một mốc ngày. Mỗi cửa sổ cuốn chiếu theo thời
gian gồm tập huấn luyện trong quá khứ, tập xác thực ở giai đoạn kế tiếp và tập kiểm
thử ở giai đoạn sau đó. Có thể sử dụng cửa sổ mở rộng để tăng dần lượng dữ liệu huấn
luyện hoặc cửa sổ trượt để giới hạn độ dài lịch sử. Loại cửa sổ, độ dài từng đoạn và
bước dịch cửa sổ phải được ghi trong bản kê. Mỗi cửa sổ phải chỉ rõ ngày cuối của tập huấn luyện, ngày bắt đầu và kết thúc của tập xác thực, ngày bắt đầu và
kết thúc của tập kiểm thử, cùng số mẫu thực tế sau khi tạo cửa sổ lịch sử.

Cấu hình đã triển khai sử dụng năm cửa sổ mở rộng, trong đó mỗi cửa sổ có 60 ngày
xác thực và 60 ngày kiểm thử; các đoạn kiểm thử không chồng lấn nhau. \textbf{Không
có tập giữ lại cuối tách riêng}, đồng thời \textbf{không có bước tìm kiếm siêu tham
số}: cấu hình mô hình được xác định trước khi chạy và không được thay đổi sau khi
quan sát kết quả. Trong mỗi cửa sổ, tập xác thực chỉ được sử dụng để chọn điểm dừng
huấn luyện, còn đoạn kiểm thử chỉ được đánh giá một lần. Do không có cấu hình nào
được tinh chỉnh trên đoạn kiểm thử, việc bổ sung tập giữ lại thứ hai không tạo thêm
bảo đảm cho quy trình này. Việc nêu rõ đặc điểm trên phản ánh chính xác thiết kế hơn
so với mô tả một tập giữ lại không được sử dụng trong kết quả.

Đối với cửa sổ mở rộng, các mẫu thuộc cùng một ngày không được phân tán giữa tập huấn
luyện và tập kiểm thử. Tất cả các mã trong cùng một ngày phải thuộc cùng một phần
chia nhằm tránh trường hợp trạng thái thị trường của một mã bị tiết lộ cho mã khác
do cách chia không đồng nhất. Mẫu được tạo từ cửa sổ lịch sử cũng phải được phân chia
theo ngày mục tiêu, thay vì chỉ dựa trên ngày bắt đầu cửa sổ. Nếu một mẫu sử dụng
phần lịch sử thuộc giai đoạn trước mốc chia nhưng có nhãn mục tiêu nằm sau mốc này,
mẫu đó phải thuộc phần chia tương ứng với nhãn.

Trong cấu hình hiện tại, các đoạn được phân chia theo ngày mục tiêu, đặt liền nhau
và không có khoảng cấm. Với nhãn của phiên kế tiếp, giá dùng để tạo đặc trưng tại
ngày quan sát đã được biết trước thời điểm bắt đầu phiên mục tiêu; nhãn của đoạn sau
không được sử dụng để huấn luyện đoạn trước. Nếu mở rộng chân trời dự báo hoặc sử
dụng đặc trưng được công bố trễ, quy trình phải bổ sung khoảng cấm và được kiểm tra
lại.

\subsection{Kiểm soát thông tin tại điểm dự báo và các phép biến đổi}

Danh mục kiểm tra rò rỉ bao gồm các hạng mục sau:

\begin{itemize}[nosep]
  \item Mốc thời gian tin và mốc cắt 15 giờ;
  \item Ánh xạ cuối tuần, ngày nghỉ và tin sau giờ;
  \item Giá điều chỉnh và công thức lợi suất mục tiêu;
  \item Bộ chuẩn hóa, bộ điền khuyết và phép chọn đặc trưng chỉ ước lượng trên tập huấn luyện;
  \item Ngưỡng $q^-_w$, $q^+_w$ chỉ tính từ lợi suất tập huấn luyện;
  \item Nhãn cảm xúc dùng để tinh chỉnh và điểm kiểm của từng phần chia;
  \item Loại trùng không để cùng một bài hoặc bản sao gần giống xuất hiện ở hai phần chia;
  \item Lựa chọn siêu tham số chỉ dùng tập huấn luyện và xác thực, không dùng đoạn kiểm thử.
\end{itemize}

Đối với mỗi phần chia, cần lưu một số ví dụ để kiểm tra thủ công, bao gồm tin trước
mốc cắt, tin sau mốc cắt, tin cuối tuần và ngày không có tin. Cần kiểm tra cả trường hợp một bài được gắn với nhiều mã, các bài gần trùng nhau giữa
hai nguồn, ngày đầu chuỗi bị thiếu chỉ báo và mẫu có nhãn ở phiên kế tiếp. Nếu không
thể xác minh một bản ghi, bản ghi đó phải bị loại khỏi phân tích chính và tỷ lệ loại
phải được báo cáo.

\subsection{Các chỉ số đánh giá}

Từ ma trận nhầm lẫn của ba lớp, độ chính xác dương tính, độ bao phủ và F1 được tính
riêng cho từng lớp. F1 vĩ mô được xác định bằng trung bình số học của F1 trên ba lớp:

\begin{equation}
  \operatorname{MacroF1}=\frac{1}{3}\sum_{c=1}^{3}F1_c.
  \label{eq:macro-f1}
\end{equation}

Độ chính xác cân bằng được tính bằng trung bình độ bao phủ của ba lớp. Đối với
AUC-ROC một-chống-còn-lại, mỗi lớp lần lượt được xem là lớp tích cực và hai lớp còn
lại là lớp tiêu cực, sau đó lấy trung bình các giá trị AUC. Độ chính xác hướng được
báo cáo dưới dạng tỷ lệ dự đoán đúng lớp, nhưng không được sử dụng riêng lẻ khi phân
bố lớp mất cân bằng. Các chỉ số được tính trên dự đoán ngoài mẫu của từng phần chia
và sau đó được tổng hợp theo quy tắc đã khóa.

\subsection{Loại bỏ đặc trưng và khoảng tin cậy}

Phép loại bỏ đặc trưng chính so sánh cấu hình chỉ sử dụng giá với cấu hình kết hợp
giá và cảm xúc. Hai mô hình sử dụng cùng các ngày kiểm thử để xác định chênh lệch
theo cặp:

\begin{equation}
  \Delta_m= m(\mathrm{price+sentiment})-m(\mathrm{price\text{-}only}),
  \label{eq:ablation-delta}
\end{equation}

trong đó $m$ là F1 vĩ mô, độ chính xác cân bằng hoặc AUC-ROC một-chống-còn-lại.
Khoảng tin cậy được ước lượng bằng hai phương pháp lấy mẫu lặp và kết quả của cả hai
phương pháp đều được báo cáo. Phương pháp thứ nhất thực hiện lấy mẫu lặp trên các cửa
sổ. Phương pháp thứ hai thực hiện lấy mẫu lặp trên các \emph{ngày} kiểm thử, trong
đó tất cả các mã thuộc cùng một ngày được lấy đồng thời để duy trì tương quan chéo
trong một phiên. Phương pháp thứ hai được ưu tiên khi số cửa sổ nhỏ, vì khoảng tin
cậy được xây dựng từ chỉ năm khối sẽ rất thô và có thể loại trừ 0 một cách ngẫu
nhiên. Số lần lấy mẫu lặp và hạt giống phải được ghi lại. Nếu khoảng tin cậy của
$\Delta_m$ chứa 0, kết luận phù hợp là chưa có bằng chứng rõ ràng cho thấy cảm xúc
cải thiện chỉ số đó trong giao thức hiện tại. Kết quả này không cho phép kết luận
rằng cảm xúc không hữu ích trong mọi bối cảnh.

Phương trình~\eqref{eq:ablation-delta} chưa thể tự tách nguồn gốc của chênh lệch vì
mô hình hai nhánh khác mô hình chỉ giá về cả kiến trúc và thông tin đầu vào. Phép
đối chứng bổ sung giữ nguyên kiến trúc hai nhánh, bài tin, thời điểm và khối lượng
tin, sau đó thay vector xác suất bằng một phân phối tiên nghiệm trung tính hằng số.
Khi đó
\begin{align}
  \substack{\text{hiệu ứng kiến trúc, sự hiện diện}\\\text{và khối lượng tin}} &= m(\text{hai nhánh, tiên nghiệm trung tính})\notag\\
  &\quad - m(\text{chỉ giá}),\\
  \text{đóng góp thông tin phân cực} &= m(\text{hai nhánh, cảm xúc thật})\notag\\
  &\quad - m(\text{hai nhánh, tiên nghiệm trung tính}),
\end{align}
và tổng của hai thành phần bằng đúng $\Delta_m$. Để phép tách có hiệu lực, hai lần
chạy phải sử dụng cùng cấu hình, cùng cửa sổ và cùng hạt giống, đồng thời cho dự đoán
trùng khớp ở cấu hình chỉ giá.

\subsection{Phân tích độ bền và điều kiện phát huy tín hiệu}

Sau khi bảng mức sàn đã ổn định, kết quả có thể được phân tầng theo ngày có nhiều
hoặc ít tin, mức độ biến động thấp hoặc cao, trạng thái tăng hoặc giảm, và mã cổ
phiếu. Các phân tích này chỉ được xem là bổ sung vì làm tăng số lượng phép so sánh.
Số lần thử và tiêu chí phân tầng cần được ghi rõ nhằm tránh lựa chọn riêng nhóm có
kết quả tốt. Kết quả theo mã có thể có phương sai lớn khi một số mã có ít tin; do đó,
cần báo cáo cả số quan sát của từng nhóm.

\subsection{Giải thích mô hình}

Đối với LSTM hai nhánh, phương pháp giải thích cơ bản dựa trên phép loại bỏ đặc trưng
hoặc che khuất. Theo đó, từng nhóm cảm xúc, số tin hoặc một nhóm chỉ báo được che lần
lượt, sau đó đo mức thay đổi của chỉ số. Đối với TFT, mức độ quan trọng của biến và
thành phần chú ý được lưu dưới dạng thông tin bổ sung. Không phương pháp nào trong
số này cho phép diễn giải biến cảm xúc là nguyên nhân của lợi suất. Vì vậy, nội dung
diễn giải nên sử dụng các cụm như ``mô hình đặt trọng số cao hơn'' hoặc ``chỉ số thay
đổi khi che biến'', và không sử dụng cụm ``cảm xúc gây ra biến động''.

\subsection{Tái lập và báo cáo}

Đối với mỗi lần chạy, cần lưu hạt giống, phiên bản Python và thư viện, điểm kiểm, danh
sách mã, mã băm dữ liệu, thời điểm thu thập, cấu hình mô hình, phép chia tập, chỉ số
và số lần thử. Báo cáo phải phân biệt ba trạng thái: đã đo, chưa chạy và không áp
dụng. Các ô chưa có kết quả không được điền số tham chiếu từ công trình khác. Bên
cạnh bảng điểm, báo cáo cần trình bày độ phủ tin, số bản ghi bị loại, phân bố lớp và
các giới hạn của nguồn dữ liệu.

Bảng~\ref{tab:cau-hinh} tóm tắt các quyết định thiết kế chính của quy trình.

\begin{table}[H]
  \centering
  \caption{Tóm tắt các quyết định thiết kế chính.}
  \label{tab:cau-hinh}
  \begin{tabular}{p{4cm} p{9cm}}
    \toprule
    Thành phần & Lựa chọn \\
    \midrule
    Mô hình cảm xúc & PhoBERT-base tinh chỉnh, ba lớp và đầu ra xác suất \\
    Mô hình chuỗi thời gian & LSTM hai nhánh là mức sàn, TFT là tầng mở rộng \\
    Hợp nhất & Nối hai trạng thái cuối rồi phân loại \\
    Phạm vi & 10 mã VN30 trên HOSE, dữ liệu từ năm 2020 đến hết năm 2025 \\
    Mốc cắt & 15 giờ theo múi giờ Việt Nam, tin sau giờ đóng cửa chuyển sang phiên kế tiếp \\
    Nhãn & Lợi suất phiên kế tiếp, ngưỡng chỉ ước lượng trên tập huấn luyện \\
    Đánh giá & Cuốn chiếu theo thời gian, năm đoạn kiểm thử tách biệt, F1 vĩ mô là chỉ số chính \\
    \bottomrule
  \end{tabular}
\noindent\textit{Nguồn: Tác giả xây dựng.}
\end{table}

\chapter{TRÌNH BÀY, ĐÁNH GIÁ VÀ BÀN LUẬN KẾT QUẢ}
\label{ch:ket-qua}

Chương này trình bày các kết quả thực nghiệm đo lường được từ hệ thống. Mỗi con số đều
có thể được truy xuất về một tệp kết quả và một lệnh chạy cụ thể, được liệt kê tại
Mục~\ref{sec:kq-taplap}. Số liệu tham chiếu từ các công trình khác không được sử dụng để
thay thế kết quả thu được trên rổ cổ phiếu của đề tài. Các hạng mục chưa thực hiện được
ghi nhận rõ là chưa thực hiện, thay vì được bổ sung bằng số liệu suy đoán.

Ba lớp xu hướng được xác định theo các phân vị một phần ba và hai phần ba của lợi suất
trong cửa sổ huấn luyện, do đó các lớp gần cân bằng theo thiết kế. Vì vậy,
\textbf{mức ngẫu nhiên của bài toán là $1/3 \approx 0{,}333$, không phải $0{,}5$}.
Con số $0{,}5$ nêu trong đề cương tương ứng với cách phát biểu bài toán hai lớp; toàn bộ
chương này sử dụng mốc ba lớp. Do các lớp gần cân bằng, Macro-F1 và độ chính xác cân bằng
được sử dụng làm các chỉ số chính. Độ chính xác chỉ được báo cáo như một chỉ số bổ sung
và không được sử dụng riêng lẻ.
\section{Kiểm kê dữ liệu đầu vào}
\label{sec:kq-snapshot}

Trước khi huấn luyện, dữ liệu đầu vào được kiểm kê độc lập nhằm phân biệt số liệu quan sát
với kết quả của mô hình. Snapshot hiện tại bao gồm kho tin Vietstock, bảng liên kết tin
với mã cổ phiếu, chuỗi giá của 10 mã HOSE và tập hạt giống phục vụ bước gán nhãn.
Các số liệu trong Bảng~\ref{tab:kq-snapshot} thuộc bước chuẩn bị dữ liệu và không
được diễn giải là kết quả dự báo.

\begin{table}[H]
  \centering
  \caption{Tổng hợp số liệu đầu vào đã kiểm kê (kho tin là toàn bộ lần thu thập, gồm cả năm 2026; số dòng giá là của cửa sổ 2020--2025).}
  \label{tab:kq-snapshot}
  \begin{tabular}{p{8.8cm} c}
    \toprule
    Hạng mục & Số lượng \\
    \midrule
    Bản ghi tin thô có URL duy nhất & 12.140 \\
    Liên kết tin tức và mã trong bảng listings & 14.256 \\
    Mã cổ phiếu được bao phủ & 10 \\
    Ngày dương lịch có ít nhất một tin & 2.101 \\
    Dòng dữ liệu giá OHLCV của 10 mã & 14.990 \\
    Mẫu trong tập hạt giống CafeF & 999 \\
    Mẫu trong miền dữ liệu đã rà soát, hợp nhất & 1.306 \\
    \bottomrule
  \end{tabular}
\noindent\textit{Nguồn: Tác giả tổng hợp từ snapshot dữ liệu của đề tài.}
\end{table}

Kho tin hiện lưu thời điểm đăng và nội dung của các bản ghi. Việc nối tin với
mã cổ phiếu dựa trên bảng listings tạo ra các liên kết theo từng mã; do đó, số liên kết
có thể lớn hơn số URL tin duy nhất. Các thống kê về dữ liệu sau khi loại trùng theo nội dung,
phân bố nhãn và số mẫu cuối cùng trong từng phần chia chỉ được khóa sau khi quy trình xử lý hoàn tất.
Quá trình rà soát nhãn ghi nhận một số ít liên kết nhiễu do trang listings tổng hợp cả tin của
các mã liên quan (ví dụ, bài về MIG xuất hiện trong trang MBB và bài về FOX xuất hiện trong trang
FPT); ba trường hợp này được đánh dấu trong \texttt{annotation\_notes.csv} và được giữ
nguyên trong tập nhãn vì nhãn NEUTRAL phản ánh đúng nội dung bài, nên không phải là lỗi
gán nhãn.
Các chỉ số của PhoBERT và mô hình dự báo được trình bày trong các mục sau khi có bản kê
huấn luyện tương ứng.
\section{Thiết lập thực nghiệm}
\label{sec:kq-setup}

\subsection{Môi trường và cấu hình phần mềm}

Hai nhánh của hệ thống được triển khai trên hai môi trường khác nhau và được ghi nhận riêng biệt. Nhánh cảm xúc tinh chỉnh PhoBERT trên GPU của Kaggle, trong khi nhánh dự báo chạy trên CPU của máy cục bộ do mô hình LSTM ở quy mô này không yêu cầu GPU. Bảng~\ref{tab:kq-setup} trình bày cấu hình thực tế của hai nhánh, được đọc trực tiếp từ
\texttt{outputs/forecast\_merged\_sentiment\_tft/forecast\_results.json} và
\texttt{models/sentiment/merged-refit/manifest.json}.

\begin{table}[H]
  \centering
  \caption{Thiết lập thực nghiệm đã ghi nhận của hai nhánh.}
  \label{tab:kq-setup}
  \begin{tabular}{p{3.6cm} p{4.6cm} p{4.6cm}}
    \toprule
    Hạng mục & Nhánh cảm xúc & Nhánh dự báo \\
    \midrule
    Phần cứng & GPU Kaggle (CUDA) & CPU cục bộ \\
    Python & 3.12.13 & 3.12.14 \\
    Thư viện chính & PyTorch 2.10.0+cu128, Transformers 5.15.1, NumPy 2.0.2 & NumPy 2.5.2, pandas 3.0.5, scikit-learn, PyTorch \\
    Mô hình nền & \texttt{vinai/phobert-base} & LSTM 1 lớp, 32 đơn vị ẩn \\
    Độ dài đầu vào & 256 token, cắt \texttt{head\_tail} & cửa sổ 5 phiên \\
    Kích thước lô & 16 & 128 \\
    Tốc độ học & $2\times10^{-5}$ & $1\times10^{-3}$ \\
    Số vòng huấn luyện & 5 & tối đa 40, dừng sớm kiên nhẫn 6 \\
    Cân bằng lớp & \texttt{inverse\_frequency} theo từng lượt & không (lớp đã cân bằng theo thiết kế) \\
    Hạt giống & gốc 42, mỗi lượt $42+k$ & 42, 43, 44 (bình quân 3 hạt giống) \\
    Chia dữ liệu & phân tầng 5 lượt & 5 cửa sổ cuốn chiếu mở rộng \\
    Hash dữ liệu & nhãn \texttt{18d896d8}\ldots & giá \texttt{19fc9e0c}\ldots, bảng \texttt{3855c4cb}\ldots, cảm xúc \texttt{f5109dc1}\ldots \\
    \bottomrule
  \end{tabular}
\noindent\textit{Nguồn: Tác giả tổng hợp từ bản kê của hai lần chạy.}
\end{table}

Hạt giống được cố định trong các bước chia dữ liệu, khởi tạo mô hình, trộn lô và lấy mẫu lặp. Trong nhánh dự báo, hạt giống dùng để khởi tạo mạng được thiết lập ngay trước khi tạo mô hình. Vì vậy, hai lần chạy độc lập cho kết quả trùng khớp đến từng chữ số, bảo đảm điều kiện để phép so sánh tại Mục~\ref{sec:kq-donggop} có hiệu lực. Các đoạn kiểm thử không được sử dụng để lựa chọn siêu tham số.

\subsection{Thiết kế chia dữ liệu}

Dữ liệu giá được chia theo ngày giao dịch và sử dụng cùng mốc lịch cho các mã. Mỗi cửa sổ cuốn chiếu theo thời gian lần lượt gồm các phần huấn luyện, xác thực và kiểm thử theo đúng thứ tự thời gian. Không có tập giữ lại cuối tách riêng. Cấu hình được cố định trước lần chạy và không được thay đổi sau khi xem kết quả; mỗi đoạn kiểm thử chỉ được đánh giá một lần. Trong từng cửa sổ, scaler, imputer, ngưỡng tạo nhãn và phép chọn đặc trưng giá chỉ được ước lượng từ phần huấn luyện. Quy tắc này hạn chế khả năng điểm số tăng do các phép biến đổi thuộc nhánh giá đã tiếp cận dữ liệu tương lai. Giới hạn riêng của bộ sinh cảm xúc hồi cứu được trình bày tại Mục~\ref{sec:kq-limitations}.

Bảng~\ref{tab:kq-windows} trình bày năm cửa sổ được sử dụng. Phần huấn luyện được mở rộng dần, các khoảng ngày kiểm thử không chồng lấn, và ngưỡng tạo nhãn được ước lượng lại trong từng cửa sổ nên thay đổi theo thời gian.

\begin{table}[H]
  \centering
  \caption{Năm cửa sổ cuốn chiếu và ngưỡng nhãn học từ phần huấn luyện.}
  \label{tab:kq-windows}
  \begin{tabular}{c c c c c c}
    \toprule
    Cửa sổ & Khoảng ngày kiểm thử & Ngưỡng dưới & Ngưỡng trên & Huấn luyện & DOWN/FLAT/UP \\
    \midrule
    1 & 22/10/2024--14/01/2025 & $-0{,}00522$ & $0{,}00565$ & 11.350 & 169/311/120 \\
    2 & 15/01/2025--16/04/2025 & $-0{,}00513$ & $0{,}00557$ & 11.950 & 165/259/176 \\
    3 & 17/04/2025--14/07/2025 & $-0{,}00502$ & $0{,}00535$ & 12.550 & 139/239/222 \\
    4 & 15/07/2025--08/10/2025 & $-0{,}00495$ & $0{,}00530$ & 13.150 & 213/168/219 \\
    5 & 09/10/2025--31/12/2025 & $-0{,}00477$ & $0{,}00535$ & 13.750 & 234/143/223 \\
    \bottomrule
  \end{tabular}
\noindent\textit{Nguồn: \texttt{outputs/forecast\_merged\_sentiment\_tft/forecast\_results.json}.}
\end{table}

Mỗi cửa sổ gồm 600 chuỗi xác thực và 600 chuỗi kiểm thử. Phân bố lớp trong tập kiểm thử lệch dần theo thời gian: cửa sổ 1 tập trung nhiều hơn ở lớp FLAT (311 trên 600), trong khi cửa sổ 5 tập trung nhiều hơn ở hai lớp biên (234 DOWN và 223 UP trên 600). Nguyên nhân là các ngưỡng được học từ dữ liệu quá khứ rồi áp dụng cho dữ liệu tương lai. Do đó, khi mức độ biến động của thị trường tăng, tỷ lệ FLAT giảm. Đây là hệ quả trực tiếp của yêu cầu không sử dụng dữ liệu tương lai và cần được xem xét khi so sánh kết quả giữa các cửa sổ.

\subsection{Cấu hình huấn luyện và tiêu chí chọn mô hình}

Điểm kiểm PhoBERT được lựa chọn dựa trên F1 vĩ mô của tập xác thực thuộc tầng cảm xúc. Mô hình dự báo được lựa chọn theo F1 vĩ mô trên tập xác thực trước khi được đánh giá trên tập kiểm thử của phần chia tương ứng. Hai cấu hình có và không có cảm xúc sử dụng cùng ngân sách thử nghiệm. Khi sử dụng nhiều hạt giống, báo cáo cần trình bày giá trị trung bình, độ lệch chuẩn và quy tắc tổng hợp. Nếu chỉ sử dụng một hạt giống do giới hạn tài nguyên, điều này phải được ghi nhận là một giới hạn của kết quả.

Các đối chứng cơ sở được chạy trước TFT. Quy trình này là cần thiết vì kết quả của mô hình phức tạp chỉ có ý nghĩa khi được so sánh với các mốc đối chứng trên cùng dữ liệu và theo cùng giao thức của đề tài
\cite{Feng2018_181009936v2,Chen2021_210708721v1}.
\section{Thống kê dữ liệu}
\label{sec:kq-data}

\subsection{Quy mô kho ngữ liệu tin}

Bảng~\ref{tab:kq-news-stats} trình bày số lượng bản ghi sau từng cổng xử lý và độ phủ tin tức trong bảng
dữ liệu mã--phiên. Các thống kê này được tính từ \texttt{stf.cli verify} và trường
\texttt{provenance} của \texttt{outputs/forecast\_merged\_sentiment\_tft/forecast\_results.json}.

\begin{table}[H]
  \centering
  \caption{Thống kê kho ngữ liệu tin tức và độ phủ theo phiên.}
  \label{tab:kq-news-stats}
  \begin{tabular}{p{8.4cm} c}
    \toprule
    Chỉ tiêu & Số lượng hoặc tỷ lệ \\
    \midrule
    Bài tin có URL duy nhất trong cửa sổ 2020--2025 & 10.695 \\
    Liên kết tin--mã cổ phiếu trong cửa sổ & 12.624 \\
    Bản ghi có mốc thời gian hợp lệ & 10.695 (100\%) \\
    Bài có nội dung ngoài tiêu đề & 10.695 (100\%) \\
    Liên kết neo vào chính phiên đăng & 6.554 \\
    Liên kết chuyển sang phiên kế tiếp & 6.070 \\
    \quad trong đó rơi sang phiên năm 2026, ngoài cửa sổ & 11 \\
    Liên kết dùng được trong bảng dữ liệu & 12.613 \\
    Dòng mã--phiên của bảng dữ liệu & 14.990 \\
    Dòng mã--phiên có ít nhất một tin & 6.948 (46,35\%) \\
    Dòng mã--phiên không có tin & 8.042 (53,65\%) \\
    \bottomrule
  \end{tabular}
\noindent\textit{Nguồn: \texttt{outputs/forecast\_merged\_sentiment\_tft/forecast\_results.json}, trường \texttt{provenance.alignment\_report\_in\_window} và \texttt{provenance.news\_coverage}.}
\end{table}

Độ phủ được tính theo hai cách và cho kết quả khác nhau; do đó, cần xác định rõ định nghĩa.
Nếu chỉ đối chiếu ngày đăng tin với ngày có phiên giao dịch, độ phủ đạt 43,7\%
(6.556 trên 14.990 dòng mã--phiên có tin đăng đúng ngày). Khi bổ sung điều kiện tin
phải được đăng trước mốc cắt 15:00, độ phủ giảm còn 29,2\% (4.377 dòng). Cách thứ nhất chỉ loại
tin đăng vào cuối tuần và ngày lễ, trong khi cách thứ hai loại thêm các tin đăng sau mốc cắt.
Giá trị 46,35\% trong Bảng~\ref{tab:kq-news-stats} phản ánh độ phủ của quy trình thực tế:
tin đăng sau mốc cắt được chuyển sang phiên kế tiếp thay vì bị loại. Do đó, 6.059 liên kết
chuyển phiên trong cửa sổ vẫn được sử dụng (6.070 liên kết \texttt{next\_session} trừ 11
liên kết neo sang phiên ngoài cửa sổ). Toàn bộ chương này sử dụng định nghĩa thứ hai.

Độ phủ tin tức khác biệt đáng kể giữa các mã: MBB và TCB có trên 2.000 liên kết cho mỗi mã,
trong khi GAS chỉ có 548. Vì vậy, số tin và biến \texttt{has\_news} được giữ làm các đặc trưng
riêng. Ngày không có tin \textbf{không} được gộp vào lớp trung tính mà được gán phân phối tiên
nghiệm trung tính, kèm theo cờ \texttt{has\_news} bằng 0. Đối chứng tại
Mục~\ref{sec:kq-donggop} giữ nguyên các ngày có tin, thời điểm và số lượng tin, đồng thời
chỉ trung tính hóa xác suất cảm xúc.

\subsection{Quy mô dữ liệu giá và phân bố nhãn}

Dữ liệu giá gồm 10 mã trên HOSE, mỗi mã có đúng 1.499 phiên từ 02/01/2020 đến 31/12/2025
và không có giá trị thiếu. Vì vậy, bảng dữ liệu mã--phiên có đủ
$10 \times 1.499 = 14.990$ dòng. Tuy nhiên, số dòng thực tế được đưa vào mô hình nhỏ hơn
giá trị này. Mỗi chuỗi cần một cửa sổ lịch sử đầy đủ gồm 5 phiên và một phiên kế tiếp để
tính nhãn, nên các phiên đầu và phiên cuối của mỗi mã không tạo được mẫu. Cụ thể, ở cửa sổ
đầu tiên, số chuỗi gồm 11.350 chuỗi huấn luyện, 600 chuỗi xác thực và 600 chuỗi kiểm thử,
không đạt 14.990 chuỗi (Bảng~\ref{tab:kq-windows}).

Tệp giá gốc trên đĩa còn chứa các phiên của năm 2026, nằm ngoài cửa sổ nghiên cứu đã khóa.
Các dòng này bị loại khi xây dựng bảng dữ liệu nên không xuất hiện trong bất kỳ số liệu nào
của chương.

Nhãn UP, FLAT và DOWN được xác định theo các ngưỡng tại phân vị một phần ba và hai phần ba
của lợi suất trong cửa sổ huấn luyện. Vì vậy, tỷ lệ ba lớp trên phần huấn luyện gần bằng
nhau theo thiết kế. Trên tập kiểm thử, tỷ lệ các lớp lệch khỏi mức cân bằng do các ngưỡng
được học từ quá khứ và áp dụng cho tương lai. Bảng~\ref{tab:kq-price-stats} tổng hợp phân bố
nhãn của tập kiểm thử theo từng cửa sổ.

\begin{table}[H]
  \centering
  \caption{Phân bố nhãn ba lớp trên tập kiểm thử của năm cửa sổ.}
  \label{tab:kq-price-stats}
  \begin{tabular}{p{4.5cm} c c c c}
    \toprule
    Cửa sổ kiểm thử & DOWN & FLAT & UP & Tổng \\
    \midrule
    1 & 169 & 311 & 120 & 600 \\
    2 & 165 & 259 & 176 & 600 \\
    3 & 139 & 239 & 222 & 600 \\
    4 & 213 & 168 & 219 & 600 \\
    5 & 234 & 143 & 223 & 600 \\
    \midrule
    Tổng & 919 & 1.121 & 960 & 3.000 \\
    \bottomrule
  \end{tabular}
\noindent\textit{Nguồn: \texttt{outputs/forecast\_merged\_sentiment\_tft/forecast\_results.json}.}
\end{table}

Trên toàn bộ 3.000 dòng kiểm thử, tỷ lệ ba lớp lần lượt là 30,6\% DOWN, 37,4\% FLAT và
32,0\% UP. Do các lớp gần cân bằng, mức ngẫu nhiên là 0,333, trong khi phương án dự đoán
toàn bộ quan sát theo lớp phổ biến chỉ đạt Macro-F1 khoảng 0,16. Vì vậy, độ chính xác không
phản ánh đầy đủ khả năng nhận diện hai lớp UP và DOWN; F1 vĩ mô và độ chính xác cân bằng
được lựa chọn làm các chỉ số trung tâm.
\section{Kết quả phân loại cảm xúc}
\label{sec:kq-sentiment}

\subsection{Chất lượng nhãn trong miền dữ liệu}

Các số liệu chính thức trong mục này được tính trên tập \texttt{labeled\_merged.csv}, gồm 1.306
mẫu trong miền dữ liệu. Lượt đầu bao gồm 306 mẫu. Ở lượt mở rộng, cùng người rà soát đã
xác nhận 1.000 mẫu, không thay đổi bất kỳ nhãn sơ bộ nào, sau đó hợp nhất các mẫu này với
lượt đầu. Mô hình ngôn ngữ lớn chỉ được sử dụng để tạo nhãn sơ bộ. Do không có người gán
nhãn thứ hai, quy trình không tính Cohen's hoặc Fleiss' kappa.

Toàn bộ mẫu trong tập có \texttt{annotation\_source} bằng \texttt{human\_reviewed} và
\texttt{annotation\_status} bằng \texttt{REVIEWED}. Xác thực chéo theo tầng được thực
hiện trên tập dữ liệu này. Mục~\ref{sec:kq-mornhan} trình bày cách sử dụng các tầng nhằm
duy trì tính không chệch của ước lượng trên tập kiểm thử.

\begin{table}[H]
  \centering
  \caption{Thông tin gán nhãn tập cảm xúc trong miền dữ liệu.}
  \label{tab:kq-annotation}
  \begin{tabular}{p{7.6cm} c c}
    \toprule
    Chỉ tiêu & Lượt 1 & Lượt 2, đã rà soát \\
    \midrule
    Số người rà soát & 1 & 1 \\
    Số mẫu & 306 & 1.000 \\
    NEGATIVE & 31 (10,1\%) & 163 (16,3\%) \\
    NEUTRAL & 204 (66,7\%) & 567 (56,7\%) \\
    POSITIVE & 71 (23,2\%) & 270 (27,0\%) \\
    Dùng cho chỉ số đang trình bày & Có, lịch sử & Có, xác thực chéo theo tầng \\
    \bottomrule
  \end{tabular}
\noindent\textit{Nguồn: \texttt{to\_label\_r1.csv} và \texttt{to\_label\_batch2\_ai.csv}.}
\end{table}

Năm tập kiểm thử ngoài có tổng cộng 59 mẫu NEGATIVE, với 11 hoặc 12 mẫu trong mỗi lượt.
Khoảng tin cậy của độ bao phủ đối với lớp này vẫn rộng hơn so với hai lớp còn lại, với
nửa độ rộng xấp xỉ $\pm0{,}128$ quanh tỷ lệ 0,5. Tuy nhiên, số lượng này cho phép đo độ
bao phủ của lớp NEGATIVE, thay vì để lớp này vắng mặt trong nhiều lượt như ở kết quả
lịch sử sử dụng 306 mẫu.

\subsection{Chỉ số mô hình cảm xúc}
\label{sec:kq-sentiment-metrics}

PhoBERT được đánh giá bằng xác thực chéo theo tầng 5 lượt trên 1.306 mẫu đã được rà soát.
Hai tầng làm giàu chỉ được đưa vào phần huấn luyện; tập kiểm thử ngoài chỉ bao gồm dữ liệu
từ \texttt{eval\_random} và \texttt{baseline\_random}. Năm tập kiểm thử có tổng cộng 656
dòng, trong đó mỗi lượt gồm 131 hoặc 132 dòng. Trong mỗi lượt, 118 dòng xác thực tiếp tục
được rút từ phần huấn luyện ngoài để lựa chọn điểm kiểm. Cấu hình thực nghiệm là
\texttt{title\_context} với phương pháp cắt \texttt{head\_tail}, năm vòng huấn luyện và
trọng số lớp \texttt{inverse\_frequency} được tính riêng trên phần huấn luyện.

\begin{table}[H]
  \centering
  \caption{Kết quả phân loại cảm xúc, trung bình trên 5 lượt.}
  \label{tab:kq-sentiment}
  \begin{tabular}{l c c c c}
    \toprule
    Cấu hình và tập nhãn & Độ chính xác & F1 vĩ mô & Độ chính xác cân bằng & F1 NEG \\
    \midrule
    Không trọng số lớp, 3 vòng, 306 mẫu & 0,6667 & 0,2667 & 0,3333 & 0,0000 \\
    Có trọng số lớp, 5 vòng, 306 mẫu & 0,7353 & 0,5176 & 0,5473 & 0,0767 \\
    Có trọng số lớp, 5 vòng, 1.306 mẫu & \textbf{0,8125} & \textbf{0,7298} & \textbf{0,7689} & \textbf{0,6036} \\
    \bottomrule
  \end{tabular}
\noindent\textit{Nguồn: \texttt{outputs/sentiment-cv-merged/cv\_results.json}; hai dòng 306 mẫu là kết quả lịch sử.}
\end{table}

Cấu hình không sử dụng trọng số lớp trên 306 mẫu là đường cơ sở lịch sử. Kết quả sử dụng
trọng số nghịch đảo tần suất trên 1.306 mẫu được thu nhận từ một tập dữ liệu và giao thức
mới, nên không được xem là mô hình cảm xúc cuối cùng. Khác biệt giữa kết quả trên hai cỡ
tập dữ liệu không được diễn giải là hiệu ứng riêng của trọng số lớp.

Cấu hình đường cơ sở không sử dụng trọng số lớp trên 306 mẫu dự đoán toàn bộ mẫu là
NEUTRAL. Độ chính xác cân bằng bằng $1/3$ trong cả năm lượt, trong khi F1 của NEGATIVE và
POSITIVE bằng 0. Lần chạy mở rộng có sử dụng trọng số không còn biểu hiện suy biến. F1 vĩ
mô trung bình đạt 0,7298 với độ lệch chuẩn 0,0423, còn độ chính xác cân bằng đạt 0,7689
với độ lệch chuẩn 0,0497.

\begin{table}[H]
  \centering
  \caption{F1 và độ bao phủ theo lớp, cấu hình 1.306 mẫu có trọng số, trung bình 5 lượt.}
  \label{tab:kq-sentiment-class}
  \begin{tabular}{l c c}
    \toprule
    Lớp & F1 & Độ bao phủ \\
    \midrule
    NEGATIVE & 0,6036 & 0,6803 \\
    NEUTRAL & 0,8816 & 0,8359 \\
    POSITIVE & 0,7041 & 0,7906 \\
    \bottomrule
  \end{tabular}
\noindent\textit{Nguồn: \texttt{outputs/sentiment-cv-merged/cv\_results.csv}.}
\end{table}

Độ bao phủ của lớp NEGATIVE theo từng lượt lần lượt là $0{,}7500$; $0{,}8182$;
$0{,}6667$; $0{,}5000$; $0{,}6667$. Không có lượt nào bỏ qua hoàn toàn lớp này. Tuy
nhiên, do năm tập kiểm thử ngoài chỉ có tổng cộng 59 mẫu, mức độ bất định đối với lớp
NEGATIVE vẫn lớn hơn so với các lớp còn lại.

Artifact hiện có chỉ đánh giá cấu hình \texttt{title\_context} với
\texttt{head\_tail}; do đó, kết quả này không cung cấp phép so sánh mới giữa các phương án
đầu vào hoặc các chiến lược cắt token.

Các điểm kiểm của từng lượt là artifact của quy trình xác thực chéo và không được lựa chọn
dựa trên chỉ số của tập kiểm thử ngoài để chấm điểm tin tức. Sau khi khóa cấu hình, điểm
kiểm triển khai được tinh chỉnh lại trên toàn bộ 1.306 nhãn với năm epoch cố định, không
rút outer holdout và không báo cáo chỉ số kiểm thử mới. Manifest của điểm kiểm xác nhận
đúng mã băm của tập nhãn, phân bố gồm 194 NEGATIVE, 771 NEUTRAL và 341 POSITIVE, cùng
môi trường GPU đã sử dụng.

Các bảng dự báo lịch sử được tạo bằng bộ sinh 306 mẫu với \texttt{score-news} và
\texttt{--context-chars 400}. Tổng cộng 14.256 liên kết tin--mã đã được gán xác suất cho
ba lớp, trong đó 12.624 liên kết thuộc cửa sổ 2020--2025 và 12.613 liên kết được đưa vào
bảng dự báo. Phân bố nhãn đa số trong kết quả lịch sử gồm 2.095 NEGATIVE, 6.776 NEUTRAL
và 3.753 POSITIVE. Nhánh dự báo luôn sử dụng phân phối xác suất mềm thay vì nhãn cứng.
Các số liệu lịch sử này chỉ được dùng để giải thích artifact cũ và không được quy cho điểm
kiểm 1.306 mẫu.

Điểm kiểm mới đã chấm lại đúng 14.256 liên kết bằng \texttt{title\_context},
\texttt{head\_tail} và giới hạn 400 ký tự. Trong số này, 12.624 liên kết thuộc cửa sổ
2020--2025; 12.613 liên kết được neo vào một phiên bên trong cửa sổ theo mốc 15 giờ, còn
11 liên kết được đăng vào cuối ngày 31/12/2025 rơi sang phiên năm 2026 nên bị loại. 1.632
liên kết còn lại được đăng trong năm 2026 và nằm ngoài cửa sổ nghiên cứu. Bảng dự báo có
tin tại 6.948 trên 14.990 dòng mã--phiên. Lần chạy mới và đối chứng ghép cặp sử dụng cùng
các bài tin, thời điểm và số lượng tin; đối chứng chỉ thay ba xác suất của mỗi bài bằng
prior trung tính. Do điểm kiểm được tinh chỉnh trên toàn bộ tập nhãn, trong đó có các bài
năm 2025, kết quả dự báo này mang tính hồi cứu và không được xem là dự báo ngoài mẫu hoàn
toàn theo giao thức tại Chương~\ref{ch:phuong-phap}.
\section{Kết quả dự báo xu hướng và loại bỏ đặc trưng}
\label{sec:kq-forecast}

Toàn bộ các mô hình trong thang so sánh được chạy trên cùng năm cửa sổ cuốn chiếu, cùng ba hạt giống và
cùng tập dòng kiểm thử của nhánh giá. Mỗi cửa sổ gồm 600 chuỗi kiểm thử, tương ứng với tổng số 3.000 dòng
đánh giá cho mỗi hạt giống. Các giá trị trong Bảng~\ref{tab:ket-qua} là trung bình của 15 lần
chạy (5 cửa sổ $\times$ 3 hạt giống), kèm theo độ lệch chuẩn của F1 vĩ mô.

\begin{table}[H]
  \centering
  \caption{So sánh thang mô hình dự báo xu hướng phiên kế tiếp.}
  \label{tab:ket-qua}
  \begin{tabular}{l c c c c}
    \toprule
    Mô hình & F1 vĩ mô & Độ chính xác cân bằng & Độ chính xác & AUC-ROC OvR \\
    \midrule
    Đối chứng lớp phổ biến & 0,1453 $\pm$ 0,0264 & 0,3333 & 0,2820 & --- \\
    Đối chứng ngẫu nhiên & 0,3289 $\pm$ 0,0184 & 0,3343 & 0,3339 & --- \\
    Hồi quy logistic, chỉ giá & 0,3551 $\pm$ 0,0281 & 0,3792 & 0,4220 & 0,5601 \\
    Hồi quy logistic, giá và cảm xúc & 0,3667 $\pm$ 0,0250 & 0,3872 & 0,4270 & 0,5570 \\
    LSTM, chỉ giá & 0,3756 $\pm$ 0,0321 & 0,3915 & \textbf{0,4341} & \textbf{0,5714} \\
    \textbf{LSTM hai nhánh, giá và cảm xúc} & \textbf{0,3869} $\pm$ 0,0272 & \textbf{0,3953} & 0,4244 & 0,5655 \\
    TFT, chỉ giá & 0,3619 $\pm$ 0,0352 & 0,3776 & 0,4123 & 0,5596 \\
    TFT, giá và cảm xúc & 0,3597 $\pm$ 0,0201 & 0,3733 & 0,3936 & 0,5527 \\
    \bottomrule
  \end{tabular}
\noindent\textit{Nguồn: \texttt{outputs/forecast\_merged\_sentiment\_tft/forecast\_metrics.csv}.}
\end{table}

LSTM chỉ sử dụng giá đạt F1 vĩ mô 0,3756, so với 0,3289 của đối chứng ngẫu nhiên và 0,1453 của
đối chứng lớp phổ biến; AUC-ROC một-chống-còn-lại đạt 0,5714. Báo cáo không thực hiện
kiểm định ghép cặp với hai đối chứng, do đó kết quả này chỉ phản ánh so sánh mô tả giữa các điểm
quan sát. Kết quả của mô hình hai nhánh được giữ lại để mô tả lần chạy hồi cứu, nhưng không
cấu thành bằng chứng độc lập về giá trị của tin hoặc cảm xúc vì bộ sinh xác suất đã được huấn luyện
từ các nhãn tương lai so với những cửa sổ kiểm thử.

TFT đạt F1 vĩ mô 0,3619 khi chỉ sử dụng giá và 0,3597 khi bổ sung cảm xúc, thấp hơn
hai mô hình LSTM tương ứng (0,3756 và 0,3869). Trên bảng dữ liệu mã và ngày gồm 14.990
dòng, với cửa sổ đầu vào dài 5 phiên, kiến trúc cổng và chú ý của TFT không tạo ra
lợi thế so với LSTM hai nhánh. Kết quả này phù hợp với nhận định rằng mô hình có nhiều
tham số cần lượng dữ liệu lớn hơn để phát huy cơ chế chọn biến. Do đó, TFT được ghi nhận
là một tầng mô hình đã được thực nghiệm nhưng không vượt mức sàn, thay vì một hạng mục chưa được triển khai.

\subsection{Báo cáo chỉ số theo lớp}

Các dự đoán chi tiết theo từng mã, ngày, cửa sổ và hạt giống được lưu tại
\path{outputs/forecast_merged_sentiment_tft/forecast_predictions.csv}. Chỉ số tổng hợp
trong Bảng~\ref{tab:ket-qua} là trung bình của 15 lần chạy và không được thay thế bằng một bảng
dự đoán gộp, vì phép gộp sẽ làm thay đổi trọng số giữa các cửa sổ. Các lớp vẫn được xác định bằng ngưỡng học
từ từng tập huấn luyện. Vì vậy, F1 vĩ mô và độ chính xác cân bằng được sử dụng làm các chỉ số chính
thay cho độ chính xác đơn lẻ.

\begin{table}[H]
  \centering
  \small
  \caption{Chỉ số theo lớp, trung bình 15 lượt chạy; nhánh hai có cảm xúc là kết quả hồi cứu.}
  \label{tab:kq-theo-lop}
  \begin{tabular}{l rrr rrr}
    \toprule
    & \multicolumn{3}{c}{LSTM chỉ giá} & \multicolumn{3}{c}{LSTM hai nhánh} \\
    \cmidrule(lr){2-4}\cmidrule(lr){5-7}
    Lớp & P & R & F1 & P & R & F1 \\
    \midrule
    DOWN & 0,3775 & 0,2472 & 0,2862 & 0,3538 & 0,2741 & 0,3052 \\
    FLAT & 0,4385 & 0,6193 & 0,5106 & 0,4415 & 0,5760 & 0,4985 \\
    UP   & 0,3859 & 0,3081 & 0,3301 & 0,3921 & 0,3357 & 0,3571 \\
    \bottomrule
  \end{tabular}
\end{table}
\noindent\textit{Nguồn: \texttt{outputs/forecast\_merged\_sentiment\_tft/per\_class\_metrics.csv}, trung bình trên 15 tổ hợp cửa sổ và hạt giống.}

Các giá trị là trung bình của chỉ số được tính riêng cho từng tổ hợp cửa sổ và hạt giống,
nên nhất quán với Bảng~\ref{tab:ket-qua}; chúng không được suy ra từ một bảng dự đoán
đã gộp.

\subsection{Đo phần đóng góp của cảm xúc}
\label{sec:kq-donggop}

Phép so sánh trực tiếp giữa mô hình hai nhánh và mô hình một nhánh chỉ sử dụng giá \textbf{trộn lẫn} ảnh hưởng của
nhánh mạng bổ sung với ảnh hưởng của sự hiện diện, khối lượng và xác suất của tin. Đối chứng sử dụng chính các liên
kết tin đã chấm, giữ nguyên ngày, thời điểm, số lượng tin và cờ \texttt{has\_news}, nhưng
thay ba xác suất của mỗi bài bằng phân phối tiên nghiệm trung tính. Do mô hình chỉ sử dụng giá
không có đặc trưng về sự hiện diện hoặc khối lượng tin, hiệu số thứ nhất dưới đây không biểu thị hiệu
ứng kiến trúc thuần:
\begin{align}
  E_{\text{kiến trúc--tin}} &= m_{\text{trung tính}} - m_{\text{chỉ giá}}, \\
  I_{\text{phân cực}} &= m_{\text{cảm xúc thật}} - m_{\text{trung tính}}.
\end{align}

Tệp kết quả kiểm tra sự đồng nhất của mã băm tệp tin đầu vào, số dòng, mã băm giá, cấu
hình và khóa dự đoán theo cửa sổ, hạt giống, mã và ngày. Tệp cũng yêu cầu các dự đoán và
chỉ số của nhánh chỉ sử dụng giá phải trùng khớp. Vì vậy, phép ghép cặp không trộn lẫn khác biệt về độ
phủ tin với khác biệt về xác suất cảm xúc.

\begin{table}[H]
  \centering
  \caption{Tách hiệu ứng kiến trúc--sự hiện diện--khối lượng tin và đóng góp thông tin phân cực, F1 vĩ mô.}
  \label{tab:kq-ablation}
  \begin{tabular}{p{7.6cm} c}
    \toprule
    Mức & F1 vĩ mô \\
    \midrule
    LSTM chỉ giá & 0,3756 \\
    LSTM hai nhánh, phân phối tiên nghiệm trung tính & 0,3790 \\
    LSTM hai nhánh, cảm xúc thật & 0,3869 \\
    \midrule
    Hiệu ứng kiến trúc, sự hiện diện và khối lượng tin & $+0{,}0034$ \\
    \textbf{Đóng góp thông tin phân cực} & $\mathbf{+0{,}0080}$ \\
    Hiệu số trực tiếp (trộn lẫn hai thành phần) & $+0{,}0113$ \\
    \bottomrule
  \end{tabular}
\noindent\textit{Nguồn: \path{outputs/forecast_merged_sentiment_tft/information_gain.json}.}
\end{table}

Trung bình trên năm cửa sổ, cảm xúc thật cao hơn đối chứng trung tính $0{,}0080$ điểm F1 vĩ
mô; bốn trong năm cửa sổ có hiệu số dương. Hiệu ứng của kiến trúc và độ phủ tin đóng góp thêm
$+0{,}0034$. Do đó, hiệu số trực tiếp $+0{,}0113$ bao gồm cả hai thành phần và cần được
phân tách thay vì chỉ báo cáo một giá trị duy nhất. Tuy nhiên, dấu dương của ước lượng điểm
chưa đủ để kết luận rằng mức cải thiện có thể tái lập.

Bảng~\ref{tab:kq-ablation-ci} báo cáo khoảng tin cậy theo hai phương pháp lấy mẫu lặp quanh
\textit{cùng một} ước lượng điểm, là trung bình của năm hiệu số được tính riêng trong từng cửa sổ.
Phương pháp thứ nhất lấy mẫu lặp trên năm cửa sổ. Phương pháp thứ hai lấy mẫu lặp trên 300 \textit{ngày}
kiểm thử; mọi mã thuộc cùng một ngày được lấy đồng thời để duy trì tương quan chéo
trong phiên. Các ngày được lấy lại \textit{bên trong cửa sổ của chính nó} trước khi tính trung
bình trên năm cửa sổ. Khoảng theo ngày được ưu tiên do chỉ có năm khối cửa sổ.

\begin{table}[H]
  \centering
  \caption{Đóng góp thông tin của cảm xúc: một ước lượng điểm, hai cách lấy mẫu lặp.}
  \label{tab:kq-ablation-ci}
  \begin{tabular}{l c c c}
    \toprule
    Chỉ số & Hiệu số & Khoảng 95\% theo cửa sổ & Khoảng 95\% theo ngày \\
    \midrule
    F1 vĩ mô & $+0{,}0080$ & $[-0{,}0000; +0{,}0160]$ & $[-0{,}0025; +0{,}0186]$ \\
    Độ chính xác cân bằng & $+0{,}0017$ & $[-0{,}0031; +0{,}0070]$ & $[-0{,}0072; +0{,}0108]$ \\
    Độ chính xác & $-0{,}0104$ & $[-0{,}0307; +0{,}0050]$ & $[-0{,}0206; -0{,}0010]$ \\
    \bottomrule
  \end{tabular}
\noindent\textit{Nguồn: \texttt{outputs/forecast\_merged\_sentiment\_tft/information\_gain.json}.}
\end{table}

Một phiên bản trước của thí nghiệm đã gộp cả năm cửa sổ vào một ma trận nhầm lẫn duy nhất, sau đó
tính khoảng tin cậy quanh hiệu số gộp. Ước lượng gộp là một ước lượng hợp lệ và đã được
báo cáo cùng khoảng tương ứng, nhưng đây là \textit{một đại lượng khác}. F1 vĩ mô là hàm phi
tuyến của ma trận nhầm lẫn, nên hiệu số gộp không bằng trung bình của các hiệu số theo cửa sổ.
Điều này yêu cầu người đọc theo dõi đồng thời hai kết quả. Bảng hiện tại sử dụng một đại lượng và
gắn với đại lượng đó hai khoảng tin cậy. Đây là thay đổi về cách trình bày và \textbf{không} phải bằng chứng
cho thấy con số gộp trước đó không chính xác.

Cả hai khoảng đều chịu hai giới hạn. Các ngày được rút độc lập, vì vậy không phương pháp nào thực hiện lấy mẫu
lặp chuỗi theo khối; đồng thời, cả hai đều không mô hình hóa sự phụ thuộc giữa các phiên liên tiếp.
Tệp kết quả còn ghi nhận một giá trị $p$ một phía. Tuy nhiên, chiều của giả thuyết \textbf{không}
được xác định trước khi quan sát kết quả, nên giá trị này chỉ mang tính tham khảo và không được sử dụng để
tuyên bố ý nghĩa thống kê.

Khoảng theo ngày của F1 vĩ mô và độ chính xác cân bằng đều chứa 0. Đối với độ chính xác,
khoảng $[-0{,}0206; -0{,}0010]$ nằm hoàn toàn dưới 0, nhưng khoảng theo cửa sổ của cùng chỉ
số vẫn chứa 0. Hai phương pháp lấy mẫu lặp giả định các cấu trúc phụ thuộc khác nhau, nên không
tạo thành một kiểm định thống nhất. Theo quy tắc đã xác định tại
Chương~\ref{ch:phuong-phap}, kết luận phù hợp là \textbf{chưa có bằng chứng đủ mạnh về
cải thiện}. Kết quả không hỗ trợ kết luận rằng ``cảm xúc làm mô hình tệ hơn'' hoặc ``cảm xúc chắc chắn có ích''.

\subsection{Đánh giá ngoài mẫu với điểm kiểm đóng băng}
\label{sec:kq-point-in-time}

Kết quả tại các mục trên sử dụng điểm kiểm cảm xúc được tinh chỉnh trên toàn bộ 1.306 nhãn,
trong đó có các nhãn được công bố sau ngày kiểm thử đầu tiên. Do đó, các kết quả này mang tính hồi cứu và
không cấu thành bằng chứng ngoài mẫu. Mục này báo cáo lượt đánh giá ngoài mẫu đã được thực hiện.

Ngày quan sát kiểm thử đầu tiên của thang dự báo là 21/10/2024. Lệnh
\texttt{sentiment-refit --before-date 2024-10-21} chỉ giữ lại các nhãn được công bố
\emph{nghiêm ngặt trước} mốc này, gồm 1.049 trên 1.306 nhãn, và ghi
\texttt{label\_cutoff} vào bản kê điểm kiểm. Khi chạy dự báo, lệnh \texttt{forecast}
đối chiếu mã băm của bản kê với mã băm được lưu khi chấm tin, so sánh mốc cắt với ngày
kiểm thử đầu tiên, sau đó ghi \texttt{point\_in\_time} cùng lý do vào tệp kết quả. Cả
hai artifact của lượt ngoài mẫu đều có \texttt{point\_in\_time: true} với lý do
\texttt{verified}.

Phạm vi của chứng nhận này cần được xác định rõ. Điểm kiểm đóng băng chỉ loại bỏ các
\emph{nhãn huấn luyện} được công bố sau mốc cắt. Việc lựa chọn cấu hình
(\texttt{title\_context}, \texttt{head\_tail}, \texttt{inverse\_frequency}) vẫn
dựa trên xác thực chéo được thực hiện trên toàn bộ 1.306 nhãn. Vì vậy, lượt ngoài mẫu chỉ
kiểm chứng được khía cạnh rò rỉ nhãn, nhưng không kiểm chứng được rò rỉ trong bước chọn mô hình.

Một điểm kiểm thứ ba được sử dụng để tách hai nguyên nhân có thể làm thay đổi đóng góp cảm xúc giữa
hai lượt. Lệnh \texttt{sentiment-refit --sample-like} rút đúng số nhãn của từng lớp mà
điểm kiểm đóng băng đã sử dụng (167 NEGATIVE, 616 NEUTRAL, 266 POSITIVE), nhưng lấy mẫu từ
tất cả các mốc thời gian. Việc đối chiếu ba lượt, gồm đầy đủ 1.306, đóng băng 1.049 và đối
chứng 1.049, cho phép phân biệt ảnh hưởng của việc loại bỏ nhãn tương lai với ảnh
hưởng của việc huấn luyện trên ít dữ liệu hơn. Bản kê của lượt đối chứng ghi rõ
chiến lược lấy mẫu, hạt giống và mã băm của danh sách dòng được chọn, đồng thời đánh dấu
\texttt{point\_in\_time: false} để tránh diễn giải nhầm đây là một lượt sạch rò rỉ thứ
hai. Lượt đối chứng hiện chỉ có một hạt giống (7), nên kết quả mang tính thăm dò. Tập con
được rút là một yếu tố ngẫu nhiên và một lần rút duy nhất không thể tách ảnh hưởng
của kích thước tập khỏi ảnh hưởng của một lần rút không thuận lợi.

\begin{table}[H]
  \centering
  \caption{Đóng góp thông tin phân cực (F1 vĩ mô) của LSTM hai nhánh theo ba điểm kiểm cảm xúc, cùng cửa sổ kiểm thử.}
  \label{tab:kq-pit}
  \begin{tabular}{l c c c}
    \toprule
    Điểm kiểm & Hiệu số & KTC 95\% theo cửa sổ & KTC 95\% theo ngày \\
    \midrule
    Đầy đủ 1.306 nhãn (hồi cứu) & $+0{,}0080$ & $[-0{,}0000; +0{,}0160]$ & $[-0{,}0025; +0{,}0186]$ \\
    Đóng băng 1.049 (ngoài mẫu) & $+0{,}0094$ & $[-0{,}0022; +0{,}0243]$ & $[-0{,}0009; +0{,}0196]$ \\
    Đối chứng 1.049 (hạt giống 7) & $+0{,}0028$ & $[-0{,}0070; +0{,}0127]$ & $[-0{,}0074; +0{,}0133]$ \\
    \bottomrule
  \end{tabular}
\noindent\textit{Nguồn: \texttt{information\_gain.json} trong \path{outputs/forecast_merged_sentiment_tft}, \path{outputs/forecast_pit_sentiment} và \path{outputs/forecast_sizematched_s7_sentiment}.}
\end{table}

Cả ba ước lượng điểm đều dương và nằm trong khoảng $+0{,}003$ đến $+0{,}009$ F1 vĩ mô. Trong
lần rút sử dụng hạt giống 7, điểm kiểm đóng băng theo thời gian cho ước lượng cao hơn điểm kiểm
đối chứng cùng kích thước ($+0{,}0094$ so với $+0{,}0028$). Tuy nhiên, đây là kết quả thăm dò từ một
lần rút mẫu duy nhất và chưa cho phép tách chắc chắn ảnh hưởng của thời điểm huấn luyện khỏi ảnh hưởng của
thành phần mẫu được rút.

Cả ba khoảng theo ngày đều chứa 0. Khoảng của lượt đóng băng gần chạm 0 tại cận dưới
($-0{,}0009$), nhưng ba lượt này thuộc các phép so sánh được thực hiện trên cùng bộ
dữ liệu; nếu tính cả hai kiến trúc thì có sáu phép so sánh. Vì vậy, sau khi hiệu chỉnh đa so sánh, không còn cơ sở
để tuyên bố ý nghĩa thống kê. Kết luận không thay đổi: \textbf{chưa có bằng chứng đủ mạnh
về cải thiện}, đồng thời dữ liệu hiện tại cũng không loại trừ khả năng tồn tại một mức cải thiện nhỏ.

Đối với TFT, đóng góp thông tin lần lượt là $-0{,}0057$ (hồi cứu), $-0{,}0105$ (đóng băng) và
$-0{,}0020$ (đối chứng), với các khoảng theo ngày tương ứng là $[-0{,}0161; +0{,}0047]$,
$[-0{,}0209; +0{,}0013]$ và $[-0{,}0116; +0{,}0083]$. Không có lợi ích quan sát được với bất
kỳ bộ sinh nào, nhưng cũng không có mức suy giảm vượt quá ngưỡng nhiễu. Kết quả này phù hợp với
việc TFT chỉ sử dụng giá đạt kết quả thấp hơn LSTM chỉ sử dụng giá trên quy mô dữ liệu này.

\subsection{Tương quan hạng của tín hiệu, không qua mô hình}
\label{sec:kq-ic}

Một phép loại bỏ cho kết quả rỗng không thể phân biệt giữa hai khả năng: bộ ước lượng chưa khai
thác được tín hiệu hoặc bản thân tín hiệu yếu. Mục này bổ sung một mảnh bằng chứng độc lập với kiến trúc thông qua tương quan hạng Spearman giữa điểm cảm xúc theo ngày và lợi suất thực hiện,
\textbf{không khớp bất kỳ mô hình nào}. Lệnh \texttt{forecast-ic} tính toán trên 6.950 dòng
mã--ngày có tin thuộc 1.483 phiên, sử dụng điểm kiểm đóng băng; khoảng tin cậy được lấy mẫu lặp
theo toàn bộ phiên với 2.000 lần rút.

\begin{table}[H]
  \centering
  \caption{Tương quan hạng Spearman giữa \texttt{sent\_pos\_minus\_neg} và lợi suất thực hiện, điểm kiểm đóng băng.}
  \label{tab:kq-ic}
  \begin{tabular}{l c c}
    \toprule
    Lợi suất đối chiếu & Tương quan hạng & Khoảng tin cậy 95\% \\
    \midrule
    Cùng phiên, lợi suất thô & $+0{,}0130$ & $[-0{,}0158; +0{,}0406]$ \\
    Cùng phiên, lợi suất vượt trội & $+0{,}0233$ & $[+0{,}0004; +0{,}0465]$ \\
    Phiên kế tiếp, lợi suất thô & $+0{,}0110$ & $[-0{,}0164; +0{,}0399]$ \\
    Phiên kế tiếp, lợi suất vượt trội & $+0{,}0030$ & $[-0{,}0207; +0{,}0269]$ \\
    \bottomrule
  \end{tabular}
\noindent\textit{Nguồn: \path{outputs/signal_ic_pit.json}; $n = 6.950$ và $6.948$ dòng.}
\end{table}

Tương quan hạng cắt ngang theo ngày được tính bằng cách xếp hạng trong từng phiên rồi lấy trung bình,
tương ứng với cách sử dụng thực tế của một tín hiệu định lượng. Giá trị thu được là $+0{,}0062$, với sai số chuẩn $0{,}0159$
trên 1.002 phiên ($t = +0{,}39$). Với điểm kiểm hồi cứu, giá trị tương ứng là $-0{,}0047$
($t = -0{,}29$), tức là dấu thay đổi giữa hai bộ sinh và độ lớn hoàn toàn nằm trong phạm vi nhiễu.

\paragraph{Mục này không chứng minh điều gì.}
Các giới hạn cần được nêu rõ trước khi diễn giải. Hệ số Spearman gần 0 chỉ cho thấy \textbf{chưa phát hiện được bằng chứng} về liên hệ
\textbf{đơn điệu} giữa \textbf{đúng điểm số vô hướng này} và lợi suất; nó không bác bỏ sự tồn tại của liên hệ đó, đúng như nguyên tắc đã áp dụng cho phép
cắt bỏ đặc trưng. Hệ số này không cung cấp cận trên cho khả năng dự báo và
không đo lường lượng thông tin. Một phản ví dụ tiêu chuẩn là trường hợp
$x$ có phân phối đối xứng quanh 0 và $y = x^2$: biến $y$ được xác định hoàn toàn bởi $x$,
nhưng tương quan hạng giữa hai biến bằng 0. Tương tự, kết quả ở đây không nói gì về:

\begin{itemize}[nosep]
  \item một quan hệ phi tuyến, có điều kiện hoặc có tương tác với đặc trưng giá;
  \item một \emph{cách đo cảm xúc} khác, chẳng hạn theo khía cạnh, cường độ, chủ thể hoặc mức bất
        ngờ so với kỳ vọng thị trường, thay cho hiệu hai xác suất;
  \item một \emph{cách tổng hợp} khác, chẳng hạn tổng hợp theo sự kiện thay vì theo ngày hoặc
        gán trọng số theo độ tin cậy của bộ phân loại;
  \item một chân trời dài hơn (xem Mục~\ref{sec:kq-horizon}).
\end{itemize}

Kết quả trong bảng chỉ cho thấy rằng \emph{không phát hiện được liên hệ đơn điệu giữa điểm số này và
lợi suất vượt trội của phiên kế tiếp}. Các phép đối chiếu trong bảng chưa được đăng ký
trước và gồm tám ô, nên toàn bộ nội dung của mục này mang tính \textbf{thăm dò}.

Hình mẫu ``cùng phiên có, phiên sau không'' \emph{phù hợp với} giả thuyết rằng tin đã
được phản ánh vào giá trước thời điểm đóng cửa, nhưng phù hợp không phải là chứng minh. Ít nhất
hai cơ chế khác có thể tạo ra cùng hình mẫu. Cơ chế thứ nhất là \textbf{nhân quả ngược}: giá biến động trước,
sau đó báo chí sử dụng giọng điệu phù hợp với diễn biến này, nên tương quan cùng phiên có thể phát sinh mà không cần
tin dẫn dắt giá. Cơ chế thứ hai là \textbf{nhiễu}: điểm số ban đầu đã yếu và việc bổ sung một ngày làm pha loãng
tín hiệu đủ để đưa nó xuống dưới ngưỡng phát hiện. Đồ án không có thiết kế để phân biệt ba cơ chế này,
do đó không đưa ra kết luận về tính hiệu quả của thị trường.
\subsection{Nhãn lợi suất vượt trội}
\label{sec:kq-excess}

Trong thiết lập mặc định của đồ án, nhãn là lợi suất phiên kế tiếp của chính mã cổ phiếu đó. Đối với mười mã VN30 trong cửa sổ nghiên cứu, một phần đáng kể biến động hằng ngày là nhịp chung: tương quan lợi suất ngày trung bình theo cặp là $0{,}418$ và hệ số xác định trung bình so với trung bình đồng hạng là $0{,}480$. Phần chung đó không phải thứ tin riêng của một doanh nghiệp giải thích được. Mục này chạy lại toàn bộ thang mô hình với một nhãn thay thế, được tính bằng lợi suất của từng mã trừ đi trung bình đồng hạng trong cùng ngày mục tiêu, qua đó phản ánh phần mã đó chạy nhanh hay chậm hơn rổ.

Thay đổi này chỉ liên quan đến \emph{nhãn}, không liên quan đến đặc trưng. Giá trị trung bình đồng hạng dùng để khấu trừ chỉ được quan sát vào đúng ngày mục tiêu, đồng thời với lợi suất, nên không có thông tin nào xuất hiện sớm hơn so với thiết lập ban đầu. Ngưỡng tam phân vẫn được khớp riêng trên phần huấn luyện của từng cửa sổ; toàn bộ cửa sổ, hạt giống và dòng kiểm thử được giữ nguyên.

\begin{table}[H]
  \centering
  \caption{F1 vĩ mô theo hai cách gán nhãn, cùng điểm kiểm đóng băng và cùng cửa sổ kiểm thử.}
  \label{tab:kq-excess}
  \begin{tabular}{l c c}
    \toprule
    Nhánh & Nhãn lợi suất thô & Nhãn lợi suất vượt trội \\
    \midrule
    Lớp phổ biến nhất & 0,1453 & 0,1627 \\
    Ngẫu nhiên theo phân bố lớp & 0,3289 & 0,3240 \\
    Hồi quy logistic, chỉ giá & 0,3551 & 0,3403 \\
    Hồi quy logistic, giá và cảm xúc & 0,3666 & 0,3396 \\
    LSTM chỉ giá & 0,3756 & 0,3399 \\
    LSTM giá và cảm xúc & 0,3883 & 0,3434 \\
    TFT chỉ giá & 0,3619 & 0,3379 \\
    TFT giá và cảm xúc & 0,3549 & 0,3263 \\
    \bottomrule
  \end{tabular}
\noindent\textit{Nguồn: \path{outputs/forecast_pit_sentiment} và \path{outputs/forecast_excess_pit_sentiment}.}
\end{table}

Kết quả cho thấy, với nhãn lợi suất vượt trội, mọi nhánh được huấn luyện đều đạt F1 vĩ mô \emph{thấp hơn} so với khi sử dụng nhãn lợi suất thô; chênh lệch giữa các nhánh cũng thu hẹp. Trên bộ dữ liệu này, bài toán xếp hạng tương đối trong phiên khó hơn bài toán dự báo chiều biến động thô.

Không nên diễn giải bảng này như một phép phân rã. Hiệu số F1 vĩ mô giữa hai cột phản ánh chênh lệch giữa \textbf{hai bài toán gán nhãn khác nhau, được huấn luyện riêng}, không phải phép phân rã phương sai hoặc phân rã nhân quả. Ngoài ra, khi tập kiểm thử lệch lớp, mức ngẫu nhiên của F1 vĩ mô không chính xác bằng $1/3$; do đó, không thể quy đổi hai cột về cùng một tỷ lệ phần trăm.

Đóng góp của thông tin phân cực đối với nhãn vượt trội, được đo bằng cùng đối chứng trung tính ghép cặp, là $-0{,}0005$ với khoảng theo ngày $[-0{,}0091; +0{,}0080]$ ở lượt đóng băng và $+0{,}0031$ với khoảng $[-0{,}0065; +0{,}0124]$ ở lượt hồi cứu. Cả hai khoảng đều chứa 0. Đối với TFT, lượt đóng băng cho kết quả $-0{,}0126$ với khoảng $[-0{,}0234; -0{,}0010]$. Đây là một ô trong nhiều phép so sánh được thực hiện trên cùng bộ dữ liệu, nên không thể diễn giải riêng kết quả này là có ý nghĩa thống kê.

Cuối cùng, mục này và Mục~\ref{sec:kq-ic} \textbf{sử dụng chung một bộ dữ liệu}. Hai mục không phải là các lần lặp lại độc lập và không củng cố lẫn nhau về mặt thống kê; chúng chỉ cung cấp hai góc nhìn khác nhau trên cùng một mẫu.

\subsection{Ablation chân trời dự báo}
\label{sec:kq-horizon}

Chân trời một phiên là trường hợp khó nhất: hàm ý của một bài tin phải biểu hiện trong đúng một nhịp từ giá đóng cửa này đến giá đóng cửa kế tiếp; nếu không thì không hiện ra ở đâu cả. Mục này kiểm tra liệu kết quả rỗng tại $h = 1$ có đặc thù cho nhịp thời gian này hay không bằng cách chạy lại toàn bộ thang mô hình và đối chứng trung tính ghép cặp tại $h = 3$ và $h = 5$. Điểm kiểm đóng băng, cửa sổ cuốn chiếu, hạt giống và mọi siêu tham số được giữ nguyên.

\begin{table}[H]
  \centering
  \caption{F1 vĩ mô và đóng góp thông tin phân cực theo chân trời dự báo, điểm kiểm đóng băng.}
  \label{tab:kq-horizon}
  \begin{tabular}{c c c c c c}
    \toprule
    Chân trời & Logistic chỉ giá & LSTM chỉ giá & LSTM giá--cảm xúc & Đóng góp & KTC 95\% theo ngày \\
    \midrule
    $h = 1$ & 0,3551 & 0,3756 & 0,3883 & $+0{,}0094$ & $[-0{,}0009; +0{,}0196]$ \\
    $h = 3$ & 0,3570 & 0,3668 & 0,3903 & $+0{,}0134$ & $[+0{,}0020; +0{,}0249]$ \\
    $h = 5$ & 0,3374 & 0,3575 & 0,3658 & $-0{,}0007$ & $[-0{,}0120; +0{,}0099]$ \\
    \bottomrule
  \end{tabular}
\noindent\textit{Nguồn: \path{outputs/forecast_pit_sentiment}, \path{outputs/forecast_h3_pit_sentiment}, \path{outputs/forecast_h5_pit_sentiment}.}
\end{table}

\paragraph{Cảnh báo nhãn chồng lấn.}
Khoảng tin cậy tại $h > 1$ \textbf{không} thể được diễn giải tương tự khoảng tại $h = 1$. Với chân trời $h$ phiên, hai dòng liên tiếp của cùng một mã sử dụng chung $h - 1$ phiên, nên các nhãn bị chồng lấn và số quan sát độc lập hữu hiệu chỉ còn khoảng $n/h$. Thủ tục lấy mẫu lặp vẫn xem mỗi ngày kiểm thử là một khối độc lập; vì vậy, tại $h = 3$ và $h = 5$, khoảng ước lượng \emph{hẹp hơn mức bằng chứng cho phép}. Do đó, khoảng $[+0{,}0020; +0{,}0249]$ tại $h = 3$ \textbf{không} được xem là bằng chứng về một hiệu ứng khác 0.

Sự chồng lấn không gây rò rỉ giữa tập huấn luyện và tập kiểm thử. Phép chia cuốn chiếu vẫn được thực hiện theo ngày mục tiêu, nên kết quả của một dòng huấn luyện luôn xảy ra trước mọi ngày mục tiêu trong tập kiểm thử. Sự chồng lấn ảnh hưởng đến phương sai, trong khi tính nhân quả không bị ảnh hưởng.

Với cảnh báo trên, bảng cho thấy đóng góp ước lượng lần lượt là $+0{,}0094$, $+0{,}0134$ và $-0{,}0007$ khi chân trời được mở rộng. Xu hướng này \textbf{không đơn điệu}, đồng thời biên độ dao động nằm trong cùng khoảng nhiễu đã quan sát ở các lượt khác. Mặt khác, mô hình chỉ giá suy giảm đều theo chân trời, từ 0,3756 xuống 0,3575, phù hợp với kỳ vọng khi mục tiêu dự báo xa dần. Kết quả này \textbf{không xác lập được sự phụ thuộc của đóng góp cảm xúc vào chân trời dự báo}; đồng thời, kết quả rỗng ở $h = 1$ không phải là hệ quả riêng của việc chọn một phiên.

\subsection{Chẩn đoán kinh tế}
\label{sec:kq-kinh-te}

\begin{tcolorbox}[colback=red!5!white, colframe=red!60!black, title=\textbf{Kết quả trong mục này không giao dịch được}]
Đặc trưng của ngày $t$ bao gồm giá đóng cửa $\text{close}_t$ và mọi bài tin công bố tới mốc cắt 15:00 của ngày $t$. Sổ lệnh mô phỏng trong mục này vào lệnh tại \emph{chính} $\text{close}_t$, tức mức giá đã được tín hiệu sử dụng và chỉ được biết sau khi phiên giao dịch kết thúc. Nhà đầu tư không thể khớp lệnh tại mức giá đó.

Vì vậy, mọi kết quả bên dưới là \textbf{cận trên theo giả định khớp lệnh hoàn hảo, tức thời và không có trượt giá}, không phải lợi nhuận đạt được. Một phương án có thể giao dịch phải vào lệnh từ phiên mở cửa kế tiếp trở đi và sẽ bỏ lỡ phần biến động qua đêm, vốn thường chiếm phần lớn biên lợi thế ở tần suất ngày. Mục này trả lời câu hỏi ``F1 vĩ mô $0{,}38$ tương ứng với mức lớn hay nhỏ tính bằng điểm cơ bản'', \textbf{không} trả lời câu hỏi ``chiến lược này có sinh lời hay không''.
\end{tcolorbox}

Thiết lập mô phỏng như sau: trong mỗi phiên, danh mục mua với trọng số đều các mã được dự báo \texttt{UP}, bán khống với trọng số đều các mã được dự báo \texttt{DOWN}, và tái cân bằng hằng ngày trong 300 phiên kiểm thử của lượt đóng băng. Lợi suất thực hiện được tính lại từ tệp giá thay vì suy ngược từ nhãn, nên một dòng bị gán nhãn sai được phản ánh thành một khoản lỗ. Chi phí được tính trên vòng quay thực tế so với các vị thế đã trôi giá, sau khi quy về cùng cơ sở vốn, với mức 20 điểm cơ bản cho mỗi đơn vị vòng quay.

Sổ lệnh này \textbf{không trung hòa thị trường}. Khi cả hai vế đều có lệnh, các trọng số triệt tiêu nhau; tuy nhiên, $23{,}2\%$ số phiên chỉ có một vế. Trong những phiên này, vế còn lại bị giảm một nửa nhưng vẫn duy trì rủi ro thị trường, với trạng thái ròng tuyệt đối trung bình là $0{,}101$.

\begin{table}[H]
  \centering
  \caption{Chẩn đoán kinh tế của các nhánh trên 300 phiên kiểm thử, điểm kiểm đóng băng, phí 20 điểm cơ bản.}
  \label{tab:kq-backtest}
  \begin{tabular}{l c c c c}
    \toprule
    Nhánh & Sharpe trước phí & Sharpe sau phí & Phí hòa vốn & $\alpha$ so với rổ ($t$) \\
    \midrule
    LSTM chỉ giá & $+2{,}49$ & $-1{,}05$ & 14,1 bps & $+0{,}00202$ $(+2{,}38)$ \\
    LSTM giá--cảm xúc & $+3{,}44$ & $-1{,}12$ & 15,1 bps & $+0{,}00252$ $(+3{,}39)$ \\
    TFT chỉ giá & $+1{,}93$ & $-2{,}01$ & 9,8 bps & $+0{,}00129$ $(+1{,}71)$ \\
    TFT giá--cảm xúc & $+3{,}13$ & $-0{,}79$ & 16,0 bps & $+0{,}00211$ $(+2{,}92)$ \\
    Ngẫu nhiên theo phân bố lớp & $-2{,}00$ & $-12{,}51$ & --- & $-0{,}00101$ $(-2{,}03)$ \\
    \midrule
    Rổ đều, tái cân bằng hằng ngày & $+1{,}95$ & --- & --- & --- \\
    Rổ đều, mua và nắm giữ & $+2{,}61$ & --- & --- & --- \\
    \bottomrule
  \end{tabular}
\noindent\textit{Nguồn: \path{outputs/forecast_pit_sentiment/backtest_*.json}. $\alpha$ là hệ số chặn khi hồi quy lợi suất trước phí của sổ lệnh lên lợi suất rổ đều.}
\end{table}

\paragraph{Mức nhiễu tham chiếu.}
Nhánh ngẫu nhiên đạt Sharpe $-2{,}00$, cho thấy 300 phiên \emph{không} đủ để phân biệt các mức Sharpe gần nhau. Mức tham chiếu phù hợp là phân phối null hoán vị, trong đó ngày, rổ mã và toàn bộ phương pháp hạch toán được giữ nguyên, còn nhãn dự báo được xáo trộn trong từng phiên qua 200 lần rút. Đối với nhánh LSTM giá--cảm xúc, phân phối null có Sharpe trung bình $+0{,}01$, độ lệch chuẩn $0{,}88$ và khoảng 95\% $[-1{,}77; +1{,}70]$. Kết quả này cho thấy hai điểm. Thứ nhất, phân phối null tập trung quanh 0, nên phương pháp hạch toán không có thiên lệch. Thứ hai, Sharpe trước phí $+3{,}44$ nằm ngoài khoảng này.

\paragraph{Ba kết quả.}
Thứ nhất, kết quả cho thấy biên lợi thế phân loại \emph{trước phí} tồn tại và không chỉ phản ánh trạng thái ròng một chiều. Sau khi hồi quy theo rổ, hệ số chặn vẫn dương với $t = +3{,}39$, trong khi $\beta = +0{,}186$. Thứ hai, chi phí \emph{triệt tiêu hoàn toàn biên lợi thế này}: ngưỡng hòa vốn vào khoảng 15 điểm cơ bản cho mỗi đơn vị vòng quay, trong khi tổng phí môi giới và thuế bán thực tế tại HOSE cao hơn đáng kể; do đó, Sharpe sau phí âm ở mọi nhánh. Thứ ba, khi so sánh với mốc thụ động, chiến lược mua và nắm giữ rổ đều đạt Sharpe $+2{,}61$. Vì vậy, ngay cả kết quả trước phí cũng không vượt xa một chiến lược không sử dụng mô hình.

Với nhãn lợi suất vượt trội tại Mục~\ref{sec:kq-excess}, Sharpe trước phí của nhánh LSTM giá--cảm xúc chỉ còn $+0{,}42$ và hệ số chặn so với rổ là $+0{,}00012$ với $t = +0{,}15$. Quan sát này nhất quán với việc thang mô hình đạt điểm thấp hơn khi sử dụng nhãn lợi suất vượt trội. Tuy nhiên, kết quả được tính trên cùng một mẫu dữ liệu nên không phải là một xác nhận độc lập.


\subsection{Giới hạn của kiểm tra theo điều kiện thị trường}
\label{sec:kq-strat}

Đối chứng ghép cặp mới đã trực tiếp kiểm soát sự hiện diện và khối lượng tin trên toàn bộ khóa dự đoán. Quy trình này không tạo ra một ước lượng ghép cặp được phân tầng riêng theo \texttt{has\_news}; vì vậy, kết quả phân tầng của artifact lịch sử không được sử dụng để đưa ra kết luận về những ngày có tin. Các phân tích phân tầng theo giai đoạn thị trường, ngày sự kiện hoặc mật độ tin chưa được thực hiện.

\subsection{Hạng mục chưa thực hiện}

Một hạng mục trong đề cương chưa được thực hiện trong phạm vi thời gian của đồ án và không được thay thế bằng kết quả dự kiến:

\begin{itemize}[nosep]
  \item \textbf{Biến thể PhoBERT kết hợp CNN hoặc BiLSTM}: chưa chạy, nên bảng cảm xúc chỉ
        có cấu hình tầng phân loại tuyến tính.
\end{itemize}

TFT đã được cấu hình cho bài toán phân loại ba lớp và chạy trên năm đoạn kiểm thử cuốn chiếu (Bảng~\ref{tab:ket-qua}). Kết quả của TFT thấp hơn LSTM hai nhánh nên không làm thay đổi kết luận của đồ án. Hạng mục còn lại được chuyển sang Chương~\ref{ch:huong-phat-trien}. Lý do là thứ tự ưu tiên: phạm vi tối thiểu gồm đối chứng cơ sở, mô hình chỉ giá, mô hình hai nhánh và phép tách đóng góp cảm xúc, qua đó trực tiếp trả lời câu hỏi trung tâm của đề tài. Tuy nhiên, trước khi thay thế LSTM bằng một kiến trúc phức tạp hơn, cần hoàn tất lượt đánh giá không rò rỉ thời điểm của bộ sinh cảm xúc; mô hình mới không thể khắc phục hạn chế của đầu vào hiện tại.

\section{Bàn luận}
\label{sec:kq-ban-luan}

\subsection{Bàn luận theo câu hỏi nghiên cứu}

\textbf{Câu hỏi thứ nhất}: mức độ chính xác của PhoBERT trên bộ dữ liệu tự xây dựng.
Kết quả xác thực chéo theo tầng trên 1.306 mẫu đạt F1 vĩ mô 0,7298, độ chính xác 0,8125 và độ
chính xác cân bằng 0,7689. Trong đề cương, mốc tham chiếu 85--90\% được xác định theo
độ chính xác; do đó, kết quả 0,8125 chưa đạt mốc này. Tuy nhiên, trên năm tập kiểm
thử ngoài, 59 mẫu NEGATIVE đạt độ bao phủ 0,6803, so với 0,0619 trên 31 mẫu trong lượt
lịch sử. Kết quả này cho thấy chất lượng phân loại của cấu hình có trọng số trên tập mở rộng
tốt hơn, nhưng không phải bằng chứng rằng tín hiệu cảm xúc đã cải thiện dự báo giá.

\textbf{Câu hỏi thứ hai}: xác suất cảm xúc có cải thiện kết quả dự báo hay không và mức cải
thiện là bao nhiêu. Trong lượt hồi cứu, so với đối chứng trung tính ghép cặp, đóng góp của
thông tin phân cực là $+0{,}0080$ F1 vĩ mô, với khoảng tin cậy theo ngày
$[-0{,}0025; +0{,}0186]$ chứa 0. Lượt ngoài mẫu sử dụng điểm kiểm đóng băng
(Mục~\ref{sec:kq-point-in-time}) cho kết quả $+0{,}0094$, với khoảng
$[-0{,}0009; +0{,}0196]$; đối chứng cùng kích thước cho kết quả $+0{,}0028$, với khoảng
$[-0{,}0074; +0{,}0133]$. Cả ba ước lượng điểm đều dương, nhưng tất cả khoảng tin cậy đều
chứa 0. Đồng thời, ba lượt đánh giá này thuộc số nhiều phép so sánh được thực hiện trên cùng
bộ dữ liệu, nên chưa có cơ sở để khẳng định ý nghĩa thống kê. Bốn phân tích bổ sung
(Mục~\ref{sec:kq-ic}, \ref{sec:kq-excess}, \ref{sec:kq-horizon}, \ref{sec:kq-kinh-te})
cũng không phát hiện cải thiện có khả năng tái lập. Kết luận:
\textbf{chưa có bằng chứng rằng đặc trưng cảm xúc cải thiện dự báo xu hướng phiên kế
tiếp}. Dữ liệu hiện tại không loại trừ khả năng có cải thiện nhỏ hoặc suy giảm nhỏ.

Hiệu số trực tiếp $+0{,}0113$ đồng thời phản ánh khác biệt về kiến trúc, sự hiện diện và
khối lượng tin, cũng như xác suất cảm xúc; trong đó, riêng thành phần kiến trúc và độ phủ
tin là $+0{,}0034$. Đối chứng trung tính ghép cặp giữ nguyên sự hiện diện và khối lượng
tin. Vì vậy, chỉ hiệu số $+0{,}0080$ giữa hai nhánh được dùng để đo lường đóng góp của xác
suất cảm xúc.

\textbf{Câu hỏi thứ ba}: các điều kiện mà trong đó cảm xúc phát huy tác dụng. Đề tài chưa
thực hiện ước lượng ghép cặp phân tầng riêng theo điều kiện thị trường, ngày sự kiện hoặc
mật độ tin. Vì vậy, đề tài chưa đưa ra kết luận về các điều kiện này. Các nghiên cứu về sự
kiện và quan hệ liên mã cho thấy điều kiện thị trường có thể làm thay đổi giá trị của tín
hiệu, nhưng các kết quả đó không thay thế việc kiểm tra trên rổ HOSE của đề tài
\cite{Chen2019_191005078v1,Xu2021_210207372v2,Siala2026_260200086v3}.

\subsection{Kiểm tra tính hợp lệ và độ ổn định}

Các nội dung kiểm tra sau đã được thực hiện và ghi nhận kết quả thực tế:

\begin{itemize}[nosep]
  \item \textbf{Quy mô đánh giá}: 5 cửa sổ, 3.000 dòng đánh giá cho mỗi hạt giống, khoảng
        300 ngày kiểm thử. Đủ cho khoảng tin cậy theo ngày. Tuy nhiên, do chỉ
        có 5 khối theo cửa sổ, khoảng tin cậy theo cửa sổ còn thô và không được sử dụng làm
        căn cứ duy nhất.
  \item \textbf{Độ ổn định dấu}: hiệu số đóng góp thông tin mang dấu dương ở bốn trên năm
        cửa sổ và mang dấu âm ở một cửa sổ. Dấu của hiệu số chưa đủ ổn định vì khoảng tin
        cậy theo ngày vẫn chứa 0.
  \item \textbf{Ngày không có tin}: các ngày này được gán phân phối tiên nghiệm trung tính,
        đồng thời cờ \texttt{has\_news} được đặt bằng 0, thay vì được gộp vào lớp trung
        tính. Quy tắc này đã được kiểm tra bằng kiểm thử tự động.
  \item \textbf{Ước lượng chỉ trên tập huấn luyện}: ngưỡng nhãn và cả hai bộ chuẩn hóa đều
        được ước lượng lại trong từng cửa sổ, chỉ sử dụng phần huấn luyện. Sự thay đổi của
        ngưỡng theo từng cửa sổ (Bảng~\ref{tab:kq-windows}) xác nhận quy trình này. Bất biến
        này cũng đã được kiểm tra bằng kiểm thử tự động.
  \item \textbf{Số tham số}: mô hình một nhánh có 5.091 tham số, trong khi mô hình hai
        nhánh có 11.203 tham số, tương đương gấp 2,20 lần và tăng thêm 6.112 tham số. Đối
        chứng phân phối tiên nghiệm trung tính được dùng để tách đóng góp của thông tin phân
        cực; chênh lệch $+0{,}0034$ F1 vĩ mô là phần đóng góp của kiến trúc và các đặc trưng
        về sự hiện diện, khối lượng tin.
  \item \textbf{Tính tái lập}: hai lần chạy độc lập cho kết quả trùng khớp đến từng chữ số
        đối với các mô hình chỉ sử dụng dữ liệu giá. Trong quá trình thực nghiệm, một lỗi
        ảnh hưởng đến tính tái lập đã được phát hiện và sửa: trọng số mạng được khởi tạo
        trước khi gieo hạt giống, khiến kết quả của một cửa sổ phụ thuộc vào thứ tự chạy của
        các mô hình trước đó.
  \item \textbf{Dữ liệu bị loại và mã ít tin}: 11 liên kết được đăng vào cuối ngày
        31/12/2025 rơi sang phiên năm 2026 nên bị loại; không có bản ghi nào thiếu mốc thời
        gian hợp lệ. Độ phủ tin có mức chênh lệch lớn giữa các mã, từ 548 liên kết (GAS) đến
        2.165 liên kết (MBB). Do kết quả theo từng mã chưa được phân tách, ảnh hưởng của mật
        độ tin chưa thể được kết luận.
  \item \textbf{Giới hạn thời điểm của bộ sinh}: thang dự báo chính sử dụng điểm kiểm được
        học từ toàn bộ 1.306 nhãn, bao gồm cả các nhãn sau những ngày kiểm thử trong giai
        đoạn 2024--2025, nên kết quả có tính hồi cứu. Lượt ngoài mẫu sử dụng điểm kiểm đóng
        băng trước ngày kiểm thử đầu tiên đã được đo lường và xác minh
        (\texttt{point\_in\_time: true}); kết quả được trình bày tại
        Mục~\ref{sec:kq-point-in-time}.
  \item \textbf{Phạm vi phân tầng}: đối chứng ghép cặp kiểm soát sự hiện diện và khối lượng
        tin trên toàn bộ mẫu. Phân tầng ghép cặp theo mã, độ biến động và ngày sự kiện
        \textbf{chưa} được thực hiện, nên câu hỏi thứ ba của đề cương chưa được trả lời.
  \item \textbf{Đặc trưng cảm xúc trượt}: các đặc trưng trung bình trượt và trung bình lũy
        giảm được tính theo nguyên tắc nhân quả. Kiểm thử đã xác nhận rằng một tin ở phiên
        sau không làm thay đổi giá trị của dòng trước đó.
\end{itemize}

Không phát hiện lỗi trong quá trình ánh xạ thời điểm tin và thực hiện các phép biến đổi của
nhánh giá. Tập giữ lại của nhánh giá không bị sửa đổi, và toàn bộ cấu hình được lưu trong
tệp kết quả. Các con số của thang chính vẫn mang tính hồi cứu; đánh giá ngoài mẫu được báo
cáo riêng tại Mục~\ref{sec:kq-point-in-time}.

\subsection{Giới hạn của thực nghiệm}
\label{sec:kq-limitations}

Thứ nhất, rổ nghiên cứu chỉ gồm 10 mã và dữ liệu chỉ thuộc HOSE, nên khả năng khái quát
kết quả cho HNX, UPCOM hoặc nhóm vốn hóa nhỏ chưa được kiểm tra. Thứ hai, nguồn tin chỉ
gồm Vietstock, trong khi đề cương dự kiến sử dụng bốn nguồn; vì vậy, độ phủ tin chênh lệch
lớn giữa các mã (MBB trên 2.000 liên kết, GAS 548). Thứ ba, khung dự báo phiên kế tiếp có
thể quá ngắn để phản ánh tác động dài hạn của tin.

Thứ tư, điểm kiểm cảm xúc của thang chính được tinh chỉnh trên toàn bộ 1.306 nhãn,
bao gồm cả các bài sau những ngày kiểm thử dự báo, nên các con số của thang chính mang
tính hồi cứu. Giới hạn này đã được khắc phục bằng một điểm kiểm đóng băng trước ngày kiểm
thử đầu tiên và một điểm kiểm đối chứng khớp kích thước (Mục~\ref{sec:kq-point-in-time});
cả hai đã được đo lường và xác minh. Tuy nhiên, lượt đối chứng mới chỉ được thực hiện với
một hạt giống, nên kết quả chỉ có tính thăm dò và chưa phải là ước lượng ổn định về ảnh
hưởng của kích thước tập.

Thứ năm, ngưỡng phân vị một phần ba tạo ra lớp FLAT theo nghĩa thống kê, không theo nghĩa
kinh tế. Ngưỡng khoảng $\pm0{,}5\%$ (Bảng~\ref{tab:kq-windows}) là một lựa chọn thiết kế,
không phải định nghĩa ``đi ngang'' được thị trường công nhận. Biến thể sử dụng dải cố định
chưa được kiểm tra.

Ngoài ra, điểm số phân loại không trực tiếp phản ánh giá trị kinh tế sau khi tính đến phí
giao dịch, trượt giá và giới hạn thanh khoản. Vì vậy, kết quả của chương này không phải
khuyến nghị đầu tư và không được xem là bằng chứng về một chiến lược giao dịch có lợi
nhuận.
\section{Tập nhãn mở rộng và giao thức đánh giá}
\label{sec:kq-mornhan}

Xác thực chéo phân tầng trên tập nhãn mở rộng đã được hoàn tất và cung cấp bằng chứng đánh giá cho bộ phân loại hiện tại. Mục này trình bày thiết kế của đợt rút mẫu 1.000 bài theo tầng, được thực hiện nhằm đưa lớp NEGATIVE lên quy mô đủ để đo lường.

Lấy mẫu ngẫu nhiên đều không khả thi cho mục tiêu này. Tỷ lệ NEGATIVE trong kho chỉ khoảng 10\%, do đó cần gán nhãn khoảng 6.000 bài để thu được 600 mẫu NEGATIVE. Vì vậy, đợt rút mẫu sử dụng ba tầng với các vai trò riêng biệt (Bảng~\ref{tab:kq-strata}).

\begin{table}[H]
  \centering
  \caption{Thiết kế ba tầng của đợt rút mẫu mở rộng.}
  \label{tab:kq-strata}
  \begin{tabular}{p{3.0cm} c c p{5.2cm}}
    \toprule
    Tầng & Số mẫu & Xác suất chọn & Vai trò \\
    \midrule
    \texttt{eval\_random} & 350 & 0,0377 & Duy trì tỷ lệ lớp của kho; đây là tầng duy nhất có thể dùng để đánh giá không chệch \\
    \texttt{train\_negative} & 400 & 0,2633 & Làm giàu lớp NEGATIVE; chỉ dùng để huấn luyện \\
    \texttt{train\_positive} & 250 & 0,0924 & Làm giàu lớp POSITIVE; chỉ dùng để huấn luyện \\
    \bottomrule
  \end{tabular}
\noindent\textit{Nguồn: \texttt{to\_label\_batch2\_manifest.json}.}
\end{table}

Hai tầng làm giàu sử dụng một từ điển phân cực tài chính tiếng Việt cố định để mở rộng tập hợp lệ. Sau đó, các mẫu được rút ngẫu nhiên với trọng số dựa trên số lần khớp từ khóa và có giới hạn về tỷ lệ của từng mã, nhằm tránh trường hợp một mã có mật độ tin cao chiếm toàn bộ tầng. Từ điển và xác suất của mô hình \textbf{chỉ quyết định bài nào được người đọc}, không được dùng để gán nhãn.

Việc lựa chọn xếp hạng theo từ điển thay vì theo xác suất của điểm kiểm dựa trên kết quả thực nghiệm. Trên toàn bộ 306 mẫu đã gán nhãn, từ điển làm tăng tỷ lệ NEGATIVE từ 0,101 lên 0,289, tương ứng với hệ số 2,85. Xác suất của điểm kiểm ban đầu cho hệ số cao hơn, nhưng kết quả này được tính trên chính các dòng đã dùng để huấn luyện điểm kiểm, do đó phản ánh khả năng ghi nhớ thay vì khả năng xếp hạng. Khi đánh giá lại trên 62 dòng kiểm thử ngoài của lượt 1, tức các dòng mà điểm kiểm chưa từng quan sát, \textbf{không thể xác nhận} ưu thế này. Chỉ có 16 dòng thỏa mãn điều kiện hợp lệ và 7 mẫu thuộc lớp NEGATIVE, không đủ để so sánh hai phương pháp xếp hạng. Vì vậy, xác suất của điểm kiểm không được sử dụng làm căn cứ xếp hạng. Từ điển là một quy tắc cố định chưa từng quan sát nhãn, nên hệ số 2,85 của nó là ước lượng trung thực và được chọn làm cơ sở thiết kế.

Toàn bộ 1.000 mẫu ban đầu được gán nhãn sơ bộ bằng mô hình ngôn ngữ lớn theo hướng dẫn gán nhãn của đề tài. Mỗi mẫu đi kèm mức độ tự tin và một câu lý giải. Sau quá trình rà soát, phân bố nhãn không thay đổi: tầng \texttt{train\_negative} có 30,2\% NEGATIVE, phù hợp với hệ số làm giàu được dự báo từ từ điển; tầng \texttt{eval\_random} có 8,0\%, gần với tỷ lệ 10,1\% của tập 306 mẫu trong lượt đầu. Mức độ phù hợp này cung cấp một kiểm tra về tính nhất quán của thiết kế lấy mẫu.

Người rà soát đã xác nhận toàn bộ 1.000 nhãn và không thay đổi bất kỳ nhãn sơ bộ nào. Các nhãn cuối cùng mang hai cờ \texttt{human\_reviewed} và \texttt{REVIEWED}, sau đó được hợp nhất thành tập gồm 1.306 mẫu, trong đó có 194 mẫu NEGATIVE. Xác thực chéo theo tầng trên tập này đã hoàn tất, với F1 vĩ mô trung bình đạt 0,7298. Kết quả được báo cáo tại Mục~\ref{sec:kq-sentiment-metrics}.

Thiết kế này có một đánh đổi cần được nêu rõ. Do hai tầng làm giàu làm thay đổi tỷ lệ lớp của kho, các tầng này chỉ được sử dụng để huấn luyện. Tập đánh giá gồm tầng \texttt{eval\_random} và 306 mẫu cũ, tương ứng với 656 mẫu, trong đó có 59 mẫu NEGATIVE. Vì vậy, nửa độ rộng khoảng tin cậy của độ bao phủ lớp NEGATIVE giảm từ $\pm0{,}176$ xuống $\pm0{,}128$. Đây là một mức cải thiện thực tế, nhưng \textbf{không} đạt mức $\pm0{,}07$ như khi tính trên toàn bộ 194 mẫu. Để đạt mức này mà vẫn duy trì tỷ lệ lớp không lệch, cần gán nhãn khoảng 1.500 mẫu ngẫu nhiên, vượt quá ngân sách của đồ án.

Đây là lựa chọn có chủ ý: khoảng tin cậy rộng hơn nhưng không chệch được ưu tiên hơn khoảng tin cậy hẹp được tính trên tập mẫu đã chọn để làm giàu lớp thiểu số. Độ bao phủ của từng lớp trên các tầng làm giàu vẫn có thể được báo cáo, nhưng phải được xác định rõ là số liệu chẩn đoán trên tập mẫu đã làm giàu, không phải là ước lượng cho toàn bộ kho.
\section{Tái lập kết quả}
\label{sec:kq-taplap}

Bảng~\ref{tab:kq-artifact} liệt kê các lệnh và tệp kết quả tương ứng với từng bảng số liệu trong chương.

\begin{table}[H]
  \centering
  \caption{Nguồn artifact của các bảng kết quả.}
  \label{tab:kq-artifact}
  \begin{tabular}{p{3.4cm} p{9.4cm}}
    \toprule
    Bảng & Lệnh và tệp kết quả \\
    \midrule
    \ref{tab:kq-sentiment}, \ref{tab:kq-sentiment-class} & \texttt{stf.cli sentiment-cv}; \path{outputs/sentiment-cv-merged/cv_results.json} và tệp nén của lần chạy có trọng số \\
    \ref{tab:kq-windows}, \ref{tab:kq-price-stats}, \ref{tab:ket-qua} & \texttt{stf.cli forecast --news-sentiment ...}; \path{outputs/forecast_merged_sentiment_tft/forecast_metrics.csv} \\
    \ref{tab:kq-ablation}, \ref{tab:kq-ablation-ci} & \texttt{stf.cli forecast --neutral-news-sentiment ...}, \texttt{stf.cli forecast-compare}; \path{outputs/forecast_merged_sentiment_tft/information_gain.json} \\
    \ref{tab:kq-pit} & \texttt{stf.cli forecast} + \texttt{forecast-compare} trên \texttt{news\_sentiment\_pit.parquet} và \texttt{news\_sentiment\_sizematched\_s7.parquet}; \path{outputs/forecast_pit_sentiment/information_gain.json}, \path{outputs/forecast_sizematched_s7_sentiment/information_gain.json} \\
    \ref{tab:kq-strata} & \texttt{stf.cli label-candidates}; \texttt{to\_label\_batch2\_manifest.json} \\
    \bottomrule
  \end{tabular}
\noindent\textit{Nguồn: Tác giả tổng hợp.}
\end{table}

Kết quả của lần chạy đối chứng sử dụng phân phối tiên nghiệm trung tính được lưu riêng tại
\texttt{outputs/forecast\_merged\_control/}. Lần chạy này sử dụng cùng tệp xác suất đã được chấm như lần
chạy chính, nhưng tham số \texttt{-{}-neutral-news-sentiment} thay thế ba xác suất bằng phân phối tiên nghiệm
trung tính. Lệnh \texttt{forecast-compare} từ chối các cặp khác nhau về cấu hình, mã băm giá,
mã băm tệp cảm xúc, số dòng hoặc khóa dự đoán. Do đó, phép tách tại
Mục~\ref{sec:kq-donggop} không thể trộn lẫn các đầu vào khác nhau.

\chapter{KẾT LUẬN}
\label{ch:ket-luan}

Chương này tổng hợp những nội dung đã thực hiện, đối chiếu kết quả với mục tiêu và câu
hỏi nghiên cứu, sau đó nêu đóng góp, giới hạn và kết luận cuối cùng của đồ án. Các nhận
định chỉ dựa trên số liệu đã trình bày tại Chương~\ref{ch:ket-qua}. Kết quả hồi cứu, kết
quả ngoài mẫu, phân tích thăm dò và nội dung chưa thực hiện được phân biệt rõ để tránh
suy rộng vượt quá bằng chứng.

\section{Tóm tắt nội dung thực hiện}

Đồ án đã xây dựng một quy trình dự báo xu hướng biến động giá cổ phiếu trên sàn HOSE
bằng cách kết hợp dữ liệu giá với cảm xúc của tin tức tài chính tiếng Việt. Nhánh văn
bản sử dụng PhoBERT để tạo xác suất của ba lớp NEGATIVE, NEUTRAL và POSITIVE. Các xác
suất được tổng hợp theo mã cổ phiếu và phiên giao dịch cùng số lượng tin, độ phân tán
và cờ thể hiện sự hiện diện của tin. Nhánh giá sử dụng dữ liệu OHLCV đã điều chỉnh,
lợi suất và các đặc trưng kỹ thuật được tính theo nguyên tắc nhân quả.

Hai nguồn dữ liệu được căn chỉnh theo mã, ngày giao dịch và mốc cắt 15 giờ. Tin sau mốc
cắt, vào cuối tuần hoặc ngày nghỉ được chuyển sang phiên hợp lệ kế tiếp. Dữ liệu giá
được đánh giá bằng năm cửa sổ cuốn chiếu không chồng lấn; bộ chuẩn hóa, bộ điền khuyết
và ngưỡng tạo nhãn chỉ được ước lượng từ phần huấn luyện của từng cửa sổ. Thang mô hình
gồm đối chứng lớp phổ biến, đối chứng ngẫu nhiên, hồi quy logistic, LSTM và Bộ biến đổi
hợp nhất theo thời gian. Đối với nhánh có tin, một đối chứng giữ nguyên kiến trúc, bài
tin, thời điểm và số lượng tin nhưng thay xác suất cảm xúc bằng phân phối tiên nghiệm
trung tính được sử dụng để tách phần đóng góp của thông tin phân cực.

Kho dữ liệu trong cửa sổ 2020--2025 gồm 10.695 bài tin, 12.624 liên kết tin--mã và
14.990 dòng mã--phiên của 10 mã thuộc VN30. Tập nhãn cảm xúc trong miền dữ liệu gồm
1.306 mẫu đã được rà soát. Sau khi khóa cấu hình, điểm kiểm triển khai được tinh chỉnh
lại trên toàn bộ tập nhãn và dùng để chấm 14.256 liên kết tin--mã. Bên cạnh lượt hồi
cứu này, đồ án đã thực hiện một lượt ngoài mẫu với điểm kiểm chỉ học từ 1.049 nhãn
công bố trước ngày kiểm thử đầu tiên.

\section{Đối chiếu với mục tiêu và câu hỏi nghiên cứu}

\subsection{Mức độ hoàn thành các mục tiêu}

Mục tiêu thứ nhất được hoàn thành một phần. Quy trình thu thập, làm sạch, loại trùng,
ánh xạ theo phiên và gán nhãn cảm xúc đã được xây dựng; PhoBERT với tầng phân loại
tuyến tính đã được đánh giá bằng xác thực chéo phân tầng 5 lượt. Tuy nhiên, biến thể
PhoBERT kết hợp CNN hoặc BiLSTM chưa được thực nghiệm, nên phần so sánh giữa các tầng
phân loại chưa hoàn thành.

Mục tiêu thứ hai đã được thực hiện đối với toàn bộ thang mô hình đã khóa. Đặc trưng cảm
xúc được tạo từ phân phối xác suất, sau đó được đưa vào hồi quy logistic, LSTM hai
nhánh và TFT. Các cấu hình có cảm xúc được đặt cạnh đối chứng chỉ giá và đối chứng
trung tính trên cùng ngày huấn luyện, xác thực và kiểm thử.

Mục tiêu thứ ba được hoàn thành ở phần đánh giá cuốn chiếu và kiểm soát các phép biến
đổi của nhánh giá. Đồ án cũng đã kiểm tra riêng giới hạn thời điểm của bộ sinh cảm xúc
bằng điểm kiểm đóng băng. Tuy nhiên, phân tích ghép cặp theo giai đoạn thị trường, mật
độ tin và ngày sự kiện chưa được thực hiện; do đó, phần xác định điều kiện phát huy tín
hiệu vẫn còn thiếu.

Mục tiêu thứ tư được hoàn thành một phần. Các lần chạy lưu cấu hình, mã băm dữ liệu,
điểm kiểm, phần chia và tệp dự đoán; hệ thống dòng lệnh và dashboard có thể đọc các
artifact này để trình bày kết quả. Phép loại bỏ đặc trưng và đối chứng trung tính hỗ
trợ nhận diện giới hạn của tín hiệu cảm xúc. Tuy nhiên, tầm quan trọng của biến, cơ chế
chú ý và các phương pháp giải thích ở cấp bài tin chưa được đánh giá đầy đủ.

\subsection{Trả lời các câu hỏi nghiên cứu}

Bảng~\ref{tab:ket-luan-cau-hoi} tóm tắt câu trả lời cho bốn câu hỏi nghiên cứu. Trạng
thái ``một phần'' được sử dụng khi đồ án đã có bằng chứng cho một vế của câu hỏi nhưng
chưa thực hiện đầy đủ phép so sánh hoặc phân tích đã đặt ra.

\begin{table}[H]
  \centering
  \small
  \caption{Đối chiếu câu hỏi nghiên cứu với bằng chứng thực nghiệm.}
  \label{tab:ket-luan-cau-hoi}
  \begin{tabular}{p{1.0cm} p{4.1cm} p{4.8cm} p{2.0cm}}
    \toprule
    Câu hỏi & Nội dung cần trả lời & Bằng chứng chính & Trạng thái \\
    \midrule
    1 & Khả năng phân loại cảm xúc của PhoBERT và khác biệt giữa các tầng phân loại &
    F1 vĩ mô 0,7298; độ chính xác 0,8125; độ chính xác cân bằng 0,7689. Chưa chạy
    biến thể CNN hoặc BiLSTM. & Một phần \\
    2 & Giá trị gia tăng của cảm xúc đối với dự báo phiên kế tiếp &
    Hiệu số ngoài mẫu: $+0{,}0094$. KTC 95\% theo ngày: từ $-0{,}0009$ đến
    $+0{,}0196$; khoảng chứa 0. & Đã trả lời \\
    3 & Điều kiện thị trường hoặc mật độ tin làm thay đổi hiệu quả &
    Chưa có ước lượng ghép cặp phân tầng theo giai đoạn, mã, ngày sự kiện hoặc mật độ
    tin. & Chưa trả lời \\
    4 & Vai trò và giới hạn của đặc trưng cảm xúc trong quyết định dự báo &
    Đối chứng trung tính, loại bỏ đặc trưng và tương quan hạng cho thấy tín hiệu yếu;
    chưa có phân tích chú ý hoặc giải thích ở cấp bài. & Một phần \\
    \bottomrule
  \end{tabular}
\noindent\textit{Nguồn: Tác giả tổng hợp từ các kết quả tại Chương~\ref{ch:ket-qua}.}
\end{table}

Đối với câu hỏi thứ nhất, cấu hình PhoBERT có trọng số lớp trên 1.306 mẫu đạt F1 vĩ
mô 0,7298, độ chính xác 0,8125 và độ chính xác cân bằng 0,7689. Lớp NEGATIVE đạt F1
0,6036 và độ bao phủ 0,6803 trên tổng cộng 59 mẫu kiểm thử ngoài. Kết quả cho thấy mô
hình đã phân biệt được cả ba lớp và không còn suy biến về lớp NEUTRAL như cấu hình
lịch sử. Tuy nhiên, số mẫu NEGATIVE còn ít và không có phép so sánh với tầng CNN hoặc
BiLSTM; vì vậy, chưa thể kết luận kiến trúc tầng phân loại nào phù hợp hơn.

Đối với câu hỏi thứ hai, LSTM chỉ sử dụng giá đạt F1 vĩ mô 0,3756, cao hơn mức 0,3289
của đối chứng ngẫu nhiên và 0,1453 của đối chứng lớp phổ biến. TFT chỉ sử dụng giá đạt
0,3619, thấp hơn LSTM. Trong lượt hồi cứu, đóng góp riêng của thông tin phân cực là
$+0{,}0080$ F1 vĩ mô với khoảng 95\% theo ngày
$[-0{,}0025; +0{,}0186]$. Lượt ngoài mẫu sử dụng điểm kiểm đóng băng đạt
$+0{,}0094$, với khoảng $[-0{,}0009; +0{,}0196]$; điểm kiểm đối chứng cùng kích thước
đạt $+0{,}0028$, với khoảng $[-0{,}0074; +0{,}0133]$. Cả ba ước lượng điểm đều dương
nhưng mọi khoảng theo ngày đều chứa 0. Đối với TFT, đóng góp thông tin không dương ở
cả ba bộ sinh. Vì vậy, kết luận phù hợp là chưa có bằng chứng đủ mạnh cho thấy cảm xúc
cải thiện dự báo, đồng thời dữ liệu hiện tại cũng không loại trừ khả năng tồn tại một
mức cải thiện nhỏ.

Đối với câu hỏi thứ ba, đối chứng ghép cặp đã kiểm soát sự hiện diện và khối lượng tin
trên toàn bộ mẫu, nhưng chưa được tính riêng theo điều kiện thị trường, mã, ngày sự kiện
hoặc mật độ tin. Ablation theo chân trời cho các hiệu số $+0{,}0094$, $+0{,}0134$ và
$-0{,}0007$ tại một, ba và năm phiên, nhưng nhãn tại chân trời dài bị chồng lấn nên
khoảng tin cậy chưa phản ánh đúng số quan sát độc lập hữu hiệu. Kết quả này không cho
phép xác định điều kiện mà tín hiệu cảm xúc phát huy tác dụng.

Đối với câu hỏi thứ tư, phép đối chứng trung tính cho phép tách thông tin phân cực khỏi
ảnh hưởng của kiến trúc, sự hiện diện và khối lượng tin. Tương quan hạng giữa điểm
\texttt{sent\_pos\_minus\_neg} với lợi suất phiên kế tiếp gần 0, cho thấy chưa phát
hiện quan hệ đơn điệu của đúng điểm số này. Tuy nhiên, kết quả không loại trừ quan hệ
phi tuyến, quan hệ có điều kiện hoặc một cách biểu diễn cảm xúc khác. Đồ án chưa báo
cáo tầm quan trọng của biến hoặc chú ý của TFT, nên câu hỏi về cơ chế bên trong mô hình
mới chỉ được trả lời thông qua phép loại bỏ và đối chứng, không phải qua giải thích ở
cấp biến hoặc cấp bài tin.

\section{Đóng góp và ý nghĩa của đồ án}

Đóng góp thứ nhất là một quy trình dữ liệu có khả năng truy nguyên, từ bài tin và giá
gốc đến bảng mã--phiên dùng để huấn luyện. Quy trình lưu URL, nguồn, mốc thời gian, mã
băm và trạng thái ánh xạ; đồng thời phân biệt ngày không có tin với ngày có tin trung
tính. Cách tổ chức này giúp kiểm tra lại việc một bài viết có thực sự quan sát được tại
thời điểm dự báo hay không.

Đóng góp thứ hai là cách sử dụng phân phối xác suất cảm xúc thay cho một nhãn cứng.
Việc giữ cả ba xác suất cho phép tổng hợp mức độ không chắc chắn và xây dựng đối chứng
trung tính mà không thay đổi số bài, thời điểm hoặc cờ có tin. Nhờ đó, hiệu số giữa mô
hình chỉ giá và mô hình hai nhánh được phân tách thành ảnh hưởng của kiến trúc--độ phủ
tin và phần thông tin phân cực.

Đóng góp thứ ba là giao thức đánh giá theo thời gian với năm đoạn kiểm thử tách biệt,
tiền xử lý chỉ học từ phần huấn luyện và cơ chế xác minh điểm kiểm cảm xúc theo mốc thời
gian. Việc báo cáo song song lượt hồi cứu và lượt ngoài mẫu giúp chỉ ra rằng một điểm số
dương chưa đủ để kết luận khi khoảng tin cậy chứa 0 hoặc khi bộ sinh tín hiệu đã tiếp
cận nhãn tương lai.

Đóng góp cuối cùng là một kết quả thực nghiệm thận trọng. Trên dữ liệu đã đo, cảm xúc
không tạo ra mức cải thiện có thể tái lập. Kết quả này vẫn có ý nghĩa vì nó được thu
nhận bằng phép đối chứng ghép cặp và được kiểm tra qua nhiều góc nhìn, thay vì giả định
trước rằng mô hình ngôn ngữ hoặc cơ chế hợp nhất phức tạp hơn sẽ cho kết quả tốt hơn.
Các hướng dùng biểu diễn ngữ cảnh, sự kiện chi tiết hoặc mô hình ngôn ngữ lớn có khả
năng giữ nhiều thông tin hơn, nhưng cũng yêu cầu dữ liệu và giao thức kiểm chứng tương
ứng \cite{Chen2021_210708721v1,Elahi2024_241101368v1,Liang2024_241210823v2}.

Về thực tiễn, hệ thống có thể hỗ trợ thăm dò dữ liệu, theo dõi xác suất cảm xúc, trạng
thái kỹ thuật và lớp dự báo theo ngày. Tuy nhiên, hệ thống chưa cấu thành một chiến
lược giao dịch. Chẩn đoán kinh tế trong Chương~\ref{ch:ket-qua} sử dụng giá khớp không
khả thi và chỉ cho thấy biên lợi thế trước phí bị triệt tiêu khi áp dụng mức chi phí đã
giả định. Vì vậy, mọi đầu ra phải được hiểu là kết quả nghiên cứu, không phải khuyến
nghị mua hoặc bán.

\section{Giới hạn của đồ án}

Thứ nhất, phạm vi nghiên cứu chỉ gồm 10 mã thuộc VN30 trên HOSE, một nguồn tin Vietstock
và giai đoạn 2020--2025. Độ phủ tin chênh lệch đáng kể giữa các mã. Kết quả vì vậy
không đại diện cho toàn bộ VN30, nhóm vốn hóa nhỏ, HNX, UPCOM hoặc thị trường khác.

Thứ hai, tập cảm xúc chỉ có một người rà soát nên không có Cohen's kappa hoặc Fleiss'
kappa. Năm tập kiểm thử ngoài chỉ có tổng cộng 59 mẫu NEGATIVE. Phép đánh giá đã đo
được lớp này, nhưng khoảng bất định vẫn lớn hơn hai lớp còn lại. Ngoài ra, mô hình ngôn
ngữ lớn được dùng để tạo nhãn sơ bộ trước khi người rà soát xác nhận; thiết kế này không
thay thế được một quy trình gán nhãn độc lập bởi nhiều người.

Thứ ba, điểm kiểm đóng băng loại bỏ các nhãn được công bố sau ngày kiểm thử đầu tiên,
nhưng cấu hình đầu vào, cách cắt token và trọng số lớp vẫn được lựa chọn từ xác thực
chéo trên toàn bộ 1.306 nhãn. Vì vậy, lượt này kiểm soát rò rỉ nhãn nhưng chưa tạo
thành một quy trình lựa chọn mô hình hoàn toàn point-in-time. Đối chứng cùng kích thước
chỉ sử dụng một hạt giống, nên chưa tách ổn định ảnh hưởng của thời điểm huấn luyện khỏi
ảnh hưởng của thành phần mẫu được rút.

Thứ tư, khoảng tin cậy theo ngày chưa mô hình hóa đầy đủ phụ thuộc giữa các phiên liên
tiếp; khoảng theo cửa sổ chỉ dựa trên năm khối nên còn thô. Nhiều phép so sánh được
thực hiện trên cùng bộ dữ liệu, gồm ba điểm kiểm, hai kiến trúc và bốn phân tích bổ sung.
Do đó, một khoảng riêng lẻ có cận gần 0 không đủ để tuyên bố ý nghĩa thống kê.


Thứ năm, câu hỏi về điều kiện phát huy tín hiệu chưa được trả lời bằng ước lượng ghép
cặp phân tầng. Kết quả hiện tại chưa cho biết cảm xúc có hữu ích hơn trong giai đoạn
biến động cao, ngày nhiều tin, ngày sự kiện hoặc đối với một nhóm mã cụ thể hay không.
Các nghiên cứu về sự kiện và quan hệ liên mã cho thấy đây là những hướng có cơ sở, nhưng
kết quả từ thị trường khác không thay thế được kiểm tra trên dữ liệu của đề tài
\cite{Chen2019_191005078v1,Xu2021_210207372v2}.

Thứ sáu, nhãn tam phân được xác định theo phân vị của phần huấn luyện nên lớp FLAT mang
ý nghĩa thống kê, không phải một ngưỡng đi ngang được thị trường công nhận. Chân trời
một phiên có thể quá ngắn đối với thông tin dài hạn; ngược lại, các chân trời dài hơn
tạo nhãn chồng lấn và cần phương pháp ước lượng độ bất định khác.

Cuối cùng, chẩn đoán kinh tế chưa sử dụng mức giá có thể giao dịch và sổ theo dõi trực
tuyến chưa có dòng dự báo hợp lệ được phát hành trước phiên mục tiêu. Hệ thống mới chứng
minh đường xử lý có thể chạy từ đầu đến cuối và phát hiện lần chạy lại; chưa có bằng
chứng về hiệu năng vận hành thực tế.

\section{Kết luận cuối cùng}

Trên 10 mã HOSE trong giai đoạn 2020--2025, mô hình LSTM chỉ sử dụng giá đạt kết quả
cao hơn hai đối chứng cơ sở về F1 vĩ mô. TFT tối giản không vượt LSTM. Khi giữ nguyên
kiến trúc, bài tin, thời điểm và số lượng tin, xác suất cảm xúc tạo ra các ước lượng
điểm nhỏ và dương đối với LSTM, nhưng các khoảng tin cậy theo ngày đều chứa 0. Các phân
tích tương quan hạng, nhãn lợi suất vượt trội, chân trời dự báo và chẩn đoán kinh tế
không cung cấp bằng chứng độc lập để thay đổi kết luận này.

Vì vậy, trong phạm vi dữ liệu và giao thức đã đo, \textbf{chưa có bằng chứng đủ mạnh
cho thấy cảm xúc trong tin tài chính tiếng Việt cải thiện dự báo xu hướng của phiên
giao dịch kế tiếp}. Kết luận này không có nghĩa là cảm xúc không hữu ích trong mọi bối
cảnh; dữ liệu hiện tại vẫn cho phép tồn tại một hiệu ứng nhỏ, một quan hệ phi tuyến
hoặc một hiệu ứng chỉ xuất hiện trong một số điều kiện chưa được kiểm tra. Kết quả không
chứng minh quan hệ nhân quả và không phải khuyến nghị đầu tư.

\chapter{HƯỚNG PHÁT TRIỂN}
\label{ch:huong-phat-trien}

Các hướng phát triển trong chương này xuất phát trực tiếp từ những phần chưa hoàn thành
và các giới hạn đã nêu ở Chương~\ref{ch:ket-luan}. Thứ tự ưu tiên đặt độ tin cậy của dữ
liệu và giao thức trước độ phức tạp của mô hình. Một kiến trúc mới chỉ có ý nghĩa khi
được so sánh trên cùng phần chia, cùng ngân sách thử nghiệm và cùng điều kiện point-in-time.

\section{Hoàn thiện các mục tiêu còn thiếu}

Hạng mục đầu tiên là so sánh tầng phân loại tuyến tính của PhoBERT với ít nhất một biến
thể PhoBERT--CNN hoặc PhoBERT--BiLSTM trên cùng tập dữ liệu. Hai cấu hình phải sử dụng
cùng đầu vào, cách cắt token, phần chia và tiêu chí chọn điểm kiểm. Kết quả cần báo cáo
F1 vĩ mô, chỉ số theo lớp, số tham số, thời gian huấn luyện và độ nhạy theo hạt giống.
Nếu kiến trúc kết hợp chỉ cải thiện trên tập xác thực nhưng không cải thiện trên các tập
kiểm thử ngoài, cấu hình tuyến tính vẫn nên được giữ làm mức sàn.

Hạng mục thứ hai là thực hiện ước lượng ghép cặp theo điều kiện thị trường. Các nhóm có
thể gồm giai đoạn biến động cao và thấp, ngày nhiều tin và ít tin, ngày sự kiện và ngày
thông thường, hoặc từng nhóm mã có độ phủ tương đương. Tiêu chí phân tầng phải được
xác định trước khi quan sát hiệu số. Mỗi nhóm cần báo cáo số ngày, số dòng mã--phiên,
phân bố lớp và khoảng tin cậy; nhóm quá nhỏ phải được ghi là không đủ cơ sở kết luận.

Hạng mục thứ ba là bổ sung phương pháp giải thích phù hợp với từng kiến trúc. Đối với
LSTM, có thể che từng nhóm đặc trưng hoặc hoán vị theo khối thời gian. Đối với TFT, có
thể báo cáo tầm quan trọng của biến và trọng số chú ý trên các cửa sổ ngoài mẫu. Các
giá trị này chỉ mô tả cách mô hình sử dụng đầu vào, không chứng minh cảm xúc gây ra
biến động giá.

\section{Củng cố đánh giá point-in-time}

Ưu tiên phương pháp quan trọng nhất là đưa cả bước lựa chọn cấu hình cảm xúc vào giao
thức point-in-time. Thay vì lựa chọn \texttt{title\_context},
\texttt{head\_tail} và trọng số lớp trên toàn bộ 1.306 nhãn rồi mới loại nhãn tương
lai, mỗi cửa sổ dự báo cần sử dụng một quy trình lồng. Phần nhãn công bố trước cửa sổ
được chia thành huấn luyện và xác thực để chọn cấu hình; điểm kiểm sau đó được tinh
chỉnh lại chỉ từ dữ liệu quá khứ và dùng để chấm các bài tin của cửa sổ tương ứng.

Đối chứng cùng kích thước cần được lặp lại với nhiều hạt giống. Mỗi lần rút phải giữ
đúng số mẫu của từng lớp, lưu danh sách dòng và mã băm, sau đó chạy lại toàn bộ phép
đối chứng trung tính. Phân phối của hiệu số qua các lần rút mới cho phép tách ảnh hưởng
của thời điểm huấn luyện khỏi ảnh hưởng ngẫu nhiên của thành phần mẫu. Nếu chi phí tinh
chỉnh theo từng cửa sổ quá lớn, có thể sử dụng một điểm kiểm được khóa trước toàn bộ
giai đoạn kiểm thử, nhưng mọi quyết định về đầu vào và siêu tham số cũng phải được đưa
ra trước mốc này.

Khoảng tin cậy cần được điều chỉnh cho phụ thuộc chuỗi thời gian. Có thể sử dụng lấy mẫu
lặp theo khối ngày liên tiếp hoặc lấy mẫu lặp theo cửa sổ khi số cửa sổ được tăng lên.
Đối với chân trời nhiều phiên, độ dài khối phải phản ánh mức chồng lấn của nhãn. Các
phép kiểm tra cần được đăng ký trước hoặc ghi đầy đủ số lần thử để hạn chế việc lựa chọn
một kết quả thuận lợi từ nhiều phép so sánh.

\section{Mở rộng dữ liệu và chất lượng nhãn}

Nguồn tin có thể được mở rộng sang CafeF, VnEconomy hoặc Stockbiz. Mỗi nguồn mới cần
được kiểm tra về điều kiện truy cập, độ phủ theo mã, mốc thời gian, nội dung, bản đăng
lại và khả năng tái lập. Các bài cùng xuất phát từ một thông cáo phải được nhận diện
trước khi chia tập. Bản ghi thiếu thời gian hoặc không xác định được mã không nên được
đưa vào bảng đặc trưng chính bằng một giá trị suy đoán.

Tập đánh giá cảm xúc cần được mở rộng bằng lấy mẫu ngẫu nhiên không chệch. Theo phân bố
đã quan sát, khoảng 1.500 mẫu ngẫu nhiên là một mục tiêu phù hợp để thu hẹp khoảng tin
cậy của độ bao phủ lớp NEGATIVE. Ít nhất hai người nên gán một phần giao nhau của tập
mẫu để tính độ đồng thuận và phân tích các trường hợp bất đồng. Nhãn có thể được mở
rộng thêm chủ thể, loại sự kiện, cường độ và mức độ chắc chắn. Biểu diễn sự kiện chi
tiết có thể bổ sung cho ba lớp phân cực, nhưng cần một hướng dẫn gán nhãn và phép đánh
giá riêng \cite{Chen2019_191005078v1}.

Độ phủ giữa các mã cũng cần được cân bằng hơn. Khi bổ sung mã hoặc nguồn, báo cáo phải
nêu số bài, tỷ lệ ngày có tin, số bản ghi bị loại và phân bố nhãn theo từng mã. Không
nên gộp toàn bộ thành một điểm số nếu một số mã có mật độ tin cao hơn nhiều so với các
mã còn lại.

\section{Mở rộng biểu diễn và cơ chế hợp nhất}

Trước khi thay đổi mô hình, cần thực hiện loại bỏ đầu vào của nhánh văn bản. Các cấu
hình nên bao gồm chỉ tiêu đề, tiêu đề với phần đầu, tiêu đề với phần cuối và kết hợp
đầu--cuối trong cùng giới hạn 256 token. So sánh này giúp xác định liệu ngữ cảnh dài
hơn có bổ sung thông tin hay chỉ đưa thêm quảng cáo, nội dung lặp và sự kiện không liên
quan. Mọi cấu hình phải dùng cùng các tập nhãn và cùng quy trình point-in-time.

Khi dữ liệu đủ lớn, phép nối hai trạng thái cuối có thể được thay bằng cơ chế hợp nhất
có cổng hoặc chú ý chéo. Cổng có thể điều chỉnh đóng góp của nhánh cảm xúc theo độ biến
động, số tin và độ chắc chắn của bộ phân loại. Chú ý chéo cho phép chuỗi giá truy vấn
chuỗi tin và ngược lại. Các mô hình đa phương thức cho thấy hướng này có khả năng giữ
được tương tác phong phú hơn phép nối, nhưng số tham số và phương sai cũng tăng
\cite{Omranpour2024_241210540v1,Elahi2024_241101368v1}. Vì vậy, mỗi biến thể phải được
đặt cạnh LSTM hai nhánh trên cùng cửa sổ, hạt giống, đối chứng trung tính và khoảng tin
cậy.

Một hướng khác là biểu diễn cảm xúc theo khía cạnh hoặc theo sự kiện thay vì một hiệu
xác suất duy nhất. Có thể tách thông tin về kết quả kinh doanh, cổ tức, quản trị, pháp
lý và đầu tư, sau đó tổng hợp theo loại sự kiện. Mô hình ngôn ngữ lớn có thể hỗ trợ tạo
nhãn sơ bộ hoặc biểu diễn ngữ cảnh, nhưng phiên bản, câu lệnh, chi phí và quy trình rà
soát phải được ghi lại. Hệ thống này không nên thay thế đối chứng PhoBERT mở khi chưa
chứng minh được khả năng tái lập
\cite{Liang2024_241210823v2,Siala2026_260200086v3}.

\section{Mở rộng phạm vi và xử lý trôi dạt khái niệm}

Sau khi quy trình ổn định trên 10 mã, nghiên cứu có thể mở rộng từng bước sang toàn bộ
VN30, các nhóm ngành, HNX hoặc UPCOM. Mỗi lần mở rộng cần khóa danh sách mã, ngày bắt
đầu, lịch giao dịch và nguồn tin. Khả năng tổng quát nên được đánh giá theo mã hoặc theo
ngành, thay vì chỉ báo cáo một điểm trung bình.

Quan hệ giữa các cổ phiếu có thể được mô hình hóa bằng đồ thị ngành, quan hệ doanh
nghiệp hoặc tương quan được học từ dữ liệu quá khứ. HIST và REST là hai hướng khai thác
thông tin liên mã, trong khi DoubleAdapt tập trung vào khả năng thích nghi khi phân
phối thay đổi \cite{Xu2021_211013716v2,Xu2021_210207372v2,Zhao2023_230609862v3}. Nếu
đồ thị được suy ra từ dữ liệu, đồ thị phải được xây dựng lại trong phần huấn luyện của
từng cửa sổ để tránh sử dụng quan hệ của tương lai.

Trôi dạt khái niệm có thể được kiểm tra bằng cách so sánh cửa sổ mở rộng với cửa sổ
trượt, theo dõi chênh lệch giữa tập huấn luyện và kiểm thử, sau đó cập nhật điểm kiểm
theo một lịch cố định. Mỗi lần cập nhật cần lưu phiên bản dữ liệu, thời điểm, cấu hình
và khả năng quay lại điểm kiểm trước. Không nên cập nhật mô hình sau mỗi phiên khi chưa
có đối chứng riêng, vì khó phân biệt khả năng thích nghi với hiện tượng quá khớp.

\section{Đánh giá giá trị kinh tế có khả năng giao dịch}

Một đánh giá tiếp theo phải sử dụng mức giá có thể giao dịch. Tín hiệu được tạo sau khi
có giá đóng cửa và tin đến mốc cắt của ngày $t$ chỉ có thể vào lệnh từ phiên kế tiếp.
Backtest cần khóa quy tắc chuyển xác suất thành vị thế, giá khớp, phí, thuế, trượt giá,
thanh khoản, vòng quay và giới hạn vị thế trước khi xem kết quả. Sổ lệnh cũng cần kiểm
soát trạng thái ròng để tách khả năng chọn mã khỏi biến động chung của thị trường.

Các chỉ số có thể gồm lợi suất tích lũy, lợi suất năm hóa, tỷ số Sharpe, mức suy giảm
tối đa, vòng quay và hệ số chặn so với rổ tham chiếu. Dự đoán phải là dự đoán ngoài mẫu
từ quy trình cuốn chiếu. Không được điều chỉnh chiến lược trên các đoạn kiểm thử rồi
báo cáo lại trên chính các đoạn đó. Mô phỏng giao dịch có thể bổ sung cho chỉ số phân
loại, nhưng kết quả vẫn phụ thuộc vào giả định khớp lệnh và đặc điểm thị trường
\cite{Chen2021_210708721v1,Feng2018_181009936v2}.

\section{Hoàn thiện vận hành và khả năng tái lập}

Đường xử lý hằng ngày cần được chạy vào buổi sáng sau khi nguồn giá đã công bố dữ liệu
của phiên trước và trước phiên khớp lệnh mở cửa. Mỗi dự báo phải ghi thời điểm phát hành,
phiên mục tiêu, phiên giá mới nhất, mã băm đầu vào và điểm kiểm. Chỉ những dòng được phát
hành trước phiên mục tiêu mới được đưa vào thống kê trực tuyến; các lần chạy lại phải
được lưu riêng.

Dashboard nên trình bày dữ liệu được cập nhật đến thời điểm nào, xác suất của ba lớp,
số tin, trạng thái dữ liệu thiếu, phiên bản mô hình và lịch sử các dự báo hợp lệ. Giao
diện không nên gọi đầu ra là khuyến nghị đầu tư. Để tăng khả năng tái lập, hệ thống cần
công bố tệp cấu hình, hạt giống, mã băm dữ liệu, bản kê điểm kiểm, nhật ký chia tập và
quy trình dựng môi trường. Nếu dữ liệu gốc không thể phân phối, có thể công bố lược đồ,
mã băm và quy trình thu thập cùng các giới hạn truy cập.

\section{Lộ trình ưu tiên}

Bảng~\ref{tab:huong-phat-trien-uu-tien} sắp xếp các hướng phát triển theo mức độ cần
thiết đối với kết luận hiện tại. Thứ tự này không dựa trên độ mới của mô hình mà dựa
trên khả năng giảm bất định và hoàn thành các câu hỏi nghiên cứu còn thiếu.

\begin{table}[H]
  \centering
  \small
  \caption{Lộ trình ưu tiên cho các hướng phát triển.}
  \label{tab:huong-phat-trien-uu-tien}
  \begin{tabular}{p{0.9cm} p{3.5cm} p{4.3cm} p{3.2cm}}
    \toprule
    Ưu tiên & Hướng thực hiện & Điều kiện kiểm soát & Kết quả cần thu được \\
    \midrule
    1 & Đánh giá point-in-time lồng và đối chứng nhiều hạt giống &
    Chọn cấu hình chỉ từ nhãn quá khứ; lưu mã băm từng lần rút &
    Phân phối hiệu số và khoảng tin cậy ổn định \\
    2 & Mở rộng nhãn ngẫu nhiên và thêm người gán nhãn &
    Tập đánh giá không chệch; hướng dẫn gán nhãn cố định &
    Chỉ số theo lớp và độ đồng thuận \\
    3 & Hoàn thiện CNN/BiLSTM, phân tầng và giải thích &
    Cùng phần chia, đầu vào, ngân sách và đối chứng &
    Trả lời đầy đủ câu hỏi 1, 3 và 4 \\
    4 & Kiểm tra độ bền theo mã, thời gian và chân trời &
    Lấy mẫu lặp theo khối; ghi nhận số phép thử &
    Phạm vi áp dụng và độ ổn định của kết quả \\
    5 & Thử biểu diễn sự kiện và hợp nhất động &
    Chỉ thực hiện sau khi giao thức cơ sở ổn định &
    So sánh có kiểm soát với LSTM hai nhánh \\
    6 & Backtest có thể giao dịch và theo dõi trực tuyến &
    Giá vào lệnh khả thi, chi phí đầy đủ, phát hành trước phiên &
    Bằng chứng kinh tế và vận hành ngoài mẫu \\
    \bottomrule
  \end{tabular}
\noindent\textit{Nguồn: Tác giả đề xuất từ các giới hạn của đồ án.}
\end{table}

Lộ trình trên bảo vệ mục tiêu cốt lõi của đề tài. Trước hết cần giảm bất định của bộ
sinh cảm xúc và hoàn thành các phép so sánh đã đặt ra. Sau đó mới mở rộng biểu diễn,
kiến trúc và phạm vi thị trường. Đánh giá kinh tế và vận hành trực tuyến được thực hiện
cuối cùng vì hai nội dung này yêu cầu dự đoán ngoài mẫu có thời điểm phát hành hợp lệ,
giả định giao dịch rõ ràng và một giai đoạn quan sát đủ dài.

% ------------------------------------------------------------------ BACK
\printbibliography[title={TÀI LIỆU THAM KHẢO}, heading=bibintoc]

\appendix
\chapter{PHỤ LỤC}
\label{ch:phu-luc}
\section{Hướng dẫn gán nhãn cảm xúc trong miền dữ liệu}

Tập dữ liệu trong miền được gán một trong ba nhãn NEGATIVE, NEUTRAL và POSITIVE. Nhãn phản ánh tác động kỳ vọng của nội dung đối với doanh nghiệp hoặc mã cổ phiếu liên quan, thay vì chỉ phản ánh sắc thái của từng câu riêng lẻ. Người gán nhãn phải đọc tiêu đề và phần nội dung hiện có, ghi nhận các trường hợp chứa nhiều sự kiện hoặc có hướng tác động không rõ ràng, đồng thời không suy luận nhãn dựa trên biến động giá xảy ra sau thời điểm đăng bài.

Các mẫu được mô hình gán nhãn sơ bộ, sau đó được một người rà soát toàn bộ. Báo cáo cuối phải trình bày số người rà soát, số mẫu hợp lệ, phân bố của ba lớp, số mẫu cần chỉnh sửa và nguyên tắc xử lý các mẫu khó. Do chỉ có một người rà soát, Cohen's kappa và Fleiss' kappa không được tính.
\section{Hồ sơ thiết kế hệ thống}

\subsection{Tác nhân và chức năng}

Hệ thống có một tác nhân chính là người nghiên cứu. Người nghiên cứu cung cấp dữ liệu,
khóa cấu hình và kiểm tra các báo cáo được tạo từ quy trình. Các chức năng chính bao gồm:

\begin{itemize}[nosep]
  \item Thu thập và kiểm tra dữ liệu tin tức và dữ liệu giá;
  \item Chuẩn hóa văn bản, loại bỏ bản ghi trùng, đồng thời lưu nguồn gốc và khả năng truy xuất;
  \item Gán tin tức với mã cổ phiếu và phiên giao dịch theo mốc cắt;
  \item Suy luận xác suất của ba lớp cảm xúc bằng PhoBERT;
  \item Tạo đặc trưng theo mã và ngày, sau đó huấn luyện mô hình dự báo;
  \item Đánh giá theo thứ tự thời gian và xuất bản kê, chỉ số và biểu đồ.
\end{itemize}

\subsection{Sơ đồ luồng xử lý}

Sơ đồ dưới đây trình bày luồng xử lý của hệ thống. Khối đánh giá chấm mỗi đoạn kiểm thử cuốn chiếu đúng một lần cùng bản kê của lần chạy. Cấu hình không được điều chỉnh hồi cứu dựa trên
kết quả của các đoạn kiểm thử.

\begin{figure}[H]
  \centering
  \begin{tikzpicture}[
      node distance=5mm,
      block/.style={
        rectangle, rounded corners=2pt, draw=black!60, thick,
        text width=10cm, minimum height=9mm, align=center,
        fill=gray!8, inner sep=4pt},
      decision/.style={
        rectangle, rounded corners=2pt, draw=blue!65, very thick,
        text width=10cm, minimum height=9mm, align=center,
        fill=blue!8, inner sep=4pt},
      arr/.style={-{Stealth[length=2.5mm]}, thick, black!65}
    ]
    \node[block] (input) {Dữ liệu tin tức và dữ liệu giá};
    \node[block, below=of input] (quality) {Kiểm tra lược đồ, mốc thời gian, mã và bản ghi trùng};
    \node[decision, below=of quality] (features) {Tạo đặc trưng cảm xúc và đặc trưng thị trường theo mã và ngày};
    \node[decision, below=of features] (models) {Huấn luyện đối chứng cơ sở, mô hình chỉ sử dụng giá và LSTM hai nhánh};
    \node[block, below=of models] (evaluation) {Đánh giá cuốn chiếu theo thời gian trên năm đoạn kiểm thử, loại bỏ đặc trưng và phân tích sai số};
    \node[block, below=of evaluation] (report) {Chỉ số, biểu đồ, bản kê và báo cáo tái lập};
    \draw[arr] (input) -- (quality);
    \draw[arr] (quality) -- (features);
    \draw[arr] (features) -- (models);
    \draw[arr] (models) -- (evaluation);
    \draw[arr] (evaluation) -- (report);
  \end{tikzpicture}
  \caption{Sơ đồ luồng xử lý và báo cáo kết quả của hệ thống.}
  \label{fig:phu-luc-luong-xu-ly}
\end{figure}
\noindent\textit{Nguồn: Tác giả xây dựng.}

\subsection{Các mô-đun chính}

\begin{description}[style=nextline, nosep]
  \item[\texttt{stf.data.news}] Đọc và chuẩn hóa dữ liệu tin tức, đồng thời xử lý mốc thời gian, tiêu đề,
  nội dung và URL.
  \item[\texttt{stf.data.prices}] Đọc chuỗi giá OHLCV theo mã cổ phiếu và kiểm tra mức độ
  bao phủ của lịch giao dịch.
  \item[\texttt{stf.sentiment.dataset}] Đọc tập dữ liệu có nhãn, chuẩn hóa nhãn về ba lớp và tạo
  các tập huấn luyện, xác thực và kiểm thử theo ngày. Việc chia dữ liệu theo thời gian yêu cầu cột \texttt{date} hoặc \texttt{published\_at} hợp lệ và báo lỗi nếu thiếu.
  \item[\texttt{stf.sentiment.model}] Tinh chỉnh PhoBERT, lưu điểm kiểm và sinh xác suất cảm xúc cho ba lớp.
  \item[\texttt{stf.sentiment.metrics}] Tính F1 vĩ mô cố định trên ba lớp và các chỉ số theo từng lớp.
\end{description}
\section{Cấu hình tái lập}

Bản chụp dữ liệu đầu vào gồm \texttt{data/raw/news/articles.parquet} chứa dữ liệu tin tức,
\texttt{data/raw/news/listings.parquet} chứa liên kết giữa tin tức và mã, các tệp giá theo mã
trong \texttt{data/raw/prices/}, cùng dữ liệu nhãn trong \texttt{data/labeled/}.
Mỗi lần chạy cần lưu phiên bản mã nguồn, phiên bản thư viện, hạt giống, điểm kiểm, mốc cắt,
phạm vi ngày, danh sách mã, mã băm dữ liệu và các tham số huấn luyện. Bản kê được sử dụng
làm nguồn đối chiếu chính cho mọi số liệu trình bày trong phần kết quả thực nghiệm.

Quá trình tiền xử lý chỉ được ước lượng trên tập huấn luyện của từng cửa sổ. Mỗi đoạn kiểm thử
chỉ được truy cập khi báo cáo chỉ số của chính đoạn đó; nghiên cứu không sử dụng một tập giữ lại
cuối tách riêng. Khi thay đổi nguồn dữ liệu, mốc cắt, ngưỡng tạo nhãn hoặc phương pháp chia tập,
phải tạo bản kê mới và không gộp kết quả giữa hai cấu hình.
\section{Cài đặt và chạy hệ thống}

Môi trường được quản lý bằng \texttt{uv} và yêu cầu Python phiên bản 3.12 trở lên. Các lệnh kiểm tra ngoại tuyến tối thiểu gồm:

\begin{verbatim}
uv sync
uv run python -m stf.cli verify
uv run python -m stf.cli sentiment-smoke --n 60
\end{verbatim}

Quá trình huấn luyện PhoBERT yêu cầu tệp nhãn hoàn chỉnh, bao gồm các cột \texttt{text} và \texttt{label}. Do đó, quá trình này chỉ được thực hiện sau khi hướng dẫn gán nhãn đã được khóa và bản kê tương ứng đã được lưu. Việc huấn luyện chính thức nên được tiến hành trên GPU; tên GPU, thời gian chạy và điểm kiểm phải được lưu cùng với kết quả.
\section{Mã nguồn}

Mã nguồn của hệ thống được tổ chức trong gói \texttt{stf} thuộc kho mã nguồn của đề tài.
Các mô-đun thu thập dữ liệu, xử lý cảm xúc, đánh giá và giao diện dòng lệnh được lưu trong thư mục
\texttt{src/stf/}. Notebook dùng để huấn luyện trên GPU và hướng dẫn thực thi được lưu trong thư mục
\texttt{notebooks/}. Phụ lục này trình bày giao diện và luồng xử lý; các kết quả định lượng
chỉ được xác nhận dựa trên sản phẩm đầu ra của một lần chạy có bản kê.

\end{document}
